\section{Tensors that are not normal}

\begin{definition}[Fixed point of a basis of normal tensors]
    \label{def:ldp_non_normal_fixed_point}
    \uses{def:bnt, def:sector_weight_data, def:prep_fixed_point_tensor,
        thm:ldpolar_tensor_decomposition, def:correlation_length}
    Let $A^i=\bigoplus_{j=1}^b\operatorname{diag}(\mu_{j,1},\dots,\mu_{j,m_j})
        \otimes A_j^i$ with normal $A_j$, each in the gauge
    (\ref{eq:ldp_normal_gauge}), forming a basis of normal tensors, and put
    $\beta_j=\sum_{k=1}^{m_j}\mu_{j,k}^N$. Consider lengths $N=qM$ with
    $(\beta_1,\dots,\beta_b)\neq0$. Let $\xi_j$ be the correlation length of
    the normal tensor $A_j$, formed from the subleading eigenvalue of its own
    transfer map $\E_{A_j}$. With $\ket{\omega_j}$ the pair state of $A_j$ and
    $\alpha^{(N)}_j=\beta_j/(\sum_l\lvert\beta_l\rvert^2)^{1/2}$, the
    \emph{fixed-point state} is
    \begin{align}
        \ket{\Omega'}
         & =\sum_{j=1}^b\alpha^{(N)}_j\bigotimes_{k=1}^{M}\ket{\omega_j}_{R_kL_{k+1}}.
        \label{eq:ldp_non_normal_fixed_point}
    \end{align}
    With $V$ from Lemma~\ref{thm:ldpolar_tensor_decomposition} for the blocked map of
    $A$, the \emph{approximating state} is
    $\ket{\widetilde\phi_N}=V^{\otimes M}\ket{\Omega'}/
        \lVert V^{\otimes M}\ket{\Omega'}\rVert$, whenever the denominator is
    nonzero.
    With the canonical form of \cite{Cirac2017MPDO_AnnPhys}, this reproduces the
    fixed-point and approximating states of \cite[eqs.~(19) and~(S7)]{Malz2024Preparation}.

    When some block has multiplicity $m_j\geq2$, the source's construction
    fails: no product of pairs carrying the weights $\beta_j$ gives, after
    $V^{\otimes M}$, an approximation of $\ket{\phi_N}$. As recorded in
    \cite{gap:mswc24_multiplicity_fixed_point}, for
    $A^0=\operatorname{diag}(1,1,0)$ and $A^1=\operatorname{diag}(0,0,1)$ the
    overlap with $\ket{\phi_N}\propto2\ket{0\cdots0}+\ket{1\cdots1}$ tends to
    $1/\sqrt5$ as $M\to\infty$ for every $q$. The failure lies in the
    coefficients $\beta_j$: as shown in Lemma~\ref{lem:ldp_one_copy_state},
    Lemma~\ref{thm:ldp_approximation_error_one_copy}, and
    \cite{gap:mswc24_repeated_block_corrected_state,gap:mswc24_multiplicity_fixed_point},
    placing the pairs of each block on one of its copies
    $k_j$ and replacing $\beta_j$ with the coefficients $\beta'_j=\beta_j(c_j/\mu_{j,k_j}^q)^M$ of
    Definition~\ref{def:ldp_one_copy_weight}, where $c_j$ are the copy norms of
    Definition~\ref{def:ldp_repeated_block_sum}, restores
    convergence for all nonzero complex weights, even when the
    $q$-site states of distinct blocks overlap, at a rate that includes the
    correlation lengths of the mixed transfer maps.

    When a weight is not a nonnegative real number, the coefficients
    $\beta_j=\mu_j^N$ count its phase twice, since $V$ already carries it. For
    $A^0=\operatorname{diag}(1,0)$ and $A^1=\operatorname{diag}(0,-1)$ and odd $N$
    the approximating state is orthogonal to $\ket{\phi_N}$; as shown in
    Section~\ref{sec:ldp_complex_weights} and
    \cite{gap:mswc24_block_form_mixed_overlap}, for multiplicity one the
    coefficients $\beta_j=\lvert\mu_j\rvert^N$ repair the construction.
\end{definition}

\begin{definition}[Fixed-point state of a sector decomposition]
    \label{def:ldp_sector_fixed_point_state}
    \lean{MPSTensor.SectorDecomposition.copyCoord,
        MPSTensor.SectorDecomposition.embeddedFixedPointPair,
        MPSTensor.SectorDecomposition.canonicalFixedPointState,
        MPSTensor.embedPair}
    \leanok
    \uses{def:ldp_non_normal_fixed_point, def:sector_weight_data,
        def:prep_pair_product_state}
    Let $A$ be given by a sector decomposition with blocks $A_1,\dots,A_b$ of
    bond dimensions $D_j$, multiplicities $m_j$ and coefficients $\mu_{j,k}$,
    whose flattened direct sum $\bigoplus_{(j,k)}\mu_{j,k}A_j$ is the tensor of
    Definition~\ref{def:ldp_non_normal_fixed_point} up to the order of the
    bond coordinates, so the bond space of $A$ is $\C^D$ with
    $D=\sum_jm_jD_j$. The blocks carry no normality, basis-of-normal-tensors
    or gauge condition. For each $j$ let $\sigma_j$ be any matrix on
    $\C^{D_j}$ and choose a copy $k_j\in\{1,\dots,m_j\}$. The pair of block
    $j$ is $\sum_{r,l}(\sqrt{\sigma_j})_{rl}\ket{r}\ket{l}$ placed, on both
    legs, on the bond coordinates of copy $k_j$, which gives
    $\ket{\omega_j}\in\C^D\otimes\C^D$. With $\beta_j$ and $\alpha^{(N)}_j$ as
    in Definition~\ref{def:ldp_non_normal_fixed_point}, the fixed-point state
    $\ket{\Omega'}$ is formed from these pairs by
    (\ref{eq:ldp_non_normal_fixed_point}). Here each $\sigma_j$ is an arbitrary
    matrix on $\C^{D_j}$, and the copy $k_j$ is specified explicitly.
\end{definition}

\begin{definition}[Fixed-point state of a family of pairs]
    \label{def:ldp_pair_family_fixed_point}
    \lean{MPSTensor.bntWeight, MPSTensor.ghzAmplitude,
        MPSTensor.sum_norm_sq_pos_of_ne_zero, MPSTensor.ghzAmplitude_norm_sq,
        MPSTensor.nonNormalFixedPointState, MPSTensor.nonNormalApproxState,
        MPSTensor.pairIsometry, Matrix.tensorPower}
    \leanok
    \uses{def:ldp_non_normal_fixed_point, thm:ldpolar_tensor_decomposition,
        def:qc_ghz_state}
    Let $\omega_1,\dots,\omega_b\in\mathbb C^D\otimes\mathbb C^D$ be any vectors,
    $\alpha\in\mathbb C^b$ any coefficients, and $M\ge0$. The state
    $\ket{\Omega'}=\sum_j\alpha_j\bigotimes_{k=1}^M\ket{\omega_j}_{R_kL_{k+1}}$ is
    formed as in (\ref{eq:ldp_non_normal_fixed_point}), and for a tensor $A$ and a
    blocking length $q$, with $V$ from Lemma~\ref{thm:ldpolar_tensor_decomposition}
    for the $q$-site blocked map of $A$, the approximating state is
    $V^{\otimes M}\ket{\Omega'}/\lVert V^{\otimes M}\ket{\Omega'}\rVert$, and zero when
    the denominator vanishes. For multiplicity coefficients $\mu_{j,k}$ put
    $\beta_j=\sum_k\mu_{j,k}^N$ and $\alpha^{(N)}_j=\beta_j/(\sum_l\lvert\beta_l\rvert^2)^{1/2}$.
    For any $\beta\in\mathbb C^b$ with $\beta\neq0$, $\sum_l\lvert\beta_l\rvert^2>0$ and
    $\sum_j\lvert\beta_j/(\sum_l\lvert\beta_l\rvert^2)^{1/2}\rvert^2=1$.
    Let $W\colon\mathbb C^b\to\mathbb C^D\otimes\mathbb C^D$,
    $\ket{j}\mapsto\ket{\omega_j}$, be the map whose $j$-th column is $\omega_j$,
    let $X^{\otimes M}$ denote the $M$-fold tensor power of a matrix $X$, with entries
    $(X^{\otimes M})_{p,s}=\prod_{k=1}^M X_{p_k,s_k}$, and put
    $\ket{\chi_M}=\sum_j\alpha_j\ket{j}^{\otimes M}\in(\mathbb C^b)^{\otimes M}$.
    For $\omega_j$ given by the embedded pair state of each normal block $A_j$,
    $\alpha=\alpha^{(N)}$, and $N=qM$, this specializes to
    Definition~\ref{def:ldp_non_normal_fixed_point}.
\end{definition}

\begin{lemma}[Local orthogonality of the pairs]
    \label{thm:ldp_local_orthogonality}
    \lean{MPSTensor.SectorDecomposition.inner_embeddedFixedPointPair,
        MPSTensor.inner_embedPair_self, MPSTensor.inner_embedPair_eq_zero_of_disjoint,
        MPSTensor.SectorDecomposition.copyCoord_injective,
        MPSTensor.SectorDecomposition.copyCoord_ne_of_ne,
        MPSTensor.SectorDecomposition.toTensor_copyCoord,
        MPSTensor.pairProductState_embedPair,
        MPSTensor.mpv_toTensorFromBlocks_fixedPointTensor}
    \leanok
    \uses{def:ldp_sector_fixed_point_state}
    In Definition~\ref{def:ldp_sector_fixed_point_state}, let each $\sigma_j$ be
    positive semidefinite with $\operatorname{Tr}\sigma_j=1$. For every choice
    of the copies $k_j$ the embedded pairs satisfy
    $\braket{\omega_j}{\omega_{j'}}=\delta_{jj'}$. The chosen coordinates of
    block $j$ carry the block $\mu_{j,k_j}A_j$ of $A$, and embedding the pair
    embeds its product state. For $M\geq1$, in the bond space
    $\C^{D_j}\otimes\C^{D_j}$ of block $j$ alone, the tensor
    $\bigoplus_k\mu_{j,k}^qP_{j,\infty}$ generates on $M$ sites the state
    $\bigl(\sum_k\mu_{j,k}^{qM}\bigr)\ket{\Omega_j}=\beta_j\ket{\Omega_j}$ with
    $N=qM$ and $\ket{\Omega_j}$ formed from the unembedded pair. The
    orthogonality is asserted after \cite[eq.~(S7)]{Malz2024Preparation},
    where the reason given is that the blocks form a basis of normal tensors;
    here it follows from the placement of the pairs alone. As recorded in
    \cite{gap:mswc24_ghz_form_local_orthogonality}, identifying
    these pairs with the fixed points of the positive parts of the blocked
    tensor is not proved here and can fail when a block has multiplicity
    greater than one.
\end{lemma}

\begin{proof}
    \leanok
    \uses{thm:ldp_pair_disjoint_orthogonal, thm:prep_pair_norm,
        thm:prep_mpv_fixed_point_tensor, thm:mpv_decomposition}
    Copies of distinct blocks occupy distinct bond coordinates, so the right
    legs of $\omega_j$ and $\omega_{j'}$ with $j\neq j'$ are supported on
    disjoint sets and the pairs are orthogonal by
    Lemma~\ref{thm:ldp_pair_disjoint_orthogonal}. The coordinates of one
    copy are distinct, so embedding preserves overlaps and
    $\braket{\omega_j}{\omega_j}=\operatorname{Tr}\sigma_j=1$. The direct sum
    generates the sum of the states of its blocks with the $M$-th powers of the
    weights, and $P_{j,\infty}$ generates $\ket{\Omega_j}$.
\end{proof}

\begin{lemma}[GHZ form of the fixed-point state]
    \label{thm:ldp_ghz_form}
    \lean{MPSTensor.isIsometry_pairIsometry, Matrix.IsIsometry.tensorPower,
        MPSTensor.nonNormalFixedPointState_eq_tensorPower_mulVec_ghzState,
        MPSTensor.nonNormalFixedPointState_norm_sq,
        MPSTensor.SectorDecomposition.isIsometry_pairIsometry_embeddedFixedPointPair,
        MPSTensor.SectorDecomposition.canonicalFixedPointState_eq_tensorPower_mulVec_ghzState,
        MPSTensor.SectorDecomposition.canonicalFixedPointState_norm_sq}
    \leanok
    \uses{def:ldp_sector_fixed_point_state, def:ldp_pair_family_fixed_point,
        thm:ldp_local_orthogonality}
    With $W$, $X^{\otimes M}$ and $\chi_M$ as in
    Definition~\ref{def:ldp_pair_family_fixed_point}, formed with
    $\alpha=\alpha^{(N)}$ and the pairs of
    Definition~\ref{def:ldp_sector_fixed_point_state}, where each $\sigma_j$ is
    positive semidefinite of unit trace, $W$ and $W^{\otimes M}$ are isometries
    and $\ket{\Omega'}=W^{\otimes M}\ket{\chi_M}$, up to the regrouping of the
    legs $R_kL_{k+1}$. If $M\geq1$ and $(\beta_1,\dots,\beta_b)\neq0$,
    then $\braket{\Omega'}{\Omega'}=1$, matching the preparation of
    \cite[eq.~(19)]{Malz2024Preparation}. The same holds for the state of
    Definition~\ref{def:ldp_pair_family_fixed_point} with $\alpha=\alpha^{(N)}$
    for any orthonormal pairs $\omega_j$; the identity
    $\ket{\Omega'}=W^{\otimes M}\ket{\chi_M}$ holds for any pairs.
\end{lemma}

\begin{proof}
    \leanok
    \uses{thm:prep_pair_product_norm, thm:ldp_local_orthogonality}
    The pairs of Definition~\ref{def:ldp_sector_fixed_point_state} are
    orthonormal by Lemma~\ref{thm:ldp_local_orthogonality}. The $(j,j')$ entry of $W^\dagger W$ is $\braket{\omega_j}{\omega_{j'}}$, and the
    entries of $(W^{\otimes M})^\dagger W^{\otimes M}$ factor over the sites,
    \begin{align*}
        \bigl((W^{\otimes M})^\dagger W^{\otimes M}\bigr)_{s,t}
        =\prod_{k=1}^M(W^\dagger W)_{s_k,t_k}=\delta_{st}.
    \end{align*}
    Expanding
    $\ket{j}^{\otimes M}$ in the product basis gives
    $W^{\otimes M}\ket{j}^{\otimes M}=\bigotimes_k\ket{\omega_j}$, and summing with the
    weights $\alpha^{(N)}_j$ gives $\ket{\Omega'}$. Overlaps of the product states factor
    over the bonds,
    $\braket{\Omega_j}{\Omega_{j'}}=\braket{\omega_j}{\omega_{j'}}^M=\delta_{jj'}$, so
    $\braket{\Omega'}{\Omega'}=\sum_j\lvert\alpha^{(N)}_j\rvert^2=1$.
\end{proof}

\begin{lemma}[Pairs with disjoint supports are orthogonal]
    \label{thm:ldp_pair_disjoint_orthogonal}
    \lean{MPSTensor.inner_pair_eq_zero_of_disjoint}
    \leanok
    \uses{def:ldp_sector_fixed_point_state}
    Let $S,S'\subseteq\{0,\dots,D-1\}$ be disjoint, and let
    $\omega,\omega'\in\C^D\otimes\C^D$ satisfy $\omega_{(a,c)}=0$ for $a\notin S$
    and $\omega'_{(a,c)}=0$ for $a\notin S'$. Then
    $\braket{\omega}{\omega'}=0$. The embedded pairs of distinct blocks in
    Definition~\ref{def:ldp_sector_fixed_point_state} are supported in this way, so
    this gives the off-diagonal half of the orthogonality
    $\braket{\omega_j}{\omega_{j'}}=\delta_{jj'}$, which
    Lemma~\ref{thm:ldp_local_orthogonality} completes.
\end{lemma}

\begin{proof}
    \leanok
    Every term of the sum vanishes: if $a\in S$ then $a\notin S'$ and
    $\omega'_{(a,c)}=0$, and otherwise $\omega_{(a,c)}=0$, so
    \begin{align*}
        \braket{\omega}{\omega'}
        =\sum_{(a,c)}\overline{\omega_{(a,c)}}\,\omega'_{(a,c)}=0.
    \end{align*}
\end{proof}

\begin{lemma}[Approximation error, general case]
    \label{thm:ldp_approximation_error_general}
    \uses{def:ldp_non_normal_fixed_point, def:correlation_length}
    \notready
    With the notation of Definition~\ref{def:ldp_non_normal_fixed_point},
    including the approximating state $\widetilde\phi_N$ defined there,
    $\xi_{\mathrm{diag}}=\max_j\xi_j$ the largest correlation length of the
    blocks $A_j$, $0<\gamma<1/2$, and $q=o(N)$,
    \begin{align}
        \varepsilon(\widetilde\phi_N,\phi_N)
         & =O\!\left(\frac Nq\,e^{-\gamma q/\xi_{\mathrm{diag}}}\right).
        \label{eq:ldp_approximation_error_general}
    \end{align}
    As printed in \cite[Lemma~1$'$(ii)]{Malz2024Preparation} and analyzed in
    \cite{gap:mswc24_multiplicity_fixed_point,gap:mswc24_block_form_mixed_overlap},
    this bound is false: it fails both for a block of multiplicity two
    (Lemma~\ref{thm:ldp_approximation_error_repeated_block}, with the pair of
    that block on one of its copies; the source does not say where the pair
    sits, and the same computation, recorded in the gap note but not
    formalized, gives the failure for every placement) and for two
    blocks of multiplicity one whose states on $q$ sites are not orthogonal
    (Lemma~\ref{thm:ldp_approximation_error_overlapping_blocks}).
    For blocks of multiplicity one whose states on $q$ sites are orthogonal it
    holds without the condition $q=o(N)$
    (Lemma~\ref{thm:ldp_approximation_error_orthogonal_blocks}). For blocks of
    multiplicity one with real weights $0\leq w_j\leq1$, the largest equal to one,
    whose states overlap, the approximating state with $\beta_j=w_j^N$
    converges at the rate obtained by replacing $\xi_{\mathrm{diag}}$ with
    $\max(\xi_{\mathrm{diag}},\xi_{\mathrm{off\text{-}diag}})$, the larger of the
    correlation lengths of the blocks and of their mixed transfer maps, again
    without the condition $q=o(N)$
    (Lemmas~\ref{thm:ldp_approximation_error_overlapping_corrected}
    and~\ref{thm:ldp_approximation_error_overlapping_weights}).
\end{lemma}
%% The direct sum A is not normal, so Lemma lem:ldp_positive_part_overlap
%% does not apply to it. The source's proof rests on the block form (S5) of
%% the positive part, which holds only when the q-site states of distinct
%% blocks are orthogonal and every multiplicity is one; the two theorems
%% below refute the statement outside that case, and the statement is
%% proved in that case (thm:ldp_approximation_error_orthogonal_blocks). The source writes the
%% diagonal correlation lengths with a double index, referring to the
%% diagonal blocks E_{A_j} of the transfer map of A as opposed to the mixed
%% blocks; here xi_j denotes the correlation length of A_j. This entry carries
%% no formalization tag: the statement is false as printed, and the
%% corrected statement below adds the orthogonality of the blocks.

% **False source result:** refutes Lemma 1'(ii) of arXiv:2307.01696 for a
% block of multiplicity two; docs/paper-gaps/mswc24_multiplicity_fixed_point.tex.
% The first linked declaration certifies that the displayed tensor is, after
% ordering its bond coordinates, the canonical form (S2) of a basis of normal
% tensors in canonical form II with the weight normalization of the source,
% so the refutation holds on the domain on which Lemma 1'(ii) is stated.
\begin{lemma}[The approximating state for a repeated block]
    \label{thm:ldp_approximation_error_repeated_block}
    \lean{MPSTensor.isBNTCanonicalForm_and_not_approximationError_le_repeatedBlock,
        MPSTensor.isBNTCanonicalForm_repeatedBlockSector,
        MPSTensor.repeatedBlockSector, MPSTensor.coeff_repeatedBlockSector,
        MPSTensor.reindex_toTensor_repeatedBlockSector,
        MPSTensor.nonNormalApproxOverlap_repeatedBlockTensor,
        MPSTensor.tendsto_norm_nonNormalApproxOverlap_repeatedBlockTensor,
        MPSTensor.norm_le_of_hasEigenvalue_transferMap_repeatedBlockBasis,
        MPSTensor.not_approximationError_le_repeatedBlockTensor}
    \leanok
    \uses{def:ldp_non_normal_fixed_point, def:ldp_sector_fixed_point_state,
        def:correlation_length, def:sector_bnt_cf, def:nt_cpsv,
        def:cpsv_cfii}
    Let $A^0=\operatorname{diag}(1,1,0)$ and $A^1=\operatorname{diag}(0,0,1)$,
    the canonical form with the normal block $A_1=(1,0)$ of multiplicity two and
    weights $(1,1)$ and the normal block $A_2=(0,1)$ of multiplicity one, so that
    $\beta=(2,1)$ and $\ket{\phi_N}\propto2\ket{0\cdots0}+\ket{1\cdots1}$.
    The blocks are normal and in canonical form~II, they form a basis of normal
    tensors whose states become orthogonal as $N\to\infty$, and the weights satisfy $\lvert\mu_{j,k}\rvert\leq1$ with one of
    modulus one, so $A$ satisfies every hypothesis of
    \cite[Lemma~1$'$(ii)]{Malz2024Preparation}. Place
    the pair of the first block on its first copy. For every $q\geq1$ and
    $M\geq1$ the approximating state of
    Definition~\ref{def:ldp_non_normal_fixed_point} satisfies
    \begin{align}
        \braket{\widetilde\phi_N}{\phi_N}
         & =\frac{4\cdot2^{-M/2}+1}{\sqrt5\,(4\cdot2^{-M}+1)^{1/2}},
        \label{eq:ldp_repeated_block_overlap}
    \end{align}
    which does not depend on $q$ and tends to $1/\sqrt5$ as $M\to\infty$. The
    transfer maps of both blocks have no eigenvalue other than one, so every
    $\lambda_2$ with $0<\lvert\lambda_2\rvert<1$ bounds their subleading
    eigenvalues. For every such $\lambda_2$, every $\gamma>0$ and every constant
    $C$ there is $M\geq1$ such that, with $q=M$, $N=M^2$ and
    $y=(N/q)e^{-\gamma q/\xi_{\mathrm{diag}}}$,
    $\varepsilon(\widetilde\phi_N,\phi_N)>Cy\,e^{Cy}$. Because the sequence $q=M$
    satisfies $q=o(N)$, this tensor refutes the bound (\ref{eq:ldp_approximation_error_general}),
    as recorded in \cite{gap:mswc24_multiplicity_fixed_point}.
\end{lemma}

\begin{proof}
    \leanok
    \uses{thm:ldpolar_decomposition, thm:ldp_polar_uniqueness}
    Each block is one-dimensional with $\sum_i\lvert a_i\rvert^2=1$, so its
    transfer map is the identity. It is therefore irreducible and normal, and
    left-canonical with the diagonal positive-definite fixed point $1$, which is
    the canonical-form~II condition; the pairs below are built from this same
    fixed point. Its self-overlap is one. The
    states $\ket{0\cdots0}$ and $\ket{1\cdots1}$ of the blocks are linearly
    independent for $N\geq1$, and the blocks are not related by a gauge
    transformation and a phase, since $A_1^1=0\neq A_2^1$. With the bond
    coordinates ordered as the two copies of $A_1$ followed by the copy of
    $A_2$, and the weights $(1,1)$ and $1$, the direct sum of the copies is $A$:
    for $i=0,1$,
    \begin{align}
        A^i & =A_1^i\oplus A_1^i\oplus A_2^i
        =\operatorname{diag}\bigl(A_1^i,A_1^i,A_2^i\bigr).
        \label{eq:ldp_repeated_block_assembly}
    \end{align}

    For a tensor of diagonal matrices the blocked map factors as $B=CJ^\dagger$,
    where $J\colon\ket e\mapsto\ket{e\,e}$ and the column $e$ of $C$ is the state
    of the diagonal coordinate $e$ on $q$ sites; here these are $\ket{0\cdots0}$
    for $e=0,1$ and $\ket{1\cdots1}$ for $e=2$. Hence $B^\dagger B=JGJ^\dagger$ with
    $G$ the Gram matrix of the columns. The positive semidefinite square root of $G$
    is $Q=2^{-1/2}\bigl(\begin{smallmatrix}1&1\\1&1\end{smallmatrix}\bigr)\oplus1$,
    and $E=\tfrac12\bigl(\begin{smallmatrix}1&1\\1&1\end{smallmatrix}\bigr)\oplus1$
    is the projector onto its range, so uniqueness of positive square roots gives
    $P=JQJ^\dagger$ and $\Pi=JEJ^\dagger$. The pairs are $e_{00}$ and $e_{22}$, so
    $V^{\otimes M}\sum_j\alpha_j\ket{\Omega_j}$ has components
    $\sum_j\alpha_j\prod_kV_{\tau_k,(c_j,c_j)}$. From $V^\dagger B=P$ and
    $V^\dagger V=\Pi$, the overlap with the unnormalized target is
    $\sum_j\sum_e\overline{\alpha_j}Q_{c_je}^M$, the squared norm of the
    approximating vector is $\sum_{j,j'}\overline{\alpha_j}\alpha_{j'}E_{c_jc_{j'}}^M$,
    and $\lVert\phi_N(A)\rVert^2=\sum_{e,e'}G_{ee'}^M=5$. With
    $\alpha=(2,1)/\sqrt5$ this gives (\ref{eq:ldp_repeated_block_overlap}). Since
    $1/\sqrt5<1/2$, the error exceeds $1/2$ for large $M$, while $Cy\,e^{Cy}$ with
    $y=Mr^M$ and $r=e^{-\gamma/\xi_{\mathrm{diag}}}<1$ tends to zero.
\end{proof}

% **False source result:** refutes Lemma 1'(ii) of arXiv:2307.01696 for two
% blocks of multiplicity one; docs/paper-gaps/mswc24_block_form_mixed_overlap.tex.
% As for thm:ldp_approximation_error_repeated_block, the first linked
% declaration certifies the hypotheses of Lemma 1'(ii) for the displayed tensor.
\begin{lemma}[The approximating state for overlapping blocks]
    \label{thm:ldp_approximation_error_overlapping_blocks}
    \lean{MPSTensor.isBNTCanonicalForm_and_not_approximationError_le_overlappingBlock,
        MPSTensor.isBNTCanonicalForm_overlappingBlockSector,
        MPSTensor.reindex_toTensor_overlappingBlockSector,
        MPSTensor.nonNormalApproxOverlap_overlappingBlockTensor,
        MPSTensor.le_one_sub_norm_nonNormalApproxOverlap_overlappingBlockTensor,
        MPSTensor.one_sub_norm_nonNormalApproxOverlap_overlappingBlockTensor_le,
        MPSTensor.norm_le_of_hasEigenvalue_transferMap_overlappingBlockBasis,
        MPSTensor.not_approximationError_le_overlappingBlockTensor}
    \leanok
    \uses{def:ldp_non_normal_fixed_point, def:ldp_sector_fixed_point_state,
        def:correlation_length, def:sector_bnt_cf, def:nt_cpsv,
        def:cpsv_cfii}
    Let $A^0=\operatorname{diag}(1,3/5)$ and $A^1=\operatorname{diag}(0,4/5)$,
    the canonical form with the normal blocks $A_1=(1,0)$ and $A_2=(3/5,4/5)$, each
    of multiplicity one and weight one, so that $\beta=(1,1)$. The blocks are
    normal and in canonical form~II and form a basis of normal tensors whose
    states become orthogonal as $N\to\infty$, so $A$
    satisfies every hypothesis of \cite[Lemma~1$'$(ii)]{Malz2024Preparation}.
    Let $s=(3/5)^q$,
    the overlap of the states of the two blocks on $q$ sites,
    $u=\bigl(\sqrt{1+s}+\sqrt{1-s}\bigr)/2$ and
    $v=\bigl(\sqrt{1+s}-\sqrt{1-s}\bigr)/2$. For every $q\geq1$ and $M\geq1$ the
    approximating state of Definition~\ref{def:ldp_non_normal_fixed_point}
    satisfies
    \begin{align}
        \braket{\widetilde\phi_N}{\phi_N}
         & =\frac{u^M+v^M}{(1+s^M)^{1/2}},
        \label{eq:ldp_overlapping_blocks_overlap}
    \end{align}
    and for $M\geq3$ the error is at least $(9/25)^q/16$, while for every $M\geq1$
    it is at most $M(9/25)^q$, so the construction converges whenever
    $(N/q)(9/25)^q\to0$. The transfer maps of both
    blocks have no eigenvalue other than one, so every $\lambda_2$ bounds their
    subleading eigenvalues. If $r=e^{-\gamma/\xi_{\mathrm{diag}}}<9/25$, then for every
    constant $C$ there is $q\geq3$ such that, with $M=q$, $N=q^2$ and
    $y=(N/q)e^{-\gamma q/\xi_{\mathrm{diag}}}$,
    $\varepsilon(\widetilde\phi_N,\phi_N)>Cy\,e^{Cy}$. Because the sequence $q=M$ satisfies
    $q=o(N)$, the bound (\ref{eq:ldp_approximation_error_general}) fails even
    when every multiplicity is one, as recorded in \cite{gap:mswc24_block_form_mixed_overlap}. What
    fails is the rate: for $M\geq3$ the error lies between $(9/25)^q/16$ and
    $M(9/25)^q$, where $9/25=\tau^2$ is the square of the eigenvalue $\tau=3/5$ of
    the mixed transfer map of the two blocks.
\end{lemma}

\begin{proof}
    \leanok
    \uses{thm:ldpolar_decomposition, thm:ldp_polar_uniqueness}
    The hypotheses are checked as in the proof of
    Lemma~\ref{thm:ldp_approximation_error_repeated_block}: both blocks are
    one-dimensional with $\sum_i\lvert a_i\rvert^2=1$, their states
    $\ket{0\cdots0}$ and $(\tfrac35\ket0+\tfrac45\ket1)^{\otimes N}$ are
    linearly independent for $N\geq1$, and $A_1^1=0\neq A_2^1$.

    As in the proof of Lemma~\ref{thm:ldp_approximation_error_repeated_block},
    the Gram matrix of the blocked map on the diagonal coordinates is
    $G=\bigl(\begin{smallmatrix}1&s\\s&1\end{smallmatrix}\bigr)$, whose positive
    square root is $Q=\bigl(\begin{smallmatrix}u&v\\v&u\end{smallmatrix}\bigr)$,
    since $u^2+v^2=1$ and $2uv=s$. The determinant of $Q$ is $\sqrt{1-s^2}>0$, so
    the support projector is $JJ^\dagger$. With $\alpha=(1,1)/\sqrt2$ the overlap
    with the unnormalized target is $\sqrt2\,(u^M+v^M)$, the approximating vector
    has norm one, and $\lVert\phi_N(A)\rVert^2=2+2s^M$, which gives
    (\ref{eq:ldp_overlapping_blocks_overlap}). For $M\geq3$,
    \begin{align*}
        u^M+v^M\leq u^3+v^2=1-u^2(1-u)\leq1-\frac{s^2}{16},
    \end{align*}
    because $u^2=\bigl(1+\sqrt{1-s^2}\bigr)/2\geq1/2$ and
    $1-u\geq(1-u^2)/2\geq s^2/8$. For the upper bound, the overlap is one when
    $M=1$, since $u+v=\sqrt{1+s}$. For $M\geq2$,
    $u\geq u^2=\bigl(1+\sqrt{1-s^2}\bigr)/2\geq1-s^2/2$, so $u^M\geq1-Ms^2/2$, and
    $(1+s^M)^{1/2}\leq1+s^2/2$; when $Ms^2<1$ this gives
    \begin{align*}
        \frac{u^M+v^M}{(1+s^M)^{1/2}}\geq\frac{1-Ms^2/2}{1+s^2/2}\geq1-Ms^2,
    \end{align*}
    and when $Ms^2\geq1$ the bound $\varepsilon\leq1$ suffices. Finally $Cy\,e^{Cy}$
    with $y=qr^q$ is eventually below $(9/25)^q/16$, because $q(25r/9)^q\to0$.
\end{proof}

\begin{lemma}[The polar matrix has a first-order perturbation]
    \label{thm:ldp_overlapping_polar_norm_lower_bound}
    \lean{MPSTensor.posDef_diagGram_overlappingBlockDiag,
        MPSTensor.overlappingBlockOverlap_div_two_le_norm_sqrt_diagGram_sub_one,
        MPSTensor.overlappingBlockOverlap_div_two_le_norm_polarPos_sub_limit,
        MPSTensor.exists_norm_polarPos_overlappingBlockTensor_sub_limit_gt_sq}
    \leanok
    \uses{thm:ldp_approximation_error_overlapping_blocks}
    For the tensor of Lemma~\ref{thm:ldp_approximation_error_overlapping_blocks},
    put $s=(3/5)^q$, let $G_q$ be the Gram matrix on its diagonal bond pairs,
    and let $J\ket j=\ket{j,j}$. For every $q\ge1$, $G_q$ is positive definite
    and the positive polar matrix is $P_q=J\sqrt{G_q}J^\dagger$. In operator norm,
    \[
        \lVert\sqrt{G_q}-I\rVert\ge s/2,
        \qquad \lVert P_q-JJ^\dagger\rVert\ge s/2.
    \]
    For every $C\ge0$ and every threshold $q_0$, there is $q\ge\max(1,q_0)$
    such that $\lVert P_q-JJ^\dagger\rVert>Cs^2$. Thus this matrix difference
    is not $O((3/5)^{2q})$. The quadratic normalized-state error of
    Lemma~\ref{thm:ldp_approximation_error_overlapping_second_order}
    concerns a different quantity.
\end{lemma}
\begin{proof}\leanok
    With $\sqrt{G_q}=\bigl(\begin{smallmatrix}u&v\\v&u\end{smallmatrix}\bigr)$
    as in the preceding theorem, $u,v\ge0$, $u^2+v^2=1$ and $2uv=s$ imply
    $u\le1$ and $v\ge s/2$. Each operator norm is at least the
    absolute value of this off-diagonal entry. The Gram matrix is positive
    definite since its positive square root is invertible for $q\ge1$.
    Finally, $(3/5)^q\to0$, so beyond any threshold one may choose $q$ with
    $C(3/5)^q<1/2$, giving $Cs^2<s/2$.
\end{proof}

\begin{lemma}[Polar decomposition of an orthogonal sum]
    \label{lem:ldp_orthogonal_sum_polar}
    \lean{Matrix.polarIso_sum_of_orthogonal, Matrix.exists_polarIso_eq_mul,
        Matrix.conjTranspose_polarIso_mul_polarIso_eq_zero,
        Matrix.conjTranspose_polarIso_mul_eq_zero}
    \leanok
    \uses{thm:ldp_polar_uniqueness}
    Let $B=\sum_jc_jB_jK_j^\dagger$ with $c_j>0$, isometries $K_j$ with
    orthogonal ranges, so that $K_j^\dagger K_{j'}=\delta_{jj'}$, and
    $B_j^\dagger B_{j'}=0$ for $j\neq j'$. With $B_j=V_jP_j$ the polar
    decompositions of the summands, the polar decomposition of $B$ has
    \begin{align*}
        V=\sum_jV_jK_j^\dagger,\qquad P=\sum_jc_jK_jP_jK_j^\dagger.
    \end{align*}
    Relative to \cite[eq.~(S5)]{Malz2024Preparation} (``Since $V$ is an isometry, $P$
    inherits the block structure of $A$'') and
    \cite{gap:mswc24_block_form_mixed_overlap}, this establishes the block form
    of the positive part for blocks of multiplicity one under the orthogonality
    of their states, an additional hypothesis without which the block form
    fails.
\end{lemma}
\begin{proof}
    \leanok
    Every partial isometry is $V_j=B_jR_j$ for some $R_j$: the support
    projector $\Pi_j$ satisfies $\Pi_j=P_jR_j$ for some $R_j$, because both
    have the same range, and $V_j=V_j\Pi_j=B_jR_j$. Hence
    $V_j^\dagger V_{j'}=R_j^\dagger B_j^\dagger B_{j'}R_{j'}=0$ for $j\neq j'$. Let
    $W=\sum_jV_jK_j^\dagger$, $Q=\sum_jc_jK_jP_jK_j^\dagger\geq0$ and
    $E=\sum_jK_j\Pi_jK_j^\dagger$. The orthogonality of the $K_j$ and of the
    $V_j$ removes all cross terms, so $WQ=B$, $W^\dagger W=E$, $E^2=E=E^\dagger$,
    $EQ=Q$ and $Q\sum_jc_j^{-1}K_jR_jK_j^\dagger=E$. Thus $E$ is the projector
    onto the range of $Q$, and Lemma~\ref{thm:ldp_polar_uniqueness} gives
    $V=W$ and $P=Q$.
\end{proof}

\begin{lemma}[Approximation error for orthogonal blocks of multiplicity one]
    \label{thm:ldp_approximation_error_orthogonal_blocks}
    \lean{MPSTensor.exists_approximationError_le_mul_blockSum,
        MPSTensor.exists_approximationError_le_blockSum,
        MPSTensor.nonNormalApproxVector_blockSum,
        MPSTensor.mixedMapLM_eq_zero_of_conjTranspose_physicalMatrix_mul_eq_zero,
        MPSTensor.mpvOverlap_eq_zero,
        MPSTensor.eq_zero_of_hasEigenvalue_mixedMapLM_of_conjTranspose_mul_eq_zero}
    \leanok
    \uses{def:ldp_non_normal_fixed_point, def:ldp_pair_family_fixed_point,
        lem:ldp_orthogonal_sum_polar, lem:ldp_mixed_transfer_decay,
        thm:ldp_approximation_error_overlapping_weights, def:correlation_length}
    Let $A^i=\bigoplus_{j=1}^b\mu_jA_j^i$ with $b\geq1$, weights $\mu_j>0$, and
    normal blocks $A_j$ in the gauge (\ref{eq:ldp_normal_gauge}) with fixed
    points $\rho_j$, each placed on its own bond coordinates. Let $\lambda_2$
    bound the moduli of the eigenvalues other than one of every transfer map
    $\E_{A_j}$, so that its correlation length $\xi$ bounds every $\xi_j$, and
    let $0<\gamma<1$. Form the approximating state $\widetilde\phi_N$ of
    Definition~\ref{def:ldp_non_normal_fixed_point} with the pairs of the
    $\rho_j$ on the bond coordinates of their blocks and $\beta_j=\mu_j^N$.
    There is $C>0$ such that for all weights $\mu_j>0$, every block length $q$
    at which the $q$-site states of distinct blocks are orthogonal,
    $\sum_{i_1,\dots,i_q}\overline{(A_j^{i_1}\cdots A_j^{i_q})_{ab}}\,
        (A_{j'}^{i_1}\cdots A_{j'}^{i_q})_{a'b'}=0$ for $j\neq j'$ and all
    $a,b,a',b'$, and every $M\geq1$, with $N=qM$ and
    $y=(N/q)\,e^{-\gamma q/\xi}$,
    \begin{align}
        \varepsilon(\widetilde\phi_N,\phi_N)\leq Cy
        \qquad\text{and}\qquad
        \varepsilon(\widetilde\phi_N,\phi_N)\leq Cy\,e^{Cy}.
        \label{eq:ldp_approximation_error_orthogonal_blocks}
    \end{align}
    Relative to \cite[Lemma~1$'$(ii)]{Malz2024Preparation} and
    \cite{gap:mswc24_block_form_mixed_overlap}, this establishes the bound
    restricted to multiplicity one and to blocks whose $q$-site states are
    orthogonal, for $0<\gamma<1$ in place of $0<\gamma<1/2$. For
    multiplicity one the phase of $\mu_j$ can be absorbed into $A_j$, which
    stays normal with the same transfer map, so the condition $\mu_j>0$ only
    fixes this presentation. The orthogonality is an added hypothesis without
    which the statement fails
    (Lemma~\ref{thm:ldp_approximation_error_overlapping_blocks}), whereas
    under it no condition $q=o(N)$ is needed.
\end{lemma}
\begin{proof}
    \leanok
    If the $q$-site states of distinct blocks are orthogonal for no $q$, there is
    nothing to prove. Otherwise let them be orthogonal at $q_0$, and let $B_j$ be
    the $q_0$-site blocked map of $A_j$. The entries of the mixed transfer map
    $Z\mapsto\sum_iB_j^iZ(B_{j'}^i)^\dagger$ are sums
    $\sum_iB_j^{i}{}_{ab}\overline{B_{j'}^{i}{}_{a'b'}}$, the complex conjugates of
    the entries of $B_{j'}^\dagger B_j=0$, so this map vanishes. It is
    $\E_{jj'}^{\,q_0}$, so $\E_{jj'}$ is nilpotent and all its eigenvalues are
    zero. Hence $\lambda_2$ also bounds the moduli of the eigenvalues of the mixed
    transfer maps, and Lemma~\ref{thm:ldp_approximation_error_overlapping_weights}
    with $w_j=\mu_j$ gives both bounds, with a constant that does not depend on the
    weights.
\end{proof}

%% ==================================================
%% Subsection: Blocks with multiplicities
%% Source: arXiv:2307.01696 (Malz, Styliaris, Wei, Cirac), Supplemental
%% Material, "Proof of Lemma 1 and extension to non-normal tensors",
%% eqs. (S2)-(S7) and Lemma 1'(ii).
%% Local fix: the approximating state of eq. (S7) is replaced by the fixed
%% point of the corrected block form of the positive part;
%% docs/paper-gaps/mswc24_repeated_block_corrected_state.tex.
%% ==================================================

\subsection{Blocks with multiplicities}\label{sec:ldp_repeated_blocks}

When a block of the canonical form is repeated, the positive part of the
blocked map is not the direct sum printed in the source, and the state built
from nearest-neighbour pairs with the source's coefficients does not approximate
the target
\cite{gap:mswc24_multiplicity_fixed_point}. This subsection gives the correct
block form, the corrected approximating state, and the approximation error
for it when the states of distinct blocks on $q$ sites are orthogonal
\cite{gap:mswc24_repeated_block_corrected_state}.

\begin{definition}[Direct sum with multiplicities]
    \label{def:ldp_repeated_block_sum}
    \lean{MPSTensor.CopyWeights, MPSTensor.CopyWeights.weight_fun_ne_zero,
        MPSTensor.repeatedBlockSum, MPSTensor.copyNorm, MPSTensor.copyNorm_sq,
        MPSTensor.copyNorm_pos, MPSTensor.copyIsometry,
        MPSTensor.conjTranspose_copyIsometry}
    \leanok
    Let $A_1,\dots,A_b$ be tensors, let the isometries $E_{j,k}$,
    $1\leq k\leq m_j$, place the copy $k$ of block $j$ in the bond space, with
    $E_{j,k}^\dagger E_{j',k'}=\delta_{jj'}\delta_{kk'}\mathbf 1$, and let
    $\mu_{j,k}$ be nonzero weights, with at least one copy of every block. The
    direct sum with multiplicities is
    $A^i=\sum_{j,k}\mu_{j,k}E_{j,k}A_j^iE_{j,k}^\dagger$. For $q\geq1$ and
    $K_{j,k}=E_{j,k}\otimes E_{j,k}$, define
    \begin{align*}
        c_j=\Bigl(\sum_k\lvert\mu_{j,k}\rvert^{2q}\Bigr)^{1/2},\qquad
        L_j=\sum_k\frac{\overline{\mu_{j,k}^q}}{c_j}K_{j,k}.
    \end{align*}
\end{definition}

\begin{theorem}[Positive part of a blocked tensor with multiplicities]
    \label{thm:ldp_repeated_block_polar}
    \lean{MPSTensor.repeatedBlockSum_eq_blockSum,
        MPSTensor.flat_disjoint_coord,
        MPSTensor.physicalMatrix_blockTensor_repeatedBlockSum,
        MPSTensor.conjTranspose_pairEmbedding_mul_copy,
        MPSTensor.mpv_repeatedBlockSum,
        MPSTensor.physicalMatrix_blockTensor_repeatedBlockSum_eq_copyIsometry,
        MPSTensor.conjTranspose_copyIsometry_mul_self,
        MPSTensor.conjTranspose_copyIsometry_mul_eq_zero,
        MPSTensor.polarIso_blockTensor_repeatedBlockSum,
        MPSTensor.polarPos_blockTensor_repeatedBlockSum}
    \leanok
    \uses{def:ldp_repeated_block_sum, lem:ldp_orthogonal_sum_polar}
    Let $A$ be the direct sum with multiplicities of
    Definition~\ref{def:ldp_repeated_block_sum}. Its periodic state on $N\geq1$
    sites is $\sum_j\beta_j\ket{\phi_N(A_j)}$ with $\beta_j=\sum_k\mu_{j,k}^N$.
    For $q\geq1$ the blocked map of $A$ is $B=\sum_jc_jB_jL_j^\dagger$, with
    $B_j$ the blocked map of $A_j$, and $L_j^\dagger L_{j'}=\delta_{jj'}\mathbf 1$.
    If $B_j^\dagger B_{j'}=0$ for $j\neq j'$, then, with $B_j=V_jP_j$,
    \begin{align*}
        V=\sum_jV_jL_j^\dagger,\qquad P=\sum_jc_jL_jP_jL_j^\dagger.
    \end{align*}
    On the copies of block $j$ the positive part is therefore of rank one in
    the copy index, where \cite[eq.~(S5)]{Malz2024Preparation} prints
    $\operatorname{diag}(\mu_{j,1}^q,\dots,\mu_{j,m_j}^q)\otimes P_j$
    \cite{gap:mswc24_repeated_block_corrected_state}.
\end{theorem}
\begin{proof}
    \leanok
    Listing all copies as blocks of multiplicity one, every word of length
    $q\geq1$ is a direct sum, so $B=\sum_{j,k}\mu_{j,k}^qB_jK_{j,k}^\dagger$ and
    $\phi_N(A)=\sum_{j,k}\mu_{j,k}^N\phi_N(A_j)$. Grouping the copies of each
    block gives $B=\sum_jc_jB_jL_j^\dagger$, since
    $c_j\,\mu_{j,k}^q/c_j=\mu_{j,k}^q$. Also
    $L_j^\dagger L_{j'}=\sum_{k,k'}(\mu_{j,k}^q/c_j)
        (\overline{\mu_{j',k'}^q}/c_{j'})K_{j,k}^\dagger K_{j',k'}$, which is
    $\sum_k\lvert\mu_{j,k}\rvert^{2q}/c_j^2=1$ times the identity for $j=j'$ and
    zero otherwise. Lemma~\ref{lem:ldp_orthogonal_sum_polar} gives the polar
    decomposition.
\end{proof}

\begin{definition}[Corrected approximating state]
    \label{def:ldp_repeated_block_state}
    \lean{MPSTensor.copyFixedPointState, MPSTensor.copyApproxVector,
        MPSTensor.copyApproxOverlap, MPSTensor.copyApproxOverlap_eq}
    \leanok
    \uses{def:ldp_non_normal_fixed_point, def:ldp_repeated_block_sum,
        thm:ldp_repeated_block_polar}
    With the notation of Definition~\ref{def:ldp_repeated_block_sum} and
    Theorem~\ref{thm:ldp_repeated_block_polar}, fixed points
    $\rho_j$ of the blocks, and $\ket{\Omega_j}$ the product of the pairs of
    $\rho_j$ over $M$ sites in the bond space of block $j$, the corrected
    approximating state on $N=qM$ sites with coefficients $\alpha_j$ is
    \begin{align*}
        \ket{\widetilde\phi_N}\propto V^{\otimes M}\sum_j\alpha_j
        L_j^{\otimes M}\ket{\Omega_j}.
    \end{align*}
    When the weights $\beta_j$ of Theorem~\ref{thm:ldp_repeated_block_polar} are
    not all zero, the source's coefficients are
    $\alpha_j=\beta_j/\bigl(\sum_l\lvert\beta_l\rvert^2\bigr)^{1/2}$; if all
    $\beta_j$ vanish these coefficients are undefined.
    Replacing each $P_j$ in the corrected block form of
    Theorem~\ref{thm:ldp_repeated_block_polar} by its fixed point, just as
    \cite[eq.~(S7)]{Malz2024Preparation} does with the printed block form,
    yields this state with the weights of the source. Because the isometry
    $L_j$ acts on the two legs of each site, for a repeated block the state
    $L_j^{\otimes M}\ket{\Omega_j}$ is not a product of nearest-neighbour pairs
    of the full bond space; for multiplicity one and $\mu_j>0$, as recorded in
    \cite{gap:mswc24_repeated_block_corrected_state}, it reduces to the state
    of Definition~\ref{def:ldp_non_normal_fixed_point}.
\end{definition}

\begin{theorem}[Approximation error for orthogonal blocks with multiplicities]
    \label{thm:ldp_approximation_error_repeated_orthogonal_blocks}
    \lean{MPSTensor.exists_approximationError_le_repeatedBlockSum,
        MPSTensor.exists_approximationError_le_mul_repeatedBlockSum,
        MPSTensor.copyApproxVector_repeatedBlockSum}
    \leanok
    \uses{def:ldp_repeated_block_state, thm:ldp_repeated_block_polar,
        thm:ldp_approximation_error_orthogonal_blocks,
        thm:ldp_approximation_error_repeated_overlapping, def:correlation_length}
    With the notation of Theorem~\ref{thm:ldp_repeated_block_polar}, let the
    blocks $A_j$ be normal in the gauge (\ref{eq:ldp_normal_gauge}) with fixed
    points $\rho_j$, let $\lambda_2$ bound the moduli of the eigenvalues other
    than one of every transfer map $\E_{A_j}$, with correlation length $\xi$, and
    let $0<\gamma<1$. There are constants $C,C'>0$ such that for all nonzero
    complex weights, every block length $q$ at which the $q$-site states of
    distinct blocks are orthogonal, $B_j^\dagger B_{j'}=0$ for $j\neq j'$, and
    every $M\geq1$ with $N=qM$ and $(\beta_1,\dots,\beta_b)\neq0$, the corrected
    approximating state of Definition~\ref{def:ldp_repeated_block_state}
    satisfies, with $y=(N/q)\,e^{-\gamma q/\xi}$,
    \begin{align*}
        \varepsilon(\widetilde\phi_N,\phi_N) & \leq C'y,        &
        \varepsilon(\widetilde\phi_N,\phi_N) & \leq Cy\,e^{Cy}.
    \end{align*}
    Relative to \cite[Lemma~1$'$(ii)]{Malz2024Preparation} and
    \cite{gap:mswc24_repeated_block_corrected_state,gap:mswc24_block_form_mixed_overlap},
    this establishes the bound for the corrected approximating state with
    arbitrary multiplicities and complex weights, restricted to blocks whose
    $q$-site states are orthogonal, for $0<\gamma<1$ in place of
    $0<\gamma<1/2$, and without the condition $q=o(N)$.
    When the states are orthogonal, the corrected state is the combination
    $\sum_j\alpha_j\ket{\phi_M(V_jP_{j,\infty})}$ of the approximating states of the
    blocks.
\end{theorem}
\begin{proof}
    \leanok
    By Theorem~\ref{thm:ldp_repeated_block_polar}, $VL_j=V_j$ at every $q$ at which
    the states are orthogonal, which gives the last statement. For the bounds,
    suppose the states are orthogonal at some $q_0$; otherwise there is nothing to
    prove. As in the proof of
    Theorem~\ref{thm:ldp_approximation_error_orthogonal_blocks}, every mixed
    transfer map $\E_{jj'}$, $j\neq j'$, is then nilpotent, so $\lambda_2$ bounds the
    moduli of its eigenvalues, and
    Theorem~\ref{thm:ldp_approximation_error_repeated_overlapping} gives both bounds.
\end{proof}

%% ==================================================
%% Subsection: Blocks whose states overlap
%% Source: arXiv:2307.01696 (Malz, Styliaris, Wei, Cirac), Supplemental
%% Material, "Proof of Lemma 1 and extension to non-normal tensors",
%% eqs. (S2)-(S12) and Lemma 1'(ii).
%% Local fix: the rate of Lemma 1'(ii) is corrected by including the
%% correlation lengths of the mixed transfer maps; scope restriction to
%% blocks of multiplicity one with real weights;
%% docs/paper-gaps/mswc24_block_form_mixed_overlap.tex.
%% ==================================================

\subsection{Blocks whose states overlap}\label{sec:ldp_overlapping_blocks}

When the states of distinct blocks on $q$ sites are not orthogonal, the positive
part of the blocked map is not block diagonal, and the rate of
\cite[Lemma~1$'$(ii)]{Malz2024Preparation} fails
(Lemma~\ref{thm:ldp_approximation_error_overlapping_blocks}). The construction
itself still converges. This subsection proves the approximation error for the
source's approximating state, for blocks of multiplicity one with real weights, at the
rate set by the larger of the correlation lengths of the blocks and of their mixed
transfer maps \cite{gap:mswc24_block_form_mixed_overlap}.

\begin{lemma}[Decay of mixed transfer maps]
    \label{lem:ldp_mixed_transfer_decay}
    \lean{MPSTensor.mixedMapLM_blockTensor_apply,
        MPSTensor.conjTranspose_physicalMatrix_mul_physicalMatrix_apply,
        MPSTensor.exists_norm_mixedMapLM_pow_apply_le,
        MPSTensor.exists_norm_gram_blockTensor_mixed_le,
        MPSTensor.exists_norm_mpvOverlap_le_of_mixedMapLM,
        MPSTensor.exp_neg_div_correlationLength_le_one}
    \leanok
    \uses{def:correlation_length, thm:ldp_positive_part_rate}
    Let $X$ and $Y$ be tensors with $d$ physical indices and bond dimensions
    $D_1,D_2\geq1$, with mixed transfer map
    $\E_{XY}(Z)=\sum_iX^iZ(Y^i)^\dagger$ on $D_1\times D_2$ matrices, and let
    $X_q$, $Y_q$ be their blocked tensors on $q$ sites. Then
    $\E_{X_qY_q}=\E_{XY}^{\,q}$, and for tensors with a common number of physical
    indices the mixed Gram matrix of their physical matrices, read as maps from bond
    pairs to physical indices, is a rearrangement of the mixed transfer map,
    $(\mathcal B_X^\dagger\mathcal B_Y)_{(a,b),(a',b')}
        =\E_{YX}(\ket{b'}\bra b)_{a'a}$. For $\lambda_2\in\C$ with
    $\lvert\lambda_2\rvert\leq1$, correlation length $\xi$, and $\gamma\geq0$,
    $e^{-\gamma/\xi}\leq1$. Now let $\lvert\lambda_2\rvert<1$ and $0<\gamma<1$, and
    write $r=e^{-\gamma/\xi}$.
    \begin{enumerate}
        \item If every eigenvalue of $\E_{XY}$ has modulus at most
              $\lvert\lambda_2\rvert$, there are $C,K$ with
              $\lVert\E_{XY}^{\,n}(Z)\rVert\leq Cr^n\lVert Z\rVert$ for all $n$ and $Z$,
              and
              $\bigl\lvert\sum_s\phi_N(X)_s\,\overline{\phi_N(Y)_s}\bigr\rvert\leq Kr^N$
              for all $N$.
        \item If every eigenvalue of $\E_{YX}$ has modulus at most
              $\lvert\lambda_2\rvert$, there is $K'$ with
              $\lVert\mathcal B_X^\dagger\mathcal B_Y\rVert\leq K'r^q$ for all $q$.
    \end{enumerate}
    For distinct blocks $A_j$, $A_{j'}$ of a basis of normal tensors the spectral
    radius $\tau_{jj'}$ of $\E_{jj'}=\E_{A_jA_{j'}}$ is below one, and the source
    writes $\xi_{jj'}=-1/\ln\tau_{\max}$ with $\tau_{\max}=\tau_{jj'}$
    \cite[eqs.~(S14)--(S15)]{Malz2024Preparation}.
\end{lemma}
\begin{proof}
    \leanok
    The blocked tensors have letters $X^{i_1}\cdots X^{i_q}$, so
    $\E_{X_qY_q}(Z)=\sum_{i_1,\dots,i_q}X^{i_1}\cdots X^{i_q}Z
        (Y^{i_1}\cdots Y^{i_q})^\dagger=\E_{XY}^{\,q}(Z)$, and the entry formula is
    $\sum_i\overline{X^i_{ab}}\,Y^i_{a'b'}=\E_{YX}(\ket{b'}\bra b)_{a'a}$. Every point of
    the spectrum of $\E_{XY}$ is an eigenvalue, of modulus at most
    $\lvert\lambda_2\rvert<r$ by Lemma~\ref{thm:ldp_positive_part_rate}, so the
    spectral radius is below $r$ and Gelfand's formula gives
    $\lVert\E_{XY}^{\,n}(Z)\rVert\leq Cr^n\lVert Z\rVert$. The overlap is
    $\Tr\E_{XY}^{\,N}=\sum_{p,q}\E_{XY}^{\,N}(\ket p\bra q)_{pq}$, and every entry is
    at most a fixed multiple of the norm. The Gram matrix is bounded by the sum of
    its entries, each of which is an entry of $\E_{YX}^{\,q}$ applied to a matrix
    unit. Finally $-\gamma/\xi=\gamma\ln\lvert\lambda_2\rvert\leq0$.
\end{proof}

\begin{lemma}[Shared estimates]
    \label{lem:ldp_shared_overlap_estimates}
    \lean{Matrix.l2_opNorm_le_one_of_conjTranspose_mul_self_eq_one,
        MPSTensor.conjTranspose_mul_sum_of_isometry,
        MPSTensor.sum_mul_conjTranspose_mul_sum_of_isometry,
        Matrix.conjTranspose_tensorPower,
        Matrix.sum_star_tensorPower_mulVec_mul,
        Matrix.sum_norm_sq_tensorPower_mulVec_le,
        Matrix.tensorPower_one, Matrix.tensorPower_zero,
        MPSTensor.conjTranspose_coordEmbedding,
        MPSTensor.transpose_kronecker_apply_eq_mul_single_mul,
        MPSTensor.exists_norm_gram_sum_sub_le,
        Matrix.norm_polarPos_sub_le_sqrt,
        MPSTensor.exists_abs_sum_norm_sq_sum_mpv_sub_le,
        MPSTensor.sum_star_tensorPower_polarIso_mulVec_mul_mpv,
        MPSTensor.inv_mul_norm_le_norm_of_sum_norm_sq_le_one,
        norm_prod_range_sub_pow_le_of_forall_norm_pow_le,
        MPSTensor.exists_norm_mpvOverlap_sub_trace_pow_le}
    \leanok
    \uses{lem:ldp_mixed_transfer_decay, thm:ldp_positive_part_rate, thm:ldp_sqrt_holder,
        thm:ldp_approximation_error_normal, lem:ldp_positive_part_overlap,
        thm:overlap_trace_rect}
    Let $A_1,\dots,A_b$ be tensors with $d$ physical indices and bond dimensions $D_j$,
    normal in the gauge (\ref{eq:ldp_normal_gauge}) with positive definite fixed points
    $\rho_j$ of trace one, with $\mathcal B_j$ the blocked map of $A_j$ on $q$ sites. Let
    $\lambda_2\in\C$ with $\lvert\lambda_2\rvert<1$ bound the moduli of the eigenvalues
    other than one of every transfer map $\E_{A_j}$ and the moduli of all eigenvalues of
    the mixed transfer maps $\E_{jj'}$, $j\neq j'$, and let $\xi$ be its correlation
    length.
    \begin{enumerate}
        \item A matrix $L$ with $L^\dagger L=\Id$ has operator norm at most one, and for
              isometries $L_j$ with $L_j^\dagger L_{j'}=0$ for $j\neq j'$,
              $L_j^\dagger\sum_kL_kY_k=Y_j$ and
              $\bigl(\sum_jL_jY_jL_j^\dagger\bigr)\bigl(\sum_jL_jZ_jL_j^\dagger\bigr)
                  =\sum_jL_jY_jZ_jL_j^\dagger$.
        \item For matrices $W,W'$ with the same rows and vectors $\psi,\varphi$ on $M$
              sites, $(W^{\otimes M})^\dagger=(W^\dagger)^{\otimes M}$,
              $\braket{W^{\otimes M}\psi}{W'^{\otimes M}\varphi}
                  =\braket{\psi}{(W^\dagger W')^{\otimes M}\varphi}$; if $W^\dagger W$ is
              idempotent, then $\lVert W^{\otimes M}\psi\rVert\leq\lVert\psi\rVert$; and on
              $M\geq1$ sites $\Id^{\otimes M}=\Id$ and $0^{\otimes M}=0$.
        \item For a coordinate isometry $E$, $E^\dagger=E^{\mathsf T}$; for rectangular
              $X$ and $Y$, the entry of $X^{\mathsf T}\otimes Y$ at $((l,r),(a,c))$ is
              that of $X\ket l\bra rY$ at $(a,c)$.
        \item For $0<\gamma<1$ there is $K$ such that for all $R$, all coefficients
              $\lvert c_j\rvert\leq R$, all matrices $L_j$ of norm at most one and every
              $q$, the matrix $\mathcal B=\sum_jc_j\mathcal B_jL_j^\dagger$ satisfies
              $\bigl\lVert\mathcal B^\dagger\mathcal B-\sum_j\lvert c_j\rvert^2
                  L_j(\rho_j^{\mathsf T}\otimes\Id)L_j^\dagger\bigr\rVert
                  \leq R^2Ke^{-\gamma q/\xi}$.
        \item For every matrix $\mathcal B$ and positive semidefinite $P_\infty$, the
              positive part $P$ of $\mathcal B$ satisfies
              $\lVert P-P_\infty\rVert\leq\lVert\mathcal B^\dagger\mathcal B-P_\infty^2
                  \rVert^{1/2}$.
        \item For $0<\gamma<1$ there is $K$ such that for all $R$, all coefficients
              $\lvert\beta_j\rvert\leq R$ and every $N$,
              $\bigl\lvert\lVert\sum_j\beta_j\phi_N(A_j)\rVert^2-\sum_j\lvert\beta_j\rvert^2
                  \bigr\rvert\leq R^2Ke^{-\gamma N/\xi}$.
        \item For a tensor $A$ with blocked map $\mathcal B=VP$ on $q$ sites and any vector
              $F$ on the bond pairs of $M$ sites,
              $\braket{V^{\otimes M}F}{\phi_N(A)}=\braket{F}{\phi_M(P)}$ with $N=qM$.
              If $\lVert v\rVert\leq1$ and $t\geq0$, then
              $t^{-1}\lvert\braket vw\rvert\leq\bigl\lvert\braket vw\bigr\rvert/
                  (\lVert v\rVert t)$, the right side read as zero when $v=0$ or $t=0$.
        \item In a normed ring, if $\lVert E^k\rVert\leq c$ for every $k\geq0$ and
              $\lVert X_k-E\rVert\leq\delta$ for every $k$, then
              $\lVert X_1\cdots X_M-E^M\rVert\leq c((1+c\delta)^M-1)$. Consequently, for
              a tensor $F$ of bond dimension $D_2$ and $c\geq0$ there is $C>0$ such that
              for all $G,G'$ from bond pairs of dimension $D_1$ to bond pairs of
              dimension $D_2$ with $\lVert G-G'\rVert\leq\delta$, if the powers of the
              mixed transfer map $T'$ of $G'$ against $F$ satisfy
              $\lVert T'^kX\rVert\leq c\lVert X\rVert$, then the overlap of $G$ with $F$
              on $M$ sites is $\Tr T'^M$ up to $C(M\delta)\,e^{CM\delta}$.
    \end{enumerate}
\end{lemma}
\begin{proof}
    \leanok
    Part~(1) is $\lVert L\rVert^2=\lVert L^\dagger L\rVert$. Part~(2) is
    $\bigl(W^{\otimes M}\bigr)^\dagger W'^{\otimes M}=(W^\dagger W')^{\otimes M}$,
    together with $\Id-\Pi\geq0$ for the Hermitian idempotent
    $\Pi=(W^\dagger W)^{\otimes M}$. Part~(3) is a computation of entries. For part~(4),
    the blocks of $\mathcal B^\dagger\mathcal B$ are
    $\overline{c_j}c_{j'}L_j\mathcal B_j^\dagger\mathcal B_{j'}L_{j'}^\dagger$, with
    coefficients of modulus at most $R^2$ and the $L_j$ of norm at most one; the diagonal
    blocks are estimated by Lemma~\ref{thm:ldp_positive_part_rate} and the others by
    Lemma~\ref{lem:ldp_mixed_transfer_decay}. Part~(5) is
    Lemma~\ref{thm:ldp_sqrt_holder} at $P_\infty^2$, whose square root is $P_\infty$.
    For part~(6),
    $\lVert\sum_j\beta_j\phi_N(A_j)\rVert^2
        =\sum_{j,j'}\beta_j\overline{\beta_{j'}}\braket{\phi_N(A_{j'})}{\phi_N(A_j)}$ with
    coefficients of modulus at most $R^2$; the proof of
    Lemma~\ref{thm:ldp_approximation_error_normal} bounds the diagonal terms minus one,
    and Lemma~\ref{lem:ldp_mixed_transfer_decay} bounds the others. For part~(7), the
    target is $V^{\otimes M}\ket{\phi_M(P)}$ by $\mathcal B=VP$, and
    $(V^\dagger V)^{\otimes M}\ket{\phi_M(P)}=\ket{\phi_M(P)}$ because $V^\dagger V$ is the
    projector onto the range of $P$, so part~(2) gives the inner product; the second
    statement is $\lVert v\rVert^{-1}\geq1$. The first statement of part~(8) is the
    induction in the proof of Lemma~\ref{lem:ldp_positive_part_overlap}, which only uses
    the bound on the powers of $E$; the second is that estimate in the normed algebra of
    operators on rectangular matrices, with the overlap the trace of the $M$-th power of
    the mixed transfer map (Lemma~\ref{thm:overlap_trace_rect}).
\end{proof}

\begin{lemma}[Positive part for overlapping blocks]
    \label{lem:ldp_overlapping_positive_part}
    \lean{MPSTensor.blockSumGramLimit, MPSTensor.blockSumPosLimit,
        MPSTensor.physicalMatrix_blockTensor_blockSum_one,
        MPSTensor.gram_blockTensor_blockSum,
        MPSTensor.exists_norm_gram_blockTensor_blockSum_sub_le,
        MPSTensor.exists_norm_polarPos_blockTensor_blockSum_sub_le}
    \leanok
    \uses{lem:ldp_shared_overlap_estimates, lem:ldp_mixed_transfer_decay,
        thm:ldp_positive_part_rate, thm:ldp_sqrt_lipschitz, thm:ldp_polar_uniqueness}
    Let $A^i=\bigoplus_{j=1}^b\mu_jA_j^i$ be the direct sum with weights $\mu_j\in\C$ of
    blocks $A_j$ of bond dimensions $D_j$, each placed on its own bond coordinates by an
    isometry $E_j$, and put $K_j=E_j\otimes E_j$. Let every $A_j$ be normal in the
    gauge (\ref{eq:ldp_normal_gauge}) with positive definite fixed point $\rho_j$ of
    trace one. Let $\lambda_2\in\C$ with $\lvert\lambda_2\rvert<1$ bound the moduli of
    the eigenvalues other than one of every transfer map $\E_{A_j}$ and the moduli of
    all eigenvalues of the mixed transfer maps $\E_{jj'}$, $j\neq j'$, let $\xi$ be its
    correlation length, and let $0<\gamma<1$. For $q\geq1$ the blocked map of $A$ is
    $\mathcal B=\sum_j\mu_j^q\mathcal B_jK_j^\dagger$, with $\mathcal B_j$ the blocked map
    of $A_j$, so
    $\mathcal B^\dagger\mathcal B=\sum_{j,j'}\overline{\mu_j^q}\mu_{j'}^q
        K_j\mathcal B_j^\dagger\mathcal B_{j'}K_{j'}^\dagger$. There is $K$ such that for
    all weights with $\lvert\mu_j\rvert\leq1$ and every $q\geq1$
    \begin{align*}
        \Bigl\lVert\mathcal B^\dagger\mathcal B-\sum_j\lvert\mu_j\rvert^{2q}
        K_j\bigl(\rho_j^{\mathsf T}\otimes\Id\bigr)K_j^\dagger\Bigr\rVert
        \leq Ke^{-\gamma q/\xi}.
    \end{align*}
    For weights one, $\mathcal B=\sum_j\mathcal B_jK_j^\dagger$, and with
    \begin{align*}
        G_\infty=\sum_jK_j\bigl(\rho_j^{\mathsf T}\otimes\Id\bigr)K_j^\dagger,\qquad
        P_\infty=\sum_jK_j\bigl(\sqrt{\rho_j}^{\,\mathsf T}\otimes\Id\bigr)K_j^\dagger,
    \end{align*}
    there are $K,K'$ such that for every $q\geq1$
    \begin{align*}
        \lVert\mathcal B^\dagger\mathcal B-G_\infty\rVert\leq Ke^{-\gamma q/\xi},
        \qquad
        \lVert P-P_\infty\rVert\leq K'e^{-\gamma q/\xi},
    \end{align*}
    where $P$ is the positive part of $\mathcal B$. This replaces the block form
    asserted in \cite[eq.~(S5)]{Malz2024Preparation}, which holds only when the
    states of distinct blocks on $q$ sites are orthogonal
    \cite{gap:mswc24_block_form_mixed_overlap}.
\end{lemma}
\begin{proof}
    \leanok
    Every word of $A$ of length $q\geq1$ is the direct sum of the words of the blocks,
    multiplied by $\mu_j^q$, which gives the form of $\mathcal B$ and of its Gram
    matrix.
    The diagonal terms, with coefficients $\lvert\mu_j\rvert^{2q}\leq1$, satisfy
    $\lVert\mathcal B_j^\dagger\mathcal B_j-\rho_j^{\mathsf T}\otimes\Id\rVert
        \leq K_je^{-\gamma q/\xi}$ by Lemma~\ref{thm:ldp_positive_part_rate}, and the
    others $\lVert\mathcal B_j^\dagger\mathcal B_{j'}\rVert\leq K_{jj'}e^{-\gamma q/\xi}$
    by Lemma~\ref{lem:ldp_mixed_transfer_decay}, with coefficients of modulus at most one;
    summing gives the first bound. For weights one, let
    $Q=\Id-\sum_jK_jK_j^\dagger$, the projector onto the bond pairs whose two indices
    lie in different blocks or outside all blocks. Then $\mathcal BQ=0$, so
    $PQ=V^\dagger\mathcal BQ=0$ and $QP=0$, whence $(P+Q)^2=\mathcal B^\dagger\mathcal B+Q$;
    likewise $P_\infty Q=QP_\infty=0$ and $P_\infty^2=G_\infty$, so
    $(P_\infty+Q)^2=G_\infty+Q$. Both $P+Q$ and $P_\infty+Q$ are positive
    semidefinite, so they are the square roots. If $\rho_j\geq c_j\Id$ and
    $c=\bigl(1+\sum_j1/c_j\bigr)^{-1}$, then
    \begin{align*}
        G_\infty+Q-c\Id
        =\sum_jK_j\bigl((\rho_j-c\Id)^{\mathsf T}\otimes\Id\bigr)K_j^\dagger+(1-c)Q
        \geq0,
    \end{align*}
    so Lemma~\ref{thm:ldp_sqrt_lipschitz} gives
    $\lVert P-P_\infty\rVert=\lVert\sqrt{\mathcal B^\dagger\mathcal B+Q}
        -\sqrt{G_\infty+Q}\rVert\leq\lVert\mathcal B^\dagger\mathcal B-G_\infty\rVert/\sqrt c$.
\end{proof}

\begin{lemma}[Overlaps for overlapping blocks]
    \label{lem:ldp_overlapping_overlaps}
    \lean{MPSTensor.embeddedBlockState, MPSTensor.embedPair_fixedPointPair,
        MPSTensor.sum_star_nonNormalApproxVector_mul_mpv,
        MPSTensor.mixedMapLM_blockSumPosLimit_apply,
        MPSTensor.exists_norm_mpvOverlap_polarPosTensor_blockSum_sub_one_le,
        MPSTensor.exists_abs_norm_mpvState_blockSum_sq_sub_le}
    \leanok
    \uses{lem:ldp_overlapping_positive_part, lem:ldp_shared_overlap_estimates,
        lem:ldp_mixed_transfer_decay, def:ldp_non_normal_fixed_point, def:prep_fixed_point_tensor,
        lem:ldp_positive_part_overlap, thm:ldp_approximation_error_normal}
    \begin{enumerate}
        \item For an isometry $E$ placing the bond coordinates of a block and a
              positive semidefinite $\sigma$, the pair of $\sigma$ placed along $E$ on
              both legs is the pair of $\tilde\sigma=E\sigma E^\dagger$.
        \item For a tensor $A$, matrices $\sigma'_1,\dots,\sigma'_b$ on the bond space,
              coefficients $\alpha_j$, and $M\geq1$, the state
              $V^{\otimes M}\sum_j\alpha_j\bigotimes_k\ket{\omega(\sigma'_j)}_{R_kL_{k+1}}$
              has inner product
              $\sum_j\overline{\alpha_j}\,\braket{\phi_M(P_{\infty,\sigma'_j})}{\phi_M(P)}$
              with $\ket{\phi_N(A)}$, $N=qM$, where $P$ is the positive part of the
              blocked map and $P_{\infty,\sigma'_j}$ the fixed-point tensor of
              $\sigma'_j$.
        \item For blocks placed as in Lemma~\ref{lem:ldp_overlapping_positive_part}
              and positive semidefinite $\rho_j$, with $\tilde\rho_j=E_j\rho_jE_j^\dagger$
              and $\Pi_j=E_jE_j^\dagger$, the mixed transfer map of $P_\infty$ against the
              fixed-point tensor of $\tilde\rho_j$ is $X\mapsto\Tr(\Pi_jX)\,\tilde\rho_j$.
        \item In the setting of Lemma~\ref{lem:ldp_overlapping_positive_part}, for
              every block $j$ there is $C_j>0$ such that for every $q\geq1$ and every
              $M\geq1$, with $y=Me^{-\gamma q/\xi}$,
              \begin{align*}
                  \bigl\lvert\braket{\phi_M(P_{\infty,\tilde\rho_j})}{\phi_M(P)}-1\bigr\rvert
                  \leq C_jy\,e^{C_jy}.
              \end{align*}
        \item In the setting of Lemma~\ref{lem:ldp_overlapping_positive_part} with
              $0<\gamma<1$, there is $K$ with
              $\bigl\lvert\lVert\phi_N(A)\rVert^2-b\bigr\rvert\leq Ke^{-\gamma N/\xi}$
              for every $N\geq1$.
    \end{enumerate}
\end{lemma}
\begin{proof}
    \leanok
    Part~(1) holds because $\sqrt{E\sigma E^\dagger}=E\sqrt\sigma E^\dagger$. Part~(2) is
    Lemma~\ref{lem:ldp_shared_overlap_estimates}(7) for the product of pairs. For
    part~(3), the letter of $P_\infty$ at the physical pair $(l,r)$ is
    $\sum_k\sqrt{\tilde\rho_k}\ket l\bra r\Pi_k$, that of the fixed-point tensor is
    $\sqrt{\tilde\rho_j}\ket l\bra r$, and
    $\sum_{l,r}\ket l\bra rY\ket r\bra l=\Tr(Y)\Id$, while
    $\sqrt{\tilde\rho_k}\sqrt{\tilde\rho_j}=\delta_{kj}\tilde\rho_j$. Since
    $\Tr(\Pi_j\tilde\rho_j)=\Tr\rho_j=1$, this map is idempotent with trace one, and
    the telescoping estimate in the proof of
    Lemma~\ref{lem:ldp_positive_part_overlap}, applied with
    $\lVert P-P_\infty\rVert\leq K'e^{-\gamma q/\xi}$ from
    Lemma~\ref{lem:ldp_overlapping_positive_part}, gives part~(4). Part~(5) is
    Lemma~\ref{lem:ldp_shared_overlap_estimates}(6) with $\beta_j=1$ and $R=1$, since
    $\phi_N(A)=\sum_j\phi_N(A_j)$ for $N\geq1$.
\end{proof}

\begin{lemma}[Overlaps for weights of different moduli]
    \label{lem:ldp_overlapping_overlaps_weights}
    \lean{MPSTensor.mixedMapLM_blockSumPosLimit_smul_apply,
        MPSTensor.exists_norm_mpvOverlap_sub_pow_le_of_norm_sub_blockSumPosLimit_le,
        MPSTensor.exists_abs_norm_mpvState_blockSum_weight_sq_sub_le}
    \leanok
    \uses{lem:ldp_overlapping_overlaps, lem:ldp_overlapping_positive_part,
        lem:ldp_shared_overlap_estimates, lem:ldp_positive_part_overlap, lem:ldp_mixed_transfer_decay,
        thm:ldp_approximation_error_normal}
    \begin{enumerate}
        \item For blocks placed as in Lemma~\ref{lem:ldp_overlapping_positive_part},
              positive semidefinite $\rho_k$ and reals $t_k\geq0$, the mixed transfer map
              of $P_\infty^{t}=\sum_kK_k\bigl(\sqrt{t_k^2\rho_k}^{\,\mathsf T}\otimes\Id\bigr)
                  K_k^\dagger$ against the fixed-point tensor of $\tilde\rho_j$ is
              $X\mapsto t_j\Tr(\Pi_jX)\,\tilde\rho_j$.
        \item For blocks placed as there, positive semidefinite $\rho_k$ of trace one, and
              a block $j$, there is $C>0$ such that for all reals $t_k\geq0$ with
              $t_j\leq1$, all $\delta$, all $M\geq1$, and every matrix $P$ on the bond
              pairs with $\lVert P-P_\infty^{t}\rVert\leq\delta$,
              \begin{align*}
                  \bigl\lvert\braket{\phi_M(P_{\infty,\tilde\rho_j})}{\phi_M(P)}-t_j^M\bigr\rvert
                  \leq C(M\delta)\,e^{CM\delta}.
              \end{align*}
        \item In the setting of Lemma~\ref{lem:ldp_overlapping_positive_part} with
              $0<\gamma<1$, there is $K$ such that for all weights with
              $\lvert\mu_j\rvert\leq1$ and every $N\geq1$,
              $\bigl\lvert\lVert\phi_N(A)\rVert^2-\sum_j\lvert\mu_j^N\rvert^2\bigr\rvert
                  \leq Ke^{-\gamma N/\xi}$ for $A=\bigoplus_j\mu_jA_j$.
    \end{enumerate}
\end{lemma}
\begin{proof}
    \leanok
    For part~(1), $\sqrt{t_k^2\rho_k}=t_k\sqrt{\rho_k}$, so the letter of
    $P_\infty^{t}$ at $(l,r)$ is $\sum_kt_k\sqrt{\tilde\rho_k}\ket l\bra r\Pi_k$, and the
    computation of Lemma~\ref{lem:ldp_overlapping_overlaps}(3) carries the factor $t_j$.
    For part~(2), the mixed transfer map of $P_\infty^{t}$ is $t_jR_j$ with $R_j$ the
    idempotent of trace one of Lemma~\ref{lem:ldp_overlapping_overlaps}(3); its powers
    $t_j^kR_j$, $k\geq1$, have norm at most $\lVert R_j\rVert$, and its $M$-th power has
    trace $t_j^M$. Lemma~\ref{lem:ldp_shared_overlap_estimates}(8) with $D_1=D_2=D$, the
    fixed-point tensor of $\tilde\rho_j$, and $c=1+\lVert R_j\rVert$ gives the estimate.
    Part~(3) is Lemma~\ref{lem:ldp_shared_overlap_estimates}(6) with $R=1$, with
    $\lVert\phi_N(A)\rVert^2=\sum_{j,j'}\mu_j^N\overline{\mu_{j'}^N}
        \braket{\phi_N(A_{j'})}{\phi_N(A_j)}$ and coefficients of modulus at most one.
\end{proof}

% **Local fix (rate of the overlapping blocks):** the rate of Lemma 1'(ii) of
% arXiv:2307.01696 is corrected by including the correlation lengths of the
% mixed transfer maps; scope restriction to multiplicity one and weight one;
% docs/paper-gaps/mswc24_block_form_mixed_overlap.tex.
\begin{lemma}[Approximation error for overlapping blocks]
    \label{thm:ldp_approximation_error_overlapping_corrected}
    \lean{MPSTensor.exists_approximationError_le_overlappingBlockSum,
        MPSTensor.exists_approximationError_le_mul_overlappingBlockSum}
    \leanok
    \uses{def:ldp_non_normal_fixed_point, thm:ldp_approximation_error_complex_weights,
        def:correlation_length}
    Let $A^i=\bigoplus_{j=1}^bA_j^i$ with $b\geq1$ be the direct sum with weights one
    of normal blocks $A_j$ in the gauge (\ref{eq:ldp_normal_gauge}), with positive
    definite fixed points $\rho_j$ of trace one, each placed on its own bond
    coordinates. Let $\lambda_2\in\C$ bound the moduli of the eigenvalues other than one
    of every transfer map $\E_{A_j}$ and the moduli of all eigenvalues of the mixed
    transfer maps $\E_{jj'}$, $j\neq j'$, so that, when $\lvert\lambda_2\rvert<1$, its
    correlation length $\xi$ bounds every $\xi_j$ and every $\xi_{jj'}$, and let
    $0<\gamma<1$. Form the
    approximating state $\widetilde\phi_N$ of
    Definition~\ref{def:ldp_non_normal_fixed_point} with the pairs of the $\rho_j$ on
    the bond coordinates of their blocks and $\beta_j=1$. There are constants $C,C'>0$
    such that for every block length $q$ and every $M\geq1$, with $N=qM$ and
    $y=Me^{-\gamma q/\xi}$, which is $(N/q)e^{-\gamma q/\xi}$ for $q\geq1$,
    \begin{align}
        \varepsilon(\widetilde\phi_N,\phi_N)\leq C'y
        \qquad\text{and}\qquad
        \varepsilon(\widetilde\phi_N,\phi_N)\leq Cy\,e^{Cy}.
        \label{eq:ldp_approximation_error_overlapping_corrected}
    \end{align}
    Relative to \cite[Lemma~1$'$(ii)]{Malz2024Preparation} for blocks of multiplicity one
    and weight one, this replaces $\xi_{\mathrm{diag}}$ by
    $\max(\xi_{\mathrm{diag}},\xi_{\mathrm{off\text{-}diag}})$ and extends the range
    $0<\gamma<1/2$ to $0<\gamma<1$. The states of distinct
    blocks need not be orthogonal, and no condition $q=o(N)$ is needed. As recorded
    in \cite{gap:mswc24_block_form_mixed_overlap}
    (Lemma~\ref{thm:ldp_approximation_error_overlapping_blocks}), with
    $\xi_{\mathrm{diag}}$ alone the bound is false.
\end{lemma}
\begin{proof}
    \leanok
    Specializing Lemma~\ref{thm:ldp_approximation_error_complex_weights} to
    $\mu_j=1$, where $\beta_j=1$, yields the bound.
\end{proof}

% **Local fix (rate of the overlapping blocks):** as above, for multiplicity one with real
% weights; docs/paper-gaps/mswc24_block_form_mixed_overlap.tex.
\begin{lemma}[Approximation error for overlapping blocks with weights]
    \label{thm:ldp_approximation_error_overlapping_weights}
    \lean{MPSTensor.exists_approximationError_le_overlappingBlockSum_weight,
        MPSTensor.exists_approximationError_le_mul_overlappingBlockSum_weight}
    \leanok
    \uses{def:ldp_non_normal_fixed_point, thm:ldp_approximation_error_complex_weights,
        def:correlation_length}
    Let $A^i=\bigoplus_{j=1}^bw_jA_j^i$ be the direct sum with real weights
    $w_j\geq0$, not all zero, of normal blocks $A_j$ in the gauge
    (\ref{eq:ldp_normal_gauge}), with positive definite fixed points $\rho_j$ of trace one,
    each placed on its own bond coordinates. Let $\lambda_2\in\C$ bound the moduli of the
    eigenvalues other than one of every transfer map $\E_{A_j}$ and the moduli of all
    eigenvalues of the mixed transfer maps $\E_{jj'}$, $j\neq j'$, let $\xi$ be its
    correlation length, and let $0<\gamma<1$. Form the approximating state
    $\widetilde\phi_N$ of Definition~\ref{def:ldp_non_normal_fixed_point} with the pairs
    of the $\rho_j$ on the bond coordinates of their blocks and $\beta_j=w_j^N$. There are
    constants $C,C'>0$ such that for all such weights, every block length $q$ and every
    $M\geq1$, with $N=qM$ and $y=Me^{-\gamma q/\xi}$,
    \begin{align*}
        \varepsilon(\widetilde\phi_N,\phi_N)\leq C'y
        \qquad\text{and}\qquad
        \varepsilon(\widetilde\phi_N,\phi_N)\leq Cy\,e^{Cy}.
    \end{align*}
    As recorded in \cite{gap:mswc24_block_form_mixed_overlap}, relative to
    \cite[Lemma~1$'$(ii)]{Malz2024Preparation} for blocks of multiplicity one with weights
    of possibly different moduli and with phases absorbed into the blocks, this replaces
    $\xi_{\mathrm{diag}}$ by
    $\max(\xi_{\mathrm{diag}},\xi_{\mathrm{off\text{-}diag}})$, extends $0<\gamma<1/2$ to
    $0<\gamma<1$ with constants that do not depend on the weights, and drops the
    normalization of the source that the largest weight is one.
\end{lemma}
\begin{proof}
    \leanok
    Specializing Lemma~\ref{thm:ldp_approximation_error_complex_weights} to
    $\mu_j=w_j$, where $\lvert\mu_j\rvert^N=w_j^N$, yields the bound.
\end{proof}

%% ==================================================
%% Subsection: Repeated blocks whose states overlap
%% Source: arXiv:2307.01696 (Malz, Styliaris, Wei, Cirac), Supplemental
%% Material, "Proof of Lemma 1 and extension to non-normal tensors",
%% eqs. (S2)-(S12) and Lemma 1'(ii).
%% Local fix: corrected fixed-point state for repeated blocks,
%% docs/paper-gaps/mswc24_repeated_block_corrected_state.tex; rate including
%% the mixed transfer maps, docs/paper-gaps/mswc24_block_form_mixed_overlap.tex.
%% The bound holds for all nonzero weights, with no factor depending on them,
%% docs/paper-gaps/mswc24_repeated_overlap_small_weight_sum.tex (resolved).
%% ==================================================

\subsection{Repeated blocks whose states overlap}\label{sec:ldp_repeated_overlapping_blocks}

The two corrections of the preceding subsections combine: for blocks with
multiplicities whose states on $q$ sites overlap, the corrected approximating state
of Definition~\ref{def:ldp_repeated_block_state} converges at the rate set by the
correlation lengths of the blocks and of their mixed transfer maps
\cite{gap:mswc24_repeated_block_corrected_state,gap:mswc24_block_form_mixed_overlap}.
The estimates of Lemma~\ref{lem:ldp_shared_overlap_estimates} use only that the
embeddings are isometries with orthogonal ranges, and apply with the isometries
$L_j$ of Lemma~\ref{thm:ldp_repeated_block_polar}.

\begin{lemma}[Positive part and overlaps for repeated overlapping blocks]
    \label{lem:ldp_repeated_overlapping_overlaps}
    \lean{MPSTensor.copyPosLimitRow, MPSTensor.copyPosLimit,
        MPSTensor.conjTranspose_copyIsometry_mul_copyPosLimit,
        MPSTensor.posSemidef_copyPosLimit, MPSTensor.copyPosLimit_mul_self,
        MPSTensor.exists_norm_polarPos_blockTensor_repeatedBlockSum_sub_le,
        MPSTensor.mixedMapLM_copyPosLimitRow_apply,
        MPSTensor.exists_norm_mpvOverlap_sub_bntWeight_le_of_norm_sub_copyPosLimitRow_le,
        MPSTensor.sum_star_copyApproxVector_mul_mpv,
        MPSTensor.sum_norm_sq_copyApproxVector_le}
    \leanok
    \uses{lem:ldp_shared_overlap_estimates, thm:ldp_repeated_block_polar,
        def:ldp_repeated_block_state, lem:ldp_overlapping_overlaps_weights,
        def:prep_fixed_point_tensor}
    With the notation of Lemma~\ref{thm:ldp_repeated_block_polar}, let the blocks be
    as in Lemma~\ref{lem:ldp_overlapping_positive_part}, with $c_j$ and $L_j$ the norms
    and isometries of the copies on $q$ sites, and define
    \begin{align*}
        P_\infty=\sum_jc_jL_j\bigl(\sqrt{\rho_j}^{\,\mathsf T}\otimes\Id\bigr)L_j^\dagger,
        \qquad
        L_j^\dagger P_\infty=\sum_k\mu_{j,k}^q\bigl(\sqrt{\rho_j}^{\,\mathsf T}\otimes\Id\bigr)
        K_{j,k}^\dagger.
    \end{align*}
    \begin{enumerate}
        \item $P_\infty$ is positive semidefinite, its square is
              $\sum_jc_j^2L_j(\rho_j^{\mathsf T}\otimes\Id)L_j^\dagger$, and for
              $0<\gamma<1/2$ there is $K$ such that for all weights with
              $\lvert\mu_{j,k}\rvert\leq1$ and every $q\geq1$ the positive part $P$ of the
              blocked map satisfies $\lVert P-P_\infty\rVert\leq Ke^{-\gamma q/\xi}$.
        \item The mixed transfer map of $L_j^\dagger P_\infty$, read as a tensor whose
              physical leg is a bond pair of block $j$, against the fixed-point tensor
              of $\rho_j$ is $X\mapsto\sum_k\mu_{j,k}^q\Tr(E_{j,k}^\dagger X)\,E_{j,k}\rho_j$.
        \item For a block $j$ with at least one copy there is $C_j>0$ such that for all
              weights with $\lvert\mu_{j,k}\rvert\leq1$, all $q$, $\delta$, $M\geq1$ and
              every $G$ with $\lVert G-L_j^\dagger P_\infty\rVert\leq\delta$, the overlap
              of $G$ with the fixed-point tensor of $\rho_j$ on $M$ sites is
              $\beta_j=\sum_k\mu_{j,k}^{qM}$ up to $C_j(M\delta)\,e^{C_jM\delta}$.
        \item The unnormalized corrected state has inner product
              $\sum_j\overline{\alpha_j}\braket{\phi_M(P_{\infty,\rho_j})}
                  {\phi_M(L_j^\dagger P)}$ with $\ket{\phi_N(A)}$, and for isometries
              $L_j$ with orthogonal ranges, $\beta\neq0$ and
              $\alpha_j=\beta_j/(\sum_l\lvert\beta_l\rvert^2)^{1/2}$ it has norm at most one.
    \end{enumerate}
\end{lemma}
\begin{proof}
    \leanok
    Part~(1): $c_jL_j^\dagger=\sum_k\mu_{j,k}^qK_{j,k}^\dagger$ and
    Lemma~\ref{lem:ldp_shared_overlap_estimates}(1) give the form of $L_j^\dagger P_\infty$
    and of $P_\infty^2$; Lemma~\ref{lem:ldp_shared_overlap_estimates}(4) with
    $c_j\leq(\sum_lm_l)^{1/2}$ and $2\gamma$ in place of $\gamma$, and
    Lemma~\ref{lem:ldp_shared_overlap_estimates}(5), give the rate. Part~(2): the letter of
    $(\sqrt\rho^{\,\mathsf T}\otimes\Id)K^\dagger$ at $(l,r)$ is
    $E\sqrt\rho\ket l\bra rE^\dagger$, and $\E_{P_{\infty,\rho}}(Y)=\Tr(Y)\rho$. Part~(3): the maps
    $R_k(X)=\Tr(E_{j,k}^\dagger X)E_{j,k}\rho_j$ satisfy $R_kR_{k'}=\delta_{kk'}R_k$ and
    $\Tr R_k=1$, so $(\sum_k\mu_{j,k}^qR_k)^n=\sum_k\mu_{j,k}^{qn}R_k$ for $n\geq1$ has norm
    at most $\sum_k\lVert R_k\rVert$ and its $M$-th power has trace $\beta_j$;
    Lemma~\ref{lem:ldp_shared_overlap_estimates}(8) applies. Part~(4) follows from
    Lemma~\ref{lem:ldp_shared_overlap_estimates}(7),
    $\braket{L_j^{\otimes M}\Omega_j}{\psi}=\braket{\Omega_j}{(L_j^\dagger)^{\otimes M}\psi}$,
    and the orthonormality of the states $L_j^{\otimes M}\ket{\Omega_j}$.
\end{proof}

\begin{lemma}[A Sylvester equation with positive coefficients]
    \label{lem:ldp_sylvester_bound}
    \lean{Matrix.norm_le_of_mul_add_mul_eq_mul}
    \leanok
    Let $P$ and $S$ be positive semidefinite $n\times n$ matrices, let $\lambda>0$ and $R$
    satisfy $\lambda RR^\dagger\leq P^2$, and let $X=PW_1+W_2S$ satisfy $PX+XS=RZ$. Then
    $\lVert X\rVert\leq n^2\lambda^{-1/2}\lVert Z\rVert$.
\end{lemma}
\begin{proof}
    \leanok
    In orthonormal eigenbases $P=\sum_ap_af_af_a^\dagger$ and $S=\sum_is_ie_ie_i^\dagger$ the
    equation reads $(p_a+s_i)\braket{f_a}{Xe_i}=\braket{R^\dagger f_a}{Ze_i}$. From
    $\lambda RR^\dagger\leq P^2$, $\lVert R^\dagger f_a\rVert\leq\lambda^{-1/2}p_a$, so
    $\lvert\braket{f_a}{Xe_i}\rvert\leq\lambda^{-1/2}\lVert Z\rVert$ whenever
    $p_a+s_i>0$, since $p_a/(p_a+s_i)\leq1$. When $p_a=s_i=0$ the entry vanishes, because
    $f_a^\dagger P=0$ and $Se_i=0$. The operator norm is at most the sum of the moduli of
    the $n^2$ entries.
\end{proof}

\begin{lemma}[Positive part relative to the block weights]
    \label{lem:ldp_relative_positive_part}
    \lean{MPSTensor.norm_polarPos_sub_mul_sum_le, MPSTensor.copyScaleInv,
        MPSTensor.relativeRow, MPSTensor.relativeRowLimit,
        MPSTensor.copyScaleInv_mul_copyIsometry, MPSTensor.relativeRow_eq_smul,
        MPSTensor.relativeRowLimit_self, MPSTensor.relativeRowLimit_of_ne,
        MPSTensor.exists_norm_relativeRow_sub_le}
    \leanok
    \uses{lem:ldp_sylvester_bound, lem:ldp_repeated_overlapping_overlaps,
        lem:ldp_shared_overlap_estimates, thm:ldp_repeated_block_polar}
    \begin{enumerate}
        \item Let $L_j$ be isometries with orthogonal ranges in dimension $n$, $c_j>0$,
              $C=\sum_jc_jL_jL_j^\dagger$ and $C^{-1}=\sum_jc_j^{-1}L_jL_j^\dagger$, let
              $h_j\geq0$ with $h_j^2=g_j\geq\lambda\Id$ for some $\lambda>0$, and let
              $\mathcal B=\widetilde{\mathcal B}C$. If
              $\Delta=\widetilde{\mathcal B}^\dagger\widetilde{\mathcal B}
                  -\sum_jL_jg_jL_j^\dagger$ satisfies $\lVert\Delta\rVert\leq\lambda/2$, the
              positive part $P$ of $\mathcal B$ satisfies
              $\lVert(P-S)C^{-1}\rVert\leq n^2(\lambda/2)^{-1/2}\lVert\Delta\rVert$ with
              $S=\sum_jc_jL_jh_jL_j^\dagger$.
        \item With the notation of Lemma~\ref{thm:ldp_repeated_block_polar} and
              Lemma~\ref{lem:ldp_repeated_overlapping_overlaps}, $C^{-1}L_j=c_j^{-1}L_j$, the
              block $L_j^\dagger PC^{-1}L_{j'}$ equals $c_{j'}^{-1}L_j^\dagger PL_{j'}$, and the
              block $L_j^\dagger P_\infty C^{-1}L_{j'}$ is
              $\sqrt{\rho_j}^{\,\mathsf T}\otimes\Id$ for $j=j'$ and zero otherwise.
        \item Let the blocks be normal in the gauge (\ref{eq:ldp_normal_gauge}) with positive
              definite fixed points $\rho_j$ of trace one, let $\lambda_2$ with
              $\lvert\lambda_2\rvert<1$ bound the moduli of the eigenvalues other than one of
              every $\E_{A_j}$ and of all eigenvalues of the mixed transfer maps
              $\E_{jj'}$, $j\neq j'$, and let $0<\gamma<1$ with $e^{-\gamma/\xi}<1$. There are
              $K$ and $q_0$ such that for all nonzero weights and all $q\geq\max(q_0,1)$,
              $\lVert(P-P_\infty)C^{-1}\rVert\leq Ke^{-\gamma q/\xi}$, and every block
              $L_j^\dagger PC^{-1}L_{j'}$ is $L_j^\dagger P_\infty C^{-1}L_{j'}$ up to
              $Ke^{-\gamma q/\xi}$.
    \end{enumerate}
\end{lemma}
\begin{proof}
    \leanok
    (1) The positive part satisfies $P^2=\mathcal B^\dagger\mathcal B
        =C\widetilde{\mathcal B}^\dagger\widetilde{\mathcal B}C\geq(\lambda/2)C^2$, because
    $\widetilde{\mathcal B}^\dagger\widetilde{\mathcal B}\geq\sum_jL_jg_jL_j^\dagger
        -(\lambda/2)\Id$ and $g_j\geq\lambda\Id$. The matrix $X=(P-S)C^{-1}$ solves
    $PX+XS=C\Delta\Pi$ with $\Pi=CC^{-1}$, since $S^2=C\bigl(\sum_jL_jg_jL_j^\dagger\bigr)C$
    and $S$ commutes with $C^{-1}$. It has the form $PW_1+W_2S$ with $W_1=C^{-1}$ and
    $W_2=-C^{-1}$, and $\lVert\Pi\rVert\leq1$; Lemma~\ref{lem:ldp_sylvester_bound}
    applies with $R=C$ and $Z=\Delta\Pi$. (2) follows from the orthogonality of the ranges of the $L_j$ and
    the form of $P_\infty$ in Lemma~\ref{lem:ldp_repeated_overlapping_overlaps}. (3) By
    Lemma~\ref{thm:ldp_repeated_block_polar} the blocked map is
    $\mathcal B=\widetilde{\mathcal B}C$ with $\widetilde{\mathcal B}=\sum_j\mathcal B_jL_j^\dagger$,
    and Lemma~\ref{lem:ldp_shared_overlap_estimates}(4) with unit coefficients gives
    $\lVert\Delta\rVert\leq Ke^{-\gamma q/\xi}$ for
    $g_j=\rho_j^{\mathsf T}\otimes\Id$ and $h_j=\sqrt{\rho_j}^{\,\mathsf T}\otimes\Id$, which
    is at most $\lambda/2$ for $q\geq q_0$, where $\lambda>0$ bounds every
    $\rho_j^{\mathsf T}\otimes\Id$ from below. Part~(1) gives the first bound, and the
    blocks follow from part~(2) since the $L_j$ have norm one.
\end{proof}

\begin{lemma}[Overlaps relative to the block weights]
    \label{lem:ldp_relative_overlaps}
    \lean{MPSTensor.CopyWeights.pow,
        MPSTensor.ofPhysicalMatrixLM_mul_conjTranspose_pairEmbedding,
        MPSTensor.polarPos_mul_sum_copyIsometry,
        MPSTensor.conjTranspose_copyIsometry_mul_polarPos_eq_sum,
        MPSTensor.ofPhysicalMatrixLM_conjTranspose_copyIsometry_mul_polarPos,
        MPSTensor.mpvOverlap_conjTranspose_copyIsometry_mul_polarPos,
        MPSTensor.exists_norm_mpvOverlap_sub_ite_le}
    \leanok
    \uses{lem:ldp_relative_positive_part, lem:ldp_repeated_overlapping_overlaps,
        lem:ldp_shared_overlap_estimates, def:prep_fixed_point_tensor}
    With the notation of Lemma~\ref{lem:ldp_relative_positive_part}, let $q\geq1$ and
    $M\geq1$, with $N=qM$.
    \begin{enumerate}
        \item The letters of the tensor read from $YK^\dagger$, for a coordinate isometry $K$
              placing a bond pair, are those of the tensor read from $Y$, placed on the bond
              coordinates of $K$.
        \item $L_j^\dagger P=\sum_{j'}\sum_k\mu_{j',k}^q\,(L_j^\dagger PC^{-1}L_{j'})
                  K_{j',k}^\dagger$, so the tensor read from $L_j^\dagger P$ is the direct
              sum over the copies $(j',k)$, with the weights $\mu_{j',k}^q$, of the tensors
              read from the blocks $L_j^\dagger PC^{-1}L_{j'}$.
        \item For every tensor $F$ on the bond pairs of block $j$,
              $\braket{\phi_M(F)}{\phi_M(L_j^\dagger P)}
                  =\sum_{j'}\beta_{j'}\braket{\phi_M(F)}{\phi_M(L_j^\dagger PC^{-1}L_{j'})}$
              with $\beta_{j'}=\sum_k\mu_{j',k}^N$.
        \item For positive semidefinite $\rho_j$ of trace one and blocks $j,j'$ there is
              $C>0$ such that for all nonzero weights, all $q$, $\delta$, $M\geq1$, and every
              $G$ with $\lVert G-L_j^\dagger P_\infty C^{-1}L_{j'}\rVert\leq\delta$, the
              overlap of the tensor read from $G$ with the fixed-point tensor of $\rho_j$ on
              $M$ sites is $\delta_{jj'}$ up to $C(M\delta)\,e^{CM\delta}$.
    \end{enumerate}
\end{lemma}
\begin{proof}
    \leanok
    (1) is a computation of entries. (2) $P$ vanishes off the ranges of the $L_j$, so
    $L_j^\dagger P=\sum_{j'}L_j^\dagger PL_{j'}L_{j'}^\dagger$, and
    $c_{j'}L_{j'}^\dagger=\sum_k\mu_{j',k}^qK_{j',k}^\dagger$; part~(1) gives the direct sum.
    (3) The state of a direct sum of copies with weights $\mu_{j',k}^q$ on $M$ sites is
    $\sum_{j'}(\sum_k\mu_{j',k}^{qM})$ times the state of block $j'$. (4) By
    Lemma~\ref{lem:ldp_relative_positive_part}(2) the limit is
    $\sqrt{\rho_j}^{\,\mathsf T}\otimes\Id$ for $j=j'$, read as the fixed-point tensor of
    $\rho_j$, whose mixed transfer map against itself is the idempotent
    $X\mapsto\Tr(X)\rho_j$ of trace one, and zero otherwise, whose mixed transfer map
    vanishes. Lemma~\ref{lem:ldp_shared_overlap_estimates}(8) applies.
\end{proof}

\begin{lemma}[Second-order deficit of the relative diagonal coefficient]
    \label{lem:ldp_relative_diagonal_deficit}
    \lean{MPSTensor.relative_column_gram, MPSTensor.sum_norm_relative_column_sq,
        MPSTensor.relative_diagonal_nonneg, MPSTensor.relative_diagonal_geometry,
        MPSTensor.exists_relative_diagonal_deficit_le,
        MPSTensor.exists_norm_one_sub_relative_diagonal_le}
    \leanok
    \uses{lem:ldp_relative_positive_part, lem:ldp_relative_overlaps,
        def:prep_fixed_point_tensor}
    With the notation of Lemma~\ref{lem:ldp_relative_positive_part}, let $q\geq1$, let
    $H_{kj}=L_k^\dagger PC^{-1}L_j$ be the relative blocks, and for a positive
    semidefinite $\sigma$ let $\lVert X\rVert_\sigma^2=\Tr\sum_iX^i\sigma(X^i)^\dagger$
    for a tensor $X$ read from a matrix. Let $\sigma_j$ be a positive semidefinite fixed
    point of $\E_{A_j}$ of trace one, $F_j$ its fixed-point tensor, and
    $\alpha_j=\Tr\E_{H_{jj}F_j}(\sigma_j)$.
    \begin{enumerate}
        \item $\sum_kH_{kj}^\dagger H_{kj}=\mathcal B_j^\dagger\mathcal B_j$, where
              $\mathcal B_j$ is the physical matrix of the $q$-blocked tensor $A_j$, and
              $\sum_k\lVert H_{kj}\rVert_{\sigma_j}^2=1$.
        \item $\alpha_j$ is real, $0\leq\alpha_j\leq1$, and
              $1-\alpha_j=\tfrac12\bigl(\lVert H_{jj}-F_j\rVert_{\sigma_j}^2
                  +\sum_{k\neq j}\lVert H_{kj}\rVert_{\sigma_j}^2\bigr)$.
        \item There is $C>0$ such that for all nonzero weights, all $q\geq1$ and all $\delta$,
              if every $H_{kj}$ is within $\delta$ in operator norm of its limit from
              Lemma~\ref{lem:ldp_relative_positive_part}(2), then $\lvert\alpha_j\rvert\leq1$
              and $\lvert1-\alpha_j\rvert\leq C\delta^2$.
        \item Under the hypotheses of Lemma~\ref{lem:ldp_relative_positive_part}(3), for
              every $0<\gamma<1$ there are $C>0$ and $q_0$ such that for all nonzero
              weights and all $q\geq\max(q_0,1)$,
              $\lvert1-\alpha_j\rvert\leq Ce^{-2\gamma q/\xi}$.
    \end{enumerate}
\end{lemma}
\begin{proof}
    \leanok
    (1) $P$ vanishes off the ranges of the $L_k$, which are orthogonal, so
    $\sum_kH_{kj}^\dagger H_{kj}=c_j^{-2}L_j^\dagger P^2L_j=\mathcal B_j^\dagger\mathcal B_j$.
    The squared $\sigma_j$-norms of the $H_{kj}$ add up to
    $\Tr\E_{A_j}^q(\sigma_j)=\Tr\sigma_j=1$. (2) The coefficient is the inner product of
    the vectors $(H_{jj}^i\sqrt{\sigma_j})_i$ and $(F_j^i\sqrt{\sigma_j})_i$; both have
    norm at most one, and $\lVert F_j\rVert_{\sigma_j}=1$. It is nonnegative since $P\geq0$.
    Expanding $\lVert H_{jj}-F_j\rVert_{\sigma_j}^2$ and using (1) gives the identity.
    (3) Each $\sigma_j$-norm is bounded by a dimensional constant times the operator norm;
    the off-diagonal limits vanish. (4) Combining (3) with
    Lemma~\ref{lem:ldp_relative_positive_part}(3) gives the exponential bound.
\end{proof}

\begin{lemma}[Second-order overlaps of the relative rows]
    \label{lem:ldp_relative_overlaps_second_order}
    \lean{MPSTensor.exists_norm_relative_row_overlap_sub_ite_le_sq}
    \leanok
    \uses{lem:ldp_relative_diagonal_deficit, lem:ldp_relative_overlaps,
        lem:ldp_shared_overlap_estimates}
    With the notation of Lemma~\ref{lem:ldp_relative_diagonal_deficit}, let $\sigma_k$ be a
    positive semidefinite fixed point of $\E_{A_k}$ of trace one for every block $k$, and
    let $j,j'$ be blocks. There is $C>0$ such that for all nonzero weights, all $q\geq1$,
    all $\delta$ and all $M\geq2$: if every relative block $H_{kl}$ is within $\delta$ of
    its limit and $CM\delta^2<1$, then the overlap of the tensor read from $H_{jj'}$ with
    $F_j$ on $M$ sites is $\delta_{jj'}$ up to $CM\delta^2e^{CM\delta^2}$. No quadratic
    bound is asserted for $M=1$.
\end{lemma}
\begin{proof}
    \leanok
    For $j=j'$ the mixed transfer map of $H_{jj}$ and $F_j$ is within $O(\delta)$ of the
    idempotent $X\mapsto\Tr(X)\sigma_j$, and its trace against $\sigma_j$ is $\alpha_j$,
    which is $1$ up to $O(\delta^2)$ by
    Lemma~\ref{lem:ldp_relative_diagonal_deficit}(3); the idempotent power estimate of
    Lemma~\ref{lem:ldp_shared_overlap_estimates} gives the bound. For $j\neq j'$ the
    limit of $H_{jj'}$ is zero, so the mixed transfer map has norm $O(\delta)$ and its
    $M$-th power has norm $O(\delta^M)$, which is $O(M\delta^2)$ for $M\geq2$ under the
    smallness condition.
\end{proof}

\begin{lemma}[Approximation error for repeated overlapping blocks]
    \label{thm:ldp_approximation_error_repeated_overlapping}
    \lean{MPSTensor.exists_approximationError_le_repeatedOverlappingBlockSum,
        MPSTensor.exists_approximationError_le_mul_repeatedOverlappingBlockSum,
        MPSTensor.one_sub_norm_div_le_of_norm_sub_le,
        MPSTensor.one_sub_le_mul_mul_exp_div_of_one_le,
        MPSTensor.one_sub_le_mul_mul_exp_div_of_sum_le}
    \leanok
    \uses{lem:ldp_repeated_overlapping_overlaps, lem:ldp_shared_overlap_estimates,
        lem:ldp_relative_positive_part, lem:ldp_relative_overlaps,
        def:ldp_repeated_block_state, def:correlation_length}
    With the notation of Lemma~\ref{thm:ldp_repeated_block_polar}, let the blocks $A_j$
    be normal in the gauge (\ref{eq:ldp_normal_gauge}) with positive definite fixed points
    $\rho_j$ of trace one, let $\lambda_2\in\C$ bound the moduli of the eigenvalues other
    than one of every transfer map $\E_{A_j}$ and the moduli of all eigenvalues of the
    mixed transfer maps $\E_{jj'}$, $j\neq j'$, let $\xi$ be its correlation length, and
    let $0<\gamma<1$. There are $C,C'>0$ such that for all nonzero complex weights, every
    block length $q$ and every $M\geq1$ with $N=qM$ and $(\beta_1,\dots,\beta_b)\neq0$, the
    corrected approximating state of Definition~\ref{def:ldp_repeated_block_state}
    satisfies, with $y=Me^{-\gamma q/\xi}$,
    \begin{align*}
        \varepsilon(\widetilde\phi_N,\phi_N)\leq Cy\,e^{Cy},\qquad
        \varepsilon(\widetilde\phi_N,\phi_N)\leq C'y.
    \end{align*}
    Relative to \cite[Lemma~1$'$(ii)]{Malz2024Preparation}, this establishes the
    error bound for the corrected state with $\xi_{\mathrm{diag}}$ replaced by
    $\max(\xi_{\mathrm{diag}},\xi_{\mathrm{off\text{-}diag}})$ and with $0<\gamma<1$ in
    place of $0<\gamma<1/2$. As recorded in
    \cite{gap:mswc24_repeated_overlap_small_weight_sum}, the constants do not depend on
    the weights: the sums $\beta_j$ may cancel and $\sum_j\lvert\beta_j\rvert^2$ may be
    arbitrarily small, while the normalization $\lvert\mu_{j,k}\rvert\leq1$ of the source
    is not used. The states of distinct blocks need not be orthogonal, and no condition
    $q=o(N)$ is needed.
\end{lemma}
\begin{proof}
    \leanok
    Let $q_0$ be the threshold of Lemma~\ref{lem:ldp_relative_positive_part}(3). If
    $e^{-\gamma/\xi}\geq1$, or $q=0$, or $q<q_0$, then $C$ can be chosen with $Cy\geq1$, and
    $\varepsilon\leq1$ gives the bound. Otherwise $\lvert\lambda_2\rvert<1$. By
    Lemma~\ref{lem:ldp_repeated_overlapping_overlaps}(4) the unnormalized corrected state
    has inner product $S/\sqrt b$ with $\ket{\phi_N(A)}=\sum_j\beta_j\ket{\phi_N(A_j)}$,
    where $b=\sum_j\lvert\beta_j\rvert^2$, $S=\sum_j\overline{\beta_j}z_j$ and
    $z_j=\braket{\phi_M(P_{\infty,\rho_j})}{\phi_M(L_j^\dagger P)}$, and it has norm at most
    one, so $\lvert\braket{\widetilde\phi_N}{\phi_N}\rvert\geq\lvert S\rvert/(\sqrt bt)$ with
    $t=\lVert\phi_N(A)\rVert$. By Lemma~\ref{lem:ldp_relative_overlaps}(3),
    $z_j=\sum_{j'}\beta_{j'}w_{jj'}$, and by Lemma~\ref{lem:ldp_relative_positive_part}(3)
    and Lemma~\ref{lem:ldp_relative_overlaps}(4),
    $\lvert w_{jj'}-\delta_{jj'}\rvert\leq C_{jj'}y\,e^{C_{jj'}y}$. Hence
    \begin{align*}
        \frac Sb-1=\frac1b\sum_{j,j'}\overline{\beta_j}\beta_{j'}(w_{jj'}-\delta_{jj'}),
        \qquad
        \Bigl\lvert\frac Sb-1\Bigr\rvert\leq\sum_{j,j'}C_{jj'}y\,e^{C_{jj'}y},
    \end{align*}
    since $\lvert\beta_j\beta_{j'}\rvert\leq b$. Lemma~\ref{lem:ldp_shared_overlap_estimates}(6)
    with $R=\sqrt b$ gives $\lvert t^2/b-1\rvert\leq Ke^{-\gamma N/\xi}\leq Ky$.
    Two elementary inequalities finish the proof. First, if $\lvert S'-1\rvert\leq\delta$ and
    $\lvert t'^2-1\rvert\leq\eta$ with $t'\geq0$, then
    $1-\lvert S'\rvert/t'\leq\delta+\eta$: this is trivial when $\delta+\eta\geq1$, and
    otherwise $t'>0$, $\lvert S'\rvert\geq1-\delta$ and $t'\leq(1+t'^2)/2\leq1+\eta/2$.
    Second, let $0<s\leq1$. If $r\geq0$, $C\geq1$ and $u\geq1$, then
    $1-r\leq Cu\,e^{Cu}/s$; and if $0\leq x\leq1$, $M\geq1$, $u=Mx^q$, $K'\geq0$,
    $C_j\geq0$, $e_j\leq C_ju\,e^{C_ju}$ and $1-r\leq s^{-1}\sum_je_j+K'x^{Mq}$, then
    $1-r\leq Cu\,e^{Cu}/s$ with $C=\sum_jC_j+K'+1$. Applied with $S'=S/b$,
    $t'=t/\sqrt b$, $s=1$ and $x=e^{-\gamma/\xi}$, they give $\varepsilon\leq Cy\,e^{Cy}$,
    and the second form follows because $\varepsilon\leq1$.
\end{proof}

%% ==================================================
%% Subsection: Weights with a phase
%% Source: arXiv:2307.01696 (Malz, Styliaris, Wei, Cirac), Supplemental
%% Material, "Proof of Lemma 1 and extension to non-normal tensors",
%% eqs. (S2)-(S12) and Lemma 1'(ii).
%% False source (eqs. (S5)-(S7)): with weights that are not nonnegative reals, the
%% coefficients beta_j = mu_j^N of eq. (S4), correct in eq. (S3), count the phase of
%% the weight twice in eqs. (S6) and (S7);
%% docs/paper-gaps/mswc24_block_form_mixed_overlap.tex.
%% ==================================================

\subsection{Weights with a phase}\label{sec:ldp_complex_weights}

For blocks of multiplicity one write $\mu_j=\lvert\mu_j\rvert u_j$ with
$\lvert u_j\rvert=1$. The blocked map is $B=B_0U$, where $B_0$ is the blocked map
for the weights $\lvert\mu_j\rvert$ and the diagonal unitary $U$ multiplies the
bond coordinates of block $j$ by $u_j^q$. Hence
$B^\dagger B=U^\dagger B_0^\dagger B_0U$, and the polar decompositions satisfy
\begin{align*}
    P=U^\dagger P_0U,\qquad V=V_0U.
\end{align*}
The diagonal blocks of $P$ are those of $P_0$ and do not depend on the phases,
contrary to the block form \cite[eq.~(S5)]{Malz2024Preparation}; only the mixed
blocks of $P$, which vanish for orthogonal blocks, carry the phases
$\overline{u_j}^{\,q}u_{j'}^q$. The product
$\ket{\Omega_j}=\bigotimes_{k=1}^M\ket{\omega_j}_{R_kL_{k+1}}$ of the pairs of
block $j$ lies on the bond coordinates of block $j$, so
$V^{\otimes M}\ket{\Omega_j}=u_j^NV_0^{\otimes M}\ket{\Omega_j}$ on $N=qM$
sites. The coefficients $\beta_j=\mu_j^N$ of
\cite[eq.~(S4)]{Malz2024Preparation} are correct in the expansion
\cite[eq.~(S3)]{Malz2024Preparation} of $\ket{\phi_N}$, but under
$V^{\otimes M}$ in \cite[eqs.~(S6) and (S7)]{Malz2024Preparation} they count the
phase $u_j^N$ twice, and already for orthogonal blocks the approximating state
does not converge (Theorem~\ref{thm:ldp_approximation_error_phase_weights}). As
recorded in \cite{gap:mswc24_block_form_mixed_overlap}, with
$\beta_j=\lvert\mu_j\rvert^N$ it converges at the rate of
Theorem~\ref{thm:ldp_approximation_error_complex_weights}.

% **False source result:** refutes Lemma 1'(ii) of arXiv:2307.01696 for two
% orthogonal blocks of multiplicity one with weights 1 and -1;
% docs/paper-gaps/mswc24_block_form_mixed_overlap.tex. As for
% thm:ldp_approximation_error_repeated_block, the first linked declaration
% certifies the hypotheses of Lemma 1'(ii) for the displayed tensor.
\begin{theorem}[The approximating state for weights with a phase]
    \label{thm:ldp_approximation_error_phase_weights}
    \lean{MPSTensor.isBNTCanonicalForm_and_not_approximationError_le_phaseBlock,
        MPSTensor.phaseBlockTensor, MPSTensor.phaseBlockSector,
        MPSTensor.isBNTCanonicalForm_phaseBlockSector,
        MPSTensor.reindex_toTensor_phaseBlockSector,
        MPSTensor.embeddedFixedPointPair_phaseBlockSector,
        MPSTensor.blockSum_phaseBlock,
        MPSTensor.bntWeight_phaseBlockWeight,
        MPSTensor.diagGram_phaseBlockDiag,
        MPSTensor.nonNormalApproxOverlap_phaseBlockTensor_of_norm,
        MPSTensor.nonNormalApproxOverlap_phaseBlockTensor,
        MPSTensor.nonNormalApproxOverlap_phaseBlockTensor_norm,
        MPSTensor.one_sub_norm_nonNormalApproxOverlap_phaseBlockTensor_neg_one,
        MPSTensor.not_approximationError_le_phaseBlockTensor,
        MPSTensor.phaseBlockDiag, MPSTensor.phaseBlockWeight,
        MPSTensor.ghzAmplitude_one_vec, MPSTensor.polarPos_phaseBlockTensor,
        MPSTensor.polarSupport_phaseBlockTensor, MPSTensor.coeff_phaseBlockSector,
        MPSTensor.repeatedBlockBasisSector,
        MPSTensor.hasBNTSectorData_repeatedBlockBasisSector,
        MPSTensor.isBNTCanonicalForm_repeatedBlockBasisSector,
        MPSTensor.phaseBlockCopyCoord,
        MPSTensor.totalDim_phaseBlockSector, MPSTensor.bijective_phaseBlockCopyCoord,
        MPSTensor.phaseBlockEquiv, MPSTensor.phaseBlockEquiv_symm_apply}
    \leanok
    \uses{def:ldp_non_normal_fixed_point, def:ldp_sector_fixed_point_state,
        def:correlation_length, def:sector_bnt_cf, def:nt_cpsv,
        def:cpsv_cfii}
    Let $\mu\in\C$ with $\lvert\mu\rvert=1$, and let $A^0=\operatorname{diag}(1,0)$
    and $A^1=\operatorname{diag}(0,\mu)$, the direct sum $A_1\oplus\mu A_2$ of the
    one-dimensional normal blocks $A_1=(1,0)$ and $A_2=(0,1)$, each of multiplicity
    one, with weights $1$ and $\mu$. After ordering its bond coordinates this is the
    canonical form of a basis of normal tensors in canonical form~II with weights of
    modulus one, so $A$ satisfies every hypothesis of
    \cite[Lemma~1$'$(ii)]{Malz2024Preparation}. For every $q\geq1$ and $M\geq1$, with
    $N=qM$, the approximating state of
    Definition~\ref{def:ldp_non_normal_fixed_point}, with the pairs of the fixed
    point $1$ of each block, satisfies
    \begin{align}
        \braket{\widetilde\phi_N}{\phi_N}=\frac{1+\overline\mu^{\,N}}2
        \quad\text{for }\beta=(1,\mu^N),\qquad
        \braket{\widetilde\phi_N}{\phi_N}=1
        \quad\text{for }\beta=(1,\lvert\mu\rvert^N)=(1,1).
        \label{eq:ldp_phase_weights_overlap}
    \end{align}
    More generally, for coefficients $\alpha$ with
    $\lvert\alpha_1\rvert^2+\lvert\alpha_2\rvert^2=1$ the overlap is
    $(\overline{\alpha_1}+\overline{\alpha_2})/\sqrt2$. For $\mu=-1$ and odd $q$ and
    $M$ the approximating state with $\beta=(1,\mu^N)$ is orthogonal to $\ket{\phi_N}$,
    so $\varepsilon(\widetilde\phi_N,\phi_N)=1$. The transfer maps of both blocks have
    no eigenvalue other than one, so every $\lambda_2$ bounds their subleading
    eigenvalues. For $\mu=-1$, every $\lambda_2$ with $0<\lvert\lambda_2\rvert<1$, every
    $\gamma>0$ and every constant $C$ there is an odd $M$ such that, with $q=M$,
    $N=M^2$ and $y=(N/q)e^{-\gamma q/\xi_{\mathrm{diag}}}$,
    $\varepsilon(\widetilde\phi_N,\phi_N)>Cy\,e^{Cy}$. Because the sequence $q=M$
    satisfies $q=o(N)$, as recorded in \cite{gap:mswc24_block_form_mixed_overlap},
    the bound (\ref{eq:ldp_approximation_error_general}) fails for orthogonal blocks of
    multiplicity one when a weight is not a nonnegative real number.
\end{theorem}

\begin{proof}
    \leanok
    \uses{thm:ldpolar_decomposition, thm:ldp_polar_uniqueness}
    The hypotheses are checked as in the proof of
    Theorem~\ref{thm:ldp_approximation_error_overlapping_blocks}: both blocks are
    one-dimensional with $\sum_i\lvert a_i\rvert^2=1$, their states $\ket{0\cdots0}$
    and $\ket{1\cdots1}$ are linearly independent for $N\geq1$, and
    $A_1^1=0\neq A_2^1$. These checks hold for the blocks $A_1$ and $A_2$ with any
    multiplicities and any nonzero weights $\lvert\mu_{j,k}\rvert\leq1$, one of modulus
    one, and are shared with
    Theorem~\ref{thm:ldp_approximation_error_repeated_block}.

    For $q\geq1$ the Gram matrix of the blocked map on the diagonal coordinates is the
    identity, since $\lvert\mu\rvert^{2q}=1$ and the off-diagonal entries are
    $0^q=0$. So $P$ and the support projector both equal $JJ^\dagger$, where $J$ places
    a coordinate $e$ on the pair $(e,e)$. As in the proof of
    Theorem~\ref{thm:ldp_approximation_error_overlapping_blocks}, the overlap of the
    unnormalized approximating vector with the unnormalized target is
    $\overline{\alpha_1}+\overline{\alpha_2}$, the vector has squared norm
    $\lvert\alpha_1\rvert^2+\lvert\alpha_2\rvert^2=1$, and
    $\lVert\phi_N(A)\rVert^2=2$. For $\beta=(1,\mu^N)$,
    $\alpha=(1,\mu^N)/\sqrt2$, and for $\beta=(1,1)$, $\alpha=(1,1)/\sqrt2$, which gives
    (\ref{eq:ldp_phase_weights_overlap}). For $\mu=-1$ and $N$ odd,
    $1+\overline\mu^{\,N}=0$. Finally $Cy\,e^{Cy}$ with $y=Mr^M$ and
    $r=e^{-\gamma/\xi_{\mathrm{diag}}}<1$ is eventually below one.
\end{proof}

% **False source (eqs. (S5)-(S7), complex weights):** beta_j = |mu_j|^N in place
% of the coefficients mu_j^N that eqs. (S6) and (S7) take from eq. (S4); Local
% fix of the rate as in thm:ldp_approximation_error_overlapping_weights;
% docs/paper-gaps/mswc24_block_form_mixed_overlap.tex.
\begin{theorem}[Approximation error for complex weights]
    \label{thm:ldp_approximation_error_complex_weights}
    \lean{MPSTensor.exists_approximationError_le_overlappingBlockSum_complexWeight,
        MPSTensor.exists_approximationError_le_mul_overlappingBlockSum_complexWeight,
        MPSTensor.CopyWeights.single, MPSTensor.repeatedBlockSum_single,
        MPSTensor.nonNormalApproxOverlap_blockSum_eq_copyApproxOverlap,
        MPSTensor.blockSum_comp_of_eq_zero,
        MPSTensor.nonNormalApproxOverlap_comp_of_eq_zero,
        MPSTensor.ghzAmplitude_comp_of_eq_zero}
    \leanok
    \uses{def:ldp_non_normal_fixed_point, def:correlation_length,
        def:ldp_repeated_block_sum, def:ldp_one_copy_weight, lem:ldp_one_copy_state,
        thm:ldp_approximation_error_repeated_overlapping}
    Let $A^i=\bigoplus_{j=1}^b\mu_jA_j^i$ be the direct sum with complex weights $\mu_j$
    of normal blocks $A_j$ in the gauge (\ref{eq:ldp_normal_gauge}), with positive
    definite fixed points $\rho_j$ of trace one, each placed on its own bond coordinates.
    Let $\lambda_2\in\C$ bound the moduli of the eigenvalues other than one of every
    transfer map $\E_{A_j}$ and the moduli of all eigenvalues of the mixed transfer maps
    $\E_{jj'}$, $j\neq j'$, let $\xi$ be its correlation length, and let $0<\gamma<1$.
    Form the approximating state $\widetilde\phi_N$ of
    Definition~\ref{def:ldp_non_normal_fixed_point} with the pairs of the $\rho_j$ on the
    bond coordinates of their blocks and $\beta_j=\lvert\mu_j\rvert^N$. There are
    constants $C,C'>0$ such that for all weights $\mu$ not all zero, every block length
    $q$ and every $M\geq1$, with $N=qM$ and $y=Me^{-\gamma q/\xi}$,
    \begin{align*}
        \varepsilon(\widetilde\phi_N,\phi_N)\leq C'y
        \qquad\text{and}\qquad
        \varepsilon(\widetilde\phi_N,\phi_N)\leq Cy\,e^{Cy}.
    \end{align*}
    Relative to \cite[Lemma~1$'$(ii)]{Malz2024Preparation} for blocks of multiplicity one,
    this replaces the coefficients $\beta_j=\mu_j^N$ that
    \cite[eqs.~(S6) and (S7)]{Malz2024Preparation} take from
    \cite[eq.~(S4)]{Malz2024Preparation} by $\beta_j=\lvert\mu_j\rvert^N$, replaces
    $\xi_{\mathrm{diag}}$ by
    $\max(\xi_{\mathrm{diag}},\xi_{\mathrm{off\text{-}diag}})$, and extends the range
    $0<\gamma<1/2$ to $0<\gamma<1$. The constants do not depend on the weights, and the
    normalization $\lvert\mu_j\rvert\leq1$ with one weight of modulus one is not needed.
    As recorded in \cite{gap:mswc24_block_form_mixed_overlap}
    (Theorem~\ref{thm:ldp_approximation_error_phase_weights}), with $\beta_j=\mu_j^N$
    the bound is false.
\end{theorem}

\begin{proof}
    \leanok
    If $q=0$ then $y=M\geq1$ and the bound holds with $C\geq1$, since
    $\varepsilon\leq1$; let $q\geq1$. A block of weight zero contributes nothing to $A$,
    and its coefficient $\beta_j=0$ removes its pairs from the approximating state; the
    normalization $\sum_l\lvert\beta_l\rvert^2$ is unchanged. So the blocks of weight zero
    may be dropped, and for each of the finitely many sets of remaining blocks it
    suffices to prove the bound for nonzero weights, with a constant that does not
    depend on them; the largest of these constants serves for all weights. For nonzero
    weights, $A$ is the direct sum of Definition~\ref{def:ldp_repeated_block_sum} with
    one copy of each block. By Definition~\ref{def:ldp_one_copy_weight} the coefficients
    for the pairs on the one copy are $\beta'_j=\lvert\mu_j\rvert^N$, and by
    Lemma~\ref{lem:ldp_one_copy_state}(6) the overlap of $\widetilde\phi_N$ with
    $\ket{\phi_N}$ is that of the corrected state with the coefficients $\mu_j^N$, not all
    zero. Theorem~\ref{thm:ldp_approximation_error_repeated_overlapping} bounds it, with a
    constant independent of the weights.
\end{proof}

%% ==================================================
%% Copies with weights of different phases.
%% False source (eq. (S7), coefficients): with the pairs of each block on one
%% of its copies, the coefficients beta_j of eq. (S4) are replaced by
%% beta_j (c_j / mu_{j,k_j}^q)^M;
%% docs/paper-gaps/mswc24_repeated_block_corrected_state.tex.
%% Scope restriction (orthogonal blocks) and Local fix (rate of the overlapping
%% blocks), as for thm:ldp_approximation_error_repeated_orthogonal_blocks and
%% thm:ldp_approximation_error_repeated_overlapping;
%% docs/paper-gaps/mswc24_block_form_mixed_overlap.tex.
%% ==================================================

\subsubsection{Copies with weights of different phases}
\label{sec:ldp_copy_phases}

When a block has several copies whose weights have different phases, the phases
cannot be moved into the block. For the corrected state of
Definition~\ref{def:ldp_repeated_block_state} the coefficients are those of
\cite[eq.~(S4)]{Malz2024Preparation}, $\beta_j=\sum_k\mu_{j,k}^N$, because the
isometry $L_j$ of Definition~\ref{def:ldp_repeated_block_sum} carries the phases
$\overline{\mu_{j,k}^q}/c_j$ of the copies:
Theorems~\ref{thm:ldp_approximation_error_repeated_orthogonal_blocks}
and~\ref{thm:ldp_approximation_error_repeated_overlapping} hold for complex weights.
As recorded in \cite{gap:mswc24_repeated_block_corrected_state}, neither
$\sum_k\lvert\mu_{j,k}\rvert^N$, the coefficient of
Theorem~\ref{thm:ldp_approximation_error_complex_weights} for one copy, nor $c_j^M$
works in their place. The
source's construction itself, with the pairs of each block on one of its copies,
converges once its coefficients are changed: the partial isometry $V$ sends the
pairs placed on one copy to a multiple of the pairs placed by $L_j$.

\begin{definition}[Coefficients for the pairs on one copy]
    \label{def:ldp_one_copy_weight}
    \lean{MPSTensor.oneCopyWeight, MPSTensor.oneCopyWeight_of_subsingleton,
        MPSTensor.oneCopyWeight_ne_zero}
    \leanok
    \uses{def:ldp_repeated_block_sum, thm:ldp_repeated_block_polar}
    With the notation of Definition~\ref{def:ldp_repeated_block_sum}, choose a copy
    $k_j$ of every block $j$, and let $N=qM$. The coefficient of block $j$ is
    \begin{align*}
        \beta'_j=\beta_j\Bigl(\frac{c_j}{\mu_{j,k_j}^q}\Bigr)^M,\qquad
        \beta_j=\sum_k\mu_{j,k}^N.
    \end{align*}
    If block $j$ has the single copy $k_j$, then
    $\beta'_j=\lvert\mu_{j,k_j}\rvert^N$, the coefficient of
    Theorem~\ref{thm:ldp_approximation_error_complex_weights}. Since the weights are
    nonzero, $\beta'\neq0$ whenever $\beta\neq0$.
\end{definition}

\begin{lemma}[The pairs on one copy]
    \label{lem:ldp_one_copy_state}
    \lean{Matrix.polarIso_mul_eq_zero, Matrix.tensorPower_smul,
        MPSTensor.conjTranspose_copyIsometry_mul_pairEmbedding,
        MPSTensor.conjTranspose_copyIsometry_mul_pairEmbedding_eq_zero,
        MPSTensor.polarIso_blockTensor_repeatedBlockSum_mul_pairEmbedding,
        MPSTensor.nonNormalApproxVector_repeatedBlockSum_oneCopy,
        MPSTensor.copyApproxOverlap_ofReal_mul,
        MPSTensor.nonNormalApproxOverlap_repeatedBlockSum_oneCopy}
    \leanok
    \uses{def:ldp_one_copy_weight, def:ldp_repeated_block_state,
        def:ldp_pair_family_fixed_point, thm:ldpolar_decomposition}
    Let $A$ be the direct sum with multiplicities of
    Definition~\ref{def:ldp_repeated_block_sum}, let $q\geq1$, and let $B=VP$ be the
    polar decomposition of its blocked map. The states of distinct blocks need not be
    orthogonal.
    \begin{enumerate}
        \item If $MY=0$ for a matrix $M$ with $M=WQ$ the polar decomposition of
              Theorem~\ref{thm:ldpolar_decomposition}, so that $W^\dagger W$ is the support
              projector of $Q$, then $WY=0$.
        \item $L_j^\dagger K_{j,k}=(\mu_{j,k}^q/c_j)\mathbf 1$, and
              $L_{j'}^\dagger K_{j,k}=0$ for $j'\neq j$.
        \item $VK_{j,k}=(\mu_{j,k}^q/c_j)\,VL_j$.
        \item For any matrices $\sigma_j$ and coefficients $\alpha_j$, with
              $\ket{\Omega_j}$ the product of the pairs of $\sigma_j$ over $M$ sites,
              \begin{align*}
                  V^{\otimes M}\sum_j\alpha_jK_{j,k_j}^{\otimes M}\ket{\Omega_j}
                  =V^{\otimes M}\sum_j\alpha_j\Bigl(\frac{\mu_{j,k_j}^q}{c_j}\Bigr)^M
                  L_j^{\otimes M}\ket{\Omega_j}.
              \end{align*}
        \item Multiplying all coefficients of the corrected state by one $t>0$ does not
              change its overlap with $\ket{\phi_N}$.
        \item The approximating state of Definition~\ref{def:ldp_non_normal_fixed_point}
              with the pairs of $\sigma_j$ on the copies $k_j$ and the coefficients
              $\beta'_j$ in place of $\beta_j$ has the same overlap with
              $\ket{\phi_N}$ as the corrected state of
              Definition~\ref{def:ldp_repeated_block_state} with the coefficients
              $\beta_j$.
    \end{enumerate}
\end{lemma}
\begin{proof}
    \leanok
    (1) From $MY=0$, $(QY)^\dagger QY=Y^\dagger M^\dagger MY=0$, so $QY=0$. The support
    projector of $Q$ is $\Pi=QR=R^\dagger Q$ for some $R$, so $\Pi Y=0$, and
    $(WY)^\dagger WY=Y^\dagger\Pi Y=0$.
    (2) $L_{j'}^\dagger K_{j,k}=\sum_{k'}(\mu_{j',k'}^q/c_{j'})K_{j',k'}^\dagger K_{j,k}$,
    and $K_{j',k'}^\dagger K_{j,k}$ is the identity for the same copy and zero otherwise.
    (3) By Theorem~\ref{thm:ldp_repeated_block_polar} and~(2),
    $BK_{j,k}=\sum_{j'}c_{j'}B_{j'}L_{j'}^\dagger K_{j,k}=\mu_{j,k}^qB_j$ and
    $BL_j=c_jB_j$, so $B(K_{j,k}-(\mu_{j,k}^q/c_j)L_j)=0$, and~(1) applies.
    (4) By~(3),
    $V^{\otimes M}K_{j,k_j}^{\otimes M}=(VK_{j,k_j})^{\otimes M}
        =(\mu_{j,k_j}^q/c_j)^M(VL_j)^{\otimes M}$.
    (5) The vector after $V^{\otimes M}$ is multiplied by $t$, and the overlap is
    normalized.
    (6) With $\alpha'_j=\beta'_j/(\sum_l\lvert\beta'_l\rvert^2)^{1/2}$, part~(4) turns the
    state into the corrected state with the coefficients
    $\alpha'_j(\mu_{j,k_j}^q/c_j)^M=\beta_j/(\sum_l\lvert\beta'_l\rvert^2)^{1/2}$, a positive
    multiple of $\beta_j$, and~(5) applies; if $\beta=0$ both states vanish.
\end{proof}

% **False source (eq. (S7), coefficients):** beta'_j in place of beta_j;
% Scope restriction (orthogonal blocks) for the first two declarations; Local fix
% of the rate for the last two;
% docs/paper-gaps/mswc24_repeated_block_corrected_state.tex.
\begin{theorem}[Approximation error with the pairs on one copy]
    \label{thm:ldp_approximation_error_one_copy}
    \lean{MPSTensor.exists_approximationError_le_repeatedBlockSum_oneCopy,
        MPSTensor.exists_approximationError_le_mul_repeatedBlockSum_oneCopy,
        MPSTensor.exists_approximationError_le_repeatedOverlappingBlockSum_oneCopy,
        MPSTensor.exists_approximationError_le_mul_repeatedOverlappingBlockSum_oneCopy}
    \leanok
    \uses{lem:ldp_one_copy_state, def:ldp_one_copy_weight,
        thm:ldp_approximation_error_repeated_orthogonal_blocks,
        thm:ldp_approximation_error_repeated_overlapping, def:correlation_length}
    Let $A$ be the direct sum with multiplicities of
    Definition~\ref{def:ldp_repeated_block_sum}, with normal blocks $A_j$ in the gauge
    (\ref{eq:ldp_normal_gauge}) and positive definite fixed points $\rho_j$ of trace one,
    let $0<\gamma<1$, and let $\widetilde\phi_N$ be the approximating state of
    Definition~\ref{def:ldp_non_normal_fixed_point} with the pairs of $\rho_j$ on the
    copies $k_j$ and the coefficients $\beta'_j$ of
    Definition~\ref{def:ldp_one_copy_weight} in place of $\beta_j$. Let
    $y=Me^{-\gamma q/\xi}$.
    \begin{enumerate}
        \item Let $\lambda_2$ bound the moduli of the eigenvalues other than one of every
              transfer map $\E_{A_j}$, with correlation length $\xi$. There are
              $C,C'>0$ such that for all nonzero complex weights, every choice of the
              copies $k_j$, every block length $q\geq1$ at which the $q$-site
              states of distinct blocks are orthogonal and every $M\geq1$ with
              $\beta\neq0$, $\varepsilon(\widetilde\phi_N,\phi_N)\leq Cy\,e^{Cy}$ and
              $\varepsilon(\widetilde\phi_N,\phi_N)\leq C'y$.
        \item Let $\lambda_2$ also bound the moduli of all eigenvalues of the mixed
              transfer maps $\E_{jj'}$, $j\neq j'$. There are $C,C'>0$ such that for all
              nonzero complex weights, every choice of the copies $k_j$, every $q\geq1$
              and every $M\geq1$ with $\beta\neq0$,
              $\varepsilon(\widetilde\phi_N,\phi_N)\leq Cy\,e^{Cy}$ and
              $\varepsilon(\widetilde\phi_N,\phi_N)\leq C'y$.
    \end{enumerate}
    Relative to \cite[Lemma~1$'$(ii)]{Malz2024Preparation} for the source's construction,
    this replaces the coefficients $\beta_j$ that
    \cite[eq.~(S7)]{Malz2024Preparation} takes from
    \cite[eq.~(S4)]{Malz2024Preparation} by $\beta'_j$, for weights of any phase, in the
    scopes of Theorems~\ref{thm:ldp_approximation_error_repeated_orthogonal_blocks}
    and~\ref{thm:ldp_approximation_error_repeated_overlapping}. As recorded in
    \cite{gap:mswc24_repeated_block_corrected_state}
    (Theorems~\ref{thm:ldp_approximation_error_repeated_block}
    and~\ref{thm:ldp_approximation_error_phase_weights}), with the coefficients
    $\beta_j$ the bound is false.
\end{theorem}
\begin{proof}
    \leanok
    By Lemma~\ref{lem:ldp_one_copy_state}(6) the overlap is that of the corrected state
    with the coefficients $\beta_j$, which
    Theorems~\ref{thm:ldp_approximation_error_repeated_orthogonal_blocks}
    and~\ref{thm:ldp_approximation_error_repeated_overlapping} bound.
\end{proof}

%% ==================================================
%% Subsection: Overlapping blocks at the second-order rate
%% Source: arXiv:2307.01696 (Malz, Styliaris, Wei, Cirac), Supplemental
%% Material, "Proof of Lemma 1 and extension to non-normal tensors",
%% eqs. (S2)-(S12) and Lemma 1'(ii).
%% Local fix: the rate of Lemma 1'(ii) is corrected by including the
%% correlation lengths of the mixed transfer maps; scope restriction to
%% blocks of multiplicity one and weight one;
%% docs/paper-gaps/mswc24_block_form_mixed_overlap.tex.
%% ==================================================

\subsection{Overlapping blocks at the second-order rate}
\label{sec:ldp_overlapping_second_order}

For blocks of multiplicity one and weight one whose states overlap,
Theorem~\ref{thm:ldp_approximation_error_overlapping_corrected} bounds the error at the
rate $e^{-\gamma q/\xi}$ with $\gamma<1$. The positive part $P$ converges to $P_\infty$
only to first order in the overlaps of distinct blocks, but the overlaps
$z_j=\braket{\phi_M(P_{\infty,\tilde\rho_j})}{\phi_M(P)}$ depend on $P$ only to second order
around $P_\infty$, as in the normal case
(Theorem~\ref{thm:ldp_positive_part_overlap_second_order}). This subsection proves the rate
$e^{-2\gamma q/\xi}$ for every $0<\gamma<1$, the rate of the normal case
(Theorem~\ref{thm:ldp_approximation_error_normal}). For the blocks of
Theorem~\ref{thm:ldp_approximation_error_overlapping_blocks} the error is at least
$(9/25)^q/16$ for $M\geq3$, with $3/5$ the eigenvalue of their mixed transfer map; since
$\lambda_2=3/5$ is admissible there, the exponent $2\gamma q/\xi$ cannot be replaced by $cq/\xi$
with $c>2$ \cite{gap:mswc24_block_form_mixed_overlap}.

\begin{theorem}[Powers of an element compressed by an idempotent]
    \label{thm:ldp_idempotent_compressed_second_order}
    \lean{IsIdempotentElem.exists_norm_trace_pow_sub_one_le_sq}
    \leanok
    \uses{thm:ldp_idempotent_second_order}
    Let $e$ be an idempotent of a normed algebra over $\C$, and let $\Tr$ be linear with
    $\Tr(xy)=\Tr(yx)$, $\lvert\Tr x\rvert\le K\lVert x\rVert$ and $\Tr e=1$. There is $C>0$
    such that every $T$ with $eTe=\alpha e$ for a real $\alpha\le1$,
    $\lVert T-e\rVert\le\delta$ and $1-\alpha\le\delta^2$ satisfies, for $M\ge2$ and
    $CM\delta^2<1$,
    \begin{align*}
        \bigl\lvert\Tr T^M-1\bigr\rvert\le C(M\delta^2)\,e^{CM\delta^2}.
    \end{align*}
\end{theorem}
\begin{proof}
    \leanok
    From $CM\delta^2<1$ with $C\ge4$ one gets $\delta\le1$ and $\alpha\ge3/4$. For
    $Z=\alpha^{-1}T-e$, one has $eZe=0$, $T=\alpha(e+Z)$, and
    $\lVert Z\rVert\le\alpha^{-1}(\lVert T-e\rVert+(1-\alpha)\lVert e\rVert)
        \le2(1+\lVert e\rVert)\delta$, so the blocks of $Z$ are bounded by a multiple $z$ of
    $\delta$ with $z\le1/2$. Theorem~\ref{thm:ldp_idempotent_second_order} bounds
    $\lvert\Tr(e+Z)^M-1\rvert$ by a multiple of $M\delta^2e^{cM\delta^2}$. Finally
    $\Tr T^M-1=\alpha^M(\Tr(e+Z)^M-1)+(\alpha^M-1)$, with $0\le\alpha^M\le1$ and
    $1-\alpha^M\le M(1-\alpha)\le M\delta^2$.
\end{proof}

\begin{lemma}[Overlaps of overlapping blocks, to second order]
    \label{lem:ldp_overlapping_overlaps_second_order}
    \lean{MPSTensor.trace_mixedMapLM_ofPhysicalMatrixLM_fixedPointTensor,
        MPSTensor.im_trace_mixedMapLM_ofPhysicalMatrixLM_fixedPointTensor,
        MPSTensor.cfc_sqrt_coordEmbedding_mul_mul, MPSTensor.trace_embeddedBlockState,
        MPSTensor.trace_coordEmbedding_mul_conjTranspose_mul_embeddedBlockState,
        MPSTensor.embeddedBlockState_mul_coordEmbedding_mul_conjTranspose,
        MPSTensor.blockSum_mul_coordEmbedding,
        MPSTensor.transferMap_blockSum_embeddedBlockState,
        MPSTensor.isIdempotentElem_of_apply_eq_trace_smul,
        MPSTensor.trace_of_apply_eq_trace_smul,
        MPSTensor.transferMatrix_mul_mul_of_eq_trace_smul,
        MPSTensor.trace_mul_mixedMapLM_fixedPointTensor,
        MPSTensor.isIdempotentElem_mixedTransferMatrixLeft_blockSumPosLimit,
        MPSTensor.mixedTransferMatrixLeft_blockSumPosLimit_mul_mul,
        MPSTensor.norm_sqrtWeightLM_fixedPointTensor, MPSTensor.norm_sqrtWeightLM_eq_one,
        MPSTensor.norm_eq_of_forall_abs_norm_sub_le_pow,
        MPSTensor.exists_norm_mpvOverlap_polarPosTensor_blockSum_sub_one_le_sq}
    \leanok
    \uses{thm:ldp_unit_vector_deficit, lem:ldp_overlapping_positive_part, lem:ldp_overlapping_overlaps,
        thm:ldp_idempotent_compressed_second_order,
        thm:ldp_positive_part_overlap_second_order, def:prep_fixed_point_tensor}
    \begin{enumerate}
        \item For a $D^2\times D^2$ matrix $H$, read as a tensor $X$ with physical index the
              output pair, and a $D\times D$ matrix $\sigma$, the mixed transfer map of $X$
              against the fixed-point tensor $P_{\infty,\sigma}$ satisfies
              $\Tr\sum_iX^i\sigma(P_{\infty,\sigma}^i)^\dagger
                  =\Tr\bigl(H(\sqrt\sigma^{\,\mathsf T}\otimes\sigma)\bigr)$, which is real
              when $H$ and $\sigma$ are Hermitian.
        \item For an isometry $E$ placing a block, $\sqrt{E\sigma E^\dagger}=E\sqrt\sigma E^\dagger$
              for $\sigma\ge0$, $\Tr(E\sigma E^\dagger)=\Tr(\Pi E\sigma E^\dagger)=\Tr\sigma$, and
              $EXE^\dagger\,\Pi=EXE^\dagger$ for $\Pi=EE^\dagger$.
        \item For blocks placed as in Lemma~\ref{lem:ldp_overlapping_positive_part} with
              weights one, $A^iE_j=E_jA_j^i$, and the transfer map of $A$ satisfies
              $\E_A(E_jXE_j^\dagger)=E_j\E_{A_j}(X)E_j^\dagger$; in particular
              $\tilde\rho_j=E_j\rho_jE_j^\dagger$ is a fixed point of $\E_A$.
        \item A linear map $R:X\mapsto\Tr(QX)\sigma$ has trace $\Tr(Q\sigma)$ and is
              idempotent when $\Tr(Q\sigma)=1$, and $RTR=\Tr(QT(\sigma))R$ for every linear
              map $T$. If $\sqrt\sigma Q=\sqrt\sigma$, then
              $\Tr(QY)=\Tr Y$ for every value $Y$ of a mixed transfer map against
              $P_{\infty,\sigma}$.
        \item The mixed transfer matrix $R_j$ of $P_\infty$ against $P_{\infty,\tilde\rho_j}$ is
              idempotent with trace one, and for every $D^2\times D^2$ matrix $H$, read as a
              tensor $X$, the mixed transfer matrix $\tau$ of $X$ against
              $P_{\infty,\tilde\rho_j}$ satisfies $R_j\tau R_j=\alpha R_j$ with
              $\alpha=\Tr\sum_iX^i\tilde\rho_j(P_{\infty,\tilde\rho_j}^i)^\dagger$.
        \item For a state $\sigma$, the vector $(P_{\infty,\sigma}^i\sqrt\sigma)_i$ has norm one,
              and so has $(X^i\sqrt\sigma)_i$ whenever $\Tr\E_X(\sigma)=1$. If
              $\lvert\lVert u_q\rVert-c\rvert\le K'x^q$ and
              $\lVert u_q-v\rVert\le Kx^q$ for every $q\ge1$, with $0\le x<1$, then
              $\lVert v\rVert=c$.
        \item In the setting of Lemma~\ref{lem:ldp_overlapping_positive_part} with weights
              one, if $\lambda_2\neq0$, then for every block $j$ there is $C_j>0$ such that for every $q\ge1$ and
              every $M\ge2$ with $C_jy<1$, where $y=Me^{-2\gamma q/\xi}$,
              \begin{align*}
                  \bigl\lvert\braket{\phi_M(P_{\infty,\tilde\rho_j})}{\phi_M(P)}-1\bigr\rvert
                  \le C_jy\,e^{C_jy}.
              \end{align*}
    \end{enumerate}
\end{lemma}
\begin{proof}
    \leanok
    Part~(1) is the entrywise computation: the letter of $P_{\infty,\sigma}$ at the pair
    $(l,r)$ is $\sqrt\sigma\ket l\bra r$, and the conjugate of $\Tr(HW)$ is
    $\Tr\bigl((HW)^\dagger\bigr)=\Tr(WH)$ for Hermitian $H$ and $W$. Part~(2) follows from
    $E^\dagger E=\Id$ and the uniqueness of the positive square root, and part~(3) from
    $E_k^\dagger E_j=\delta_{kj}\Id$. In part~(4),
    $R(T(R(X)))=\Tr(QX)\Tr(QT(\sigma))\sigma$, and the values of a mixed transfer map
    against $P_{\infty,\sigma}$ are sums of terms $Y\sqrt\sigma$, on which
    $\Tr(QY\sqrt\sigma)=\Tr(Y\sqrt\sigma Q)$. Part~(5) is part~(4) with the map
    $X\mapsto\Tr(\Pi_jX)\tilde\rho_j$ of Lemma~\ref{lem:ldp_overlapping_overlaps}(3) and
    $\Tr(\Pi_j\tilde\rho_j)=1$, since $\sqrt{\tilde\rho_j}\Pi_j=\sqrt{\tilde\rho_j}$ by
    part~(2). In part~(6), the squared norm of $(X^i\sqrt\sigma)_i$ is $\Tr\E_X(\sigma)$, and
    $\lvert\lVert v\rVert-c\rvert\le(K+K')x^q$ for every $q$.

    For part~(7), $\lambda_2\neq0$ gives $0<e^{-\gamma/\xi}<1$. Let $\iota(X)=(X^i\sqrt{\tilde\rho_j})_i$. By part~(3) and
    $\E_P=\E_A^q$, $\Tr\E_P(\tilde\rho_j)=1$, so $\iota(P)$ is a unit vector, as is
    $\iota(P_{\infty,\tilde\rho_j})$. The inner product
    $\braket{\iota(P_{\infty,\tilde\rho_j})}{\iota(P_\infty)}$ is the trace of
    $\Tr(\Pi_j\tilde\rho_j)\tilde\rho_j$, which is one, and $\iota(P_\infty)$ is a unit vector
    by part~(6) with $u_q=\iota(P)$; hence $\iota(P_\infty)=\iota(P_{\infty,\tilde\rho_j})$ and
    $\lVert\iota(P)-\iota(P_{\infty,\tilde\rho_j})\rVert\le K''e^{-\gamma q/\xi}$ by
    Lemma~\ref{lem:ldp_overlapping_positive_part}. The coefficient $\alpha$ of part~(5) is
    $\braket{\iota(P_{\infty,\tilde\rho_j})}{\iota(P)}$, real by part~(1) since $P$ and
    $\tilde\rho_j$ are Hermitian, at most one, and
    $1-\alpha=\tfrac12\lVert\iota(P)-\iota(P_{\infty,\tilde\rho_j})\rVert^2$. With
    $\lVert\tau-R_j\rVert\le K_3e^{-\gamma q/\xi}$ and $z_j=\Tr\tau^M$,
    Theorem~\ref{thm:ldp_idempotent_compressed_second_order} with
    $\delta$ a multiple of $e^{-\gamma q/\xi}$ gives the bound.
\end{proof}

\begin{lemma}[The norm of the target and one block of sites]
    \label{lem:ldp_overlapping_one_block_second_order}
    \lean{MPSTensor.mpvOverlap_fixedPointTensor_embeddedBlockState,
        MPSTensor.norm_sq_sub_norm_inner_sq_div_le,
        MPSTensor.exists_norm_mpvState_sq_sub_norm_sum_mpvOverlap_sq_le}
    \leanok
    \uses{lem:ldp_overlapping_overlaps, lem:ldp_overlapping_overlaps_second_order,
        lem:ldp_mixed_transfer_decay, thm:ldp_approximation_error_normal}
    In the setting of Lemma~\ref{lem:ldp_overlapping_positive_part} with weights one, let
    $\phi_N$ be the periodic state of $A$.
    \begin{enumerate}
        \item There is $K$ with $\bigl\lvert\lVert\phi_N\rVert^2-b\bigr\rvert\le Ke^{-\gamma N/\xi}$
              for every $N\ge1$.
        \item The states $\phi_N(P_{\infty,\tilde\rho_j})$ are orthonormal for every $N\ge1$.
        \item For a vector $W$ with $\lVert W\rVert^2=c>0$ and every $u$,
              $\lVert u\rVert^2-\lvert\braket Wu\rvert^2/c\le\lVert u-W\rVert^2$.
        \item If $\lambda_2\neq0$, there is $K'$ such that for every $q\ge1$
              \begin{align*}
                  \lVert\phi_q\rVert^2-\frac1b\Bigl\lvert\sum_j
                  \braket{\phi_1(P_{\infty,\tilde\rho_j})}{\phi_1(P)}\Bigr\rvert^2
                  \le K'e^{-2\gamma q/\xi}.
              \end{align*}
    \end{enumerate}
\end{lemma}
\begin{proof}
    \leanok
    Part~(1) is Lemma~\ref{lem:ldp_overlapping_overlaps}(5). Part~(2)
    holds because the pairs of the $\tilde\rho_j$ are orthonormal and the overlap of products
    of pairs is the $N$-th power of the overlap of the pairs. Part~(3) is
    $c-2\operatorname{Re}z+\lvert z\rvert^2/c=\lvert c-z\rvert^2/c\ge0$ with $z=\braket Wu$.
    For part~(4), with $u=\phi_1(P)$ and $W=\sum_j\phi_1(P_{\infty,\tilde\rho_j})$,
    $\lVert u\rVert=\lVert\phi_q\rVert$, the sum is $\braket Wu$, and
    $\lVert W\rVert^2=b$ by part~(2), so part~(3) bounds the left side by
    $\lVert u-W\rVert^2$. The overlaps of $\phi_1(P_\infty)$ with the
    $\phi_1(P_{\infty,\tilde\rho_j})$ are the traces of the idempotents $R_j$ of
    Lemma~\ref{lem:ldp_overlapping_overlaps_second_order}(5), which are one, and
    $\lVert\phi_1(P_\infty)\rVert^2=b$ by
    Lemma~\ref{lem:ldp_overlapping_overlaps_second_order}(6) and part~(1). Hence
    $\phi_1(P_\infty)=W$, and $\lVert u-W\rVert$ is at most a multiple of
    $\lVert P-P_\infty\rVert\le K''e^{-\gamma q/\xi}$.
\end{proof}

% **Local fix (rate of the overlapping blocks):** the rate of Lemma 1'(ii) of
% arXiv:2307.01696 is corrected by including the correlation lengths of the
% mixed transfer maps; scope restriction to multiplicity one and weight one;
% docs/paper-gaps/mswc24_block_form_mixed_overlap.tex.
\begin{theorem}[Approximation error for overlapping blocks at the second-order rate]
    \label{thm:ldp_approximation_error_overlapping_second_order}
    \lean{MPSTensor.exists_approximationError_le_overlappingBlockSum_sq,
        MPSTensor.exists_approximationError_le_mul_overlappingBlockSum_sq,
        MPSTensor.one_sub_norm_nonNormalApproxOverlap_blockSum_le,
        MPSTensor.one_sub_div_le_two_mul_of_sq_sub_sq_le}
    \leanok
    \uses{def:ldp_non_normal_fixed_point, lem:ldp_overlapping_overlaps,
        lem:ldp_overlapping_overlaps_second_order,
        lem:ldp_overlapping_one_block_second_order, def:correlation_length}
    In the setting of Theorem~\ref{thm:ldp_approximation_error_overlapping_corrected}, let
    $0<\gamma<1$. There are constants $C,C'>0$ such that for every
    block length $q$ and every $M\ge1$, with $N=qM$ and $y=Me^{-2\gamma q/\xi}$, which is
    $(N/q)e^{-2\gamma q/\xi}$ for $q\ge1$,
    \begin{align*}
        \varepsilon(\widetilde\phi_N,\phi_N)\le C'y
        \qquad\text{and}\qquad
        \varepsilon(\widetilde\phi_N,\phi_N)\le Cy\,e^{Cy}.
    \end{align*}
    Relative to \cite[Lemma~1$'$(ii)]{Malz2024Preparation} and
    \cite{gap:mswc24_block_form_mixed_overlap}, this establishes the bound for
    blocks of multiplicity one and weight one with $\xi_{\mathrm{diag}}$ replaced by
    $\max(\xi_{\mathrm{diag}},\xi_{\mathrm{off\text{-}diag}})$ and the exponent $\gamma q/\xi$
    sharpened to $2\gamma q/\xi$.
\end{theorem}
\begin{proof}
    \leanok
    If $e^{-\gamma/\xi}\ge1$, if $q=0$, or if $Cy\ge1$, the bound holds because
    $\varepsilon\le1$; otherwise $0<\lvert\lambda_2\rvert<1$ and $q\ge1$. With
    $z_j=\braket{\phi_M(P_{\infty,\tilde\rho_j})}{\phi_M(P)}$,
    $s=\lvert\sum_jz_j\rvert/\sqrt b$ and $t=\lVert\phi_N(A)\rVert$, by
    Lemma~\ref{lem:ldp_overlapping_overlaps}(2) the unnormalized approximating vector has
    inner product $\sum_jz_j/\sqrt b$ with $\ket{\phi_N(A)}$ and norm at most one, so
    $\varepsilon\le1-s/t$. If $t^2\ge1/2$ and $t^2-s^2\le X$ with $X\ge0$, then
    $1-s/t\le2X$: for $s<t$, $1-s/t\le1-s^2/t^2=(t^2-s^2)/t^2$. For $M=1$,
    Lemma~\ref{lem:ldp_overlapping_one_block_second_order}(4) gives $X=K'y$, and
    $t^2\ge b-Ke^{-\gamma q/\xi}\ge1/2$ once $4K^2y<1$. For $M\ge2$,
    $N\ge2q$, so $\lvert t^2-b\rvert\le Ky$ by
    Lemma~\ref{lem:ldp_overlapping_one_block_second_order}(1), and
    $\lvert\sum_jz_j-b\rvert\le\delta=\sum_jC_jy\,e^{Cy}$ by
    Lemma~\ref{lem:ldp_overlapping_overlaps_second_order}(7), so $s^2\ge b-2\delta$ and
    $X=Ky+2\delta$. The first bound follows from the second because $\varepsilon\le1$.
\end{proof}

\subsection{Canonical block pairs}

\begin{definition}[Normalized faithful block fixed points]
    \label{def:ldp_canonical_block_fixed_points}
    \lean{MPSTensor.IsBNTCanonicalForm.basisFixedPoint,
        MPSTensor.IsBNTCanonicalForm.basisFixedPoint_spec,
        MPSTensor.IsBNTCanonicalForm.basisFixedPointPair}
    \leanok
    \uses{def:ldp_sector_fixed_point_state}
    Let $A^i=\bigoplus_j\operatorname{diag}(\mu_{j,1},\ldots,\mu_{j,m_j})
        \otimes A_j^i$ be in basis-of-normal-tensors canonical form, with each
    $A_j$ left-canonical. Let $\sigma_j$ be the unique density matrix satisfying
    $\E_{A_j}(\sigma_j)=\sigma_j$. Irreducibility implies $\sigma_j>0$.
    Following \cite[eqs.~(8) and~(S7)]{Malz2024Preparation}, choosing one retained copy $k_j$ of each block and placing
    $\ket{\omega_j}=\operatorname{vec}(\sqrt{\sigma_j})$ on both virtual
    coordinates of that copy defines the canonical block pairs.

    With $E_{j,k_j}:\C^{D_j}\to\C^D$ denoting the corresponding coordinate inclusion:
    % Formula: omega_j = (E_{j,k_j} tensor E_{j,k_j}) vec(sqrt(sigma_j)).
    % Source: MSWC24 (8), (S2), (S7), with the specified retained copy.
    % Ink: the two E boxes embed the two block indices r,l into full bond indices R,L.
    % Contractions: the two block indices of sqrt(sigma_j) with the inputs of E.
    % Boundary signature: two north virtual legs R,L on each panel; no input legs.
    \begin{center}
        \tenkzkernel
        \begin{tenkzequation}
            \begin{tenkz}[rows={wire}, cols=3, bonds=none]
                \tn[at=(1,1), skin=box, wide=3,
                    ports={90@1:virtual:$R$,90@3:virtual:$L$}]{\omega_j}
            \end{tenkz}
            $=$
            \begin{tenkz}[rows={wire,wire}, cols=3, bonds=none]
                \tn[at=(1,1), skin=box, name=e1,
                    ports={90:virtual:$R$,270:virtual}]{E_{j,k_j}}
                \tn[at=(1,3), skin=box, name=e2,
                    ports={90:virtual:$L$,270:virtual}]{E_{j,k_j}}
                \tn[at=(2,1), skin=box, wide=3, name=sigma,
                    ports={90@1:virtual,90@3:virtual}]{\sqrt{\sigma_j}}
                \tnwire{sigma.90@1}{e1.270}
                \tnwire{sigma.90@3}{e2.270}
            \end{tenkz}
        \end{tenkzequation}
    \end{center}
\end{definition}

\begin{theorem}[Orthonormal canonical block pairs]
    \label{thm:ldp_canonical_block_pair_isometry}
    \lean{MPSTensor.IsBNTCanonicalForm.inner_basisFixedPointPair,
        MPSTensor.IsBNTCanonicalForm.isIsometry_pairIsometry_basisFixedPointPair}
    \leanok
    \uses{def:ldp_canonical_block_fixed_points, thm:ldp_local_orthogonality,
        thm:ldp_ghz_form}
    The block pairs satisfy $\langle\omega_j|\omega_{j'}\rangle
        =\delta_{jj'}$. Consequently $W\ket{j}=\ket{\omega_j}$ is an isometry.
\end{theorem}
\begin{proof}
    \leanok
    Each diagonal inner product is $\operatorname{tr}\sigma_j=1$, while
    distinct blocks occupy disjoint virtual coordinates.
\end{proof}

\begin{theorem}[Canonical copy coordinates and weights]
    \label{thm:ldp_canonical_copy_coordinates}
    \lean{MPSTensor.SectorDecomposition.copyWeights,
        MPSTensor.SectorDecomposition.bntWeight_copyWeights,
        MPSTensor.SectorDecomposition.copyCoord_disjoint,
        MPSTensor.SectorDecomposition.toTensor_eq_repeatedBlockSum}
    \leanok
    \uses{def:ldp_sector_fixed_point_state, def:ldp_repeated_block_sum}
    Let $\iota_{j,k}$ include the virtual coordinates of copy $k$ of block $j$
    in the flattened canonical direct sum. Their ranges are disjoint for distinct
    pairs $(j,k)$, and the canonical tensor satisfies
    \begin{align}
        A^i & =\sum_j\sum_k\mu_{j,k}\iota_{j,k}A_j^i\iota_{j,k}^{\dagger}.
    \end{align}
    Its sector coefficient at length $N$ is $\beta_j=\sum_k\mu_{j,k}^N$.

    % Formula: A^i = sum_{j,k} mu_{j,k} iota_{j,k} A_j^i iota_{j,k}^dagger.
    % Source: MSWC24 (S2), with the flattened canonical copy coordinates.
    % Ink: iota embeds block row r into alpha; iota^dagger restricts beta to block column s.
    % Contractions: the two horizontal block indices r,s on the right panel.
    % Boundary signature: west/east virtual alpha,beta and north physical i on both panels.
    \begin{center}
        \tenkzkernel
        \begin{tenkzequation}
            \begin{tenkz}[rows={wire}, cols=1, bonds=none]
                \tn[skin=box, ports={180:virtual:$\alpha$,0:virtual:$\beta$,
                            90:physical:$i$}]{A}
            \end{tenkz}
            $=\sum_{j,k}\mu_{j,k}$
            \begin{tenkz}[rows={wire}, cols=3]
                \tn[at=(1,1), skin=box, name=e,
                    ports={180:virtual:$\alpha$,0:virtual}]{\iota_{j,k}}
                \tn[at=(1,2), skin=box, name=a,
                    ports={180:virtual,0:virtual,90:physical:$i$}]{A_j}
                \tn[at=(1,3), skin=box, name=ed,
                    ports={180:virtual,0:virtual:$\beta$}]{\iota_{j,k}^{\dagger}}
            \end{tenkz}
        \end{tenkzequation}
    \end{center}
\end{theorem}
\begin{proof}
    \leanok
    Expanding the block-diagonal matrix as the sum of its inclusions and reindexing
    the flattened list by the pairs $(j,k)$ yields both identities.
\end{proof}

\begin{definition}[Corrected canonical one-copy state]
    \label{def:ldp_corrected_canonical_one_copy}
    \lean{MPSTensor.IsBNTCanonicalForm.oneCopyFixedPointState}
    \leanok
    \uses{def:ldp_canonical_block_fixed_points, def:ldp_one_copy_weight}
    At $N=qM$, put $\beta'_j=\beta_j(c_j/\mu_{j,k_j}^q)^M$ and
    $\alpha'_j=\beta'_j/(\sum_l|\beta'_l|^2)^{1/2}$, with
    $c_j=(\sum_k|\mu_{j,k}|^{2q})^{1/2}$. If every $\beta'_j$ vanishes, set
    $\alpha'=0$. The corrected state is
    \begin{align}
        \ket{\Omega'_{\mathrm{corr}}}
         & =\sum_j\alpha'_j\bigotimes_{s=1}^M\ket{\omega_j}_{R_sL_{s+1}}.
    \end{align}
    As recorded in \cite{gap:mswc24_repeated_block_corrected_state,gap:mswc24_block_form_mixed_overlap}, the coefficient correction is necessary even for multiplicity one, where
    $\beta'_j=|\mu_j|^{qM}$.
\end{definition}

\begin{theorem}[Canonical GHZ form and normalization]
    \label{thm:ldp_corrected_canonical_ghz}
    \lean{MPSTensor.IsBNTCanonicalForm.oneCopyFixedPointState_eq_tensorPower_mulVec_ghzState,
        MPSTensor.IsBNTCanonicalForm.oneCopyFixedPointState_norm_sq}
    \leanok
    \uses{def:ldp_corrected_canonical_one_copy, thm:ldp_canonical_block_pair_isometry,
        thm:ldp_ghz_form}
    In bond coordinates the corrected state equals
    $W^{\otimes M}\sum_j\alpha'_j\ket{j}^{\otimes M}$.
    If $M\geq1$ and $(\beta_j)_j\neq0$, its norm is one.

    % Formula: Omega'_corr = W tensor W tensor W |chi_3(alpha')>, in bond coordinates.
    % Source: MSWC24 (19), with corrected one-copy coefficients beta'_j.
    % Ink: chi_3 = sum_j alpha'_j |j,j,j>; each W sends j to omega_j.
    % Contractions: the three block-label indices between chi_3 and W.
    % Boundary signature: six north virtual legs R1,L2,R2,L3,R3,L1 on each panel.
    % The displayed M=3 instance has no omitted sites or undeclared trace closures.
    For three bonds the identity is
    \begin{center}
        \tenkzkernel
        \begin{tenkzequation}
            \begin{tenkz}[rows={wire}, cols=6, bonds=none]
                \tn[at=(1,1), skin=box, wide=6,
                    ports={90@1:virtual:$R_1$,90@2:virtual:$L_2$,
                            90@3:virtual:$R_2$,90@4:virtual:$L_3$,
                            90@5:virtual:$R_3$,90@6:virtual:$L_1$}]{\Omega'_{\mathrm{corr}}}
            \end{tenkz}
            $=$
            \begin{tenkz}[rows={wire,wire}, cols=6, bonds=none]
                \tn[at=(1,1), skin=box, wide=2, name=w1,
                    ports={90@1:virtual:$R_1$,90@2:virtual:$L_2$,270:virtual}]{W}
                \tn[at=(1,3), skin=box, wide=2, name=w2,
                    ports={90@1:virtual:$R_2$,90@2:virtual:$L_3$,270:virtual}]{W}
                \tn[at=(1,5), skin=box, wide=2, name=w3,
                    ports={90@1:virtual:$R_3$,90@2:virtual:$L_1$,270:virtual}]{W}
                \tn[at=(2,1), skin=box, wide=6, name=chi,
                    ports={90@1:virtual,90@3:virtual,90@5:virtual}]{\chi_3(\alpha')}
                \tnwire{chi.90@1}{w1.270}
                \tnwire{chi.90@3}{w2.270}
                \tnwire{chi.90@5}{w3.270}
            \end{tenkz}
        \end{tenkzequation}
    \end{center}
\end{theorem}
\begin{proof}
    \leanok
    Expanding the tensor power on the GHZ state gives the bond-coordinate formula, and orthonormality of the pairs gives
    squared norm $\sum_j|\alpha'_j|^2=1$; the nonzero copy weights ensure that
    $(\beta'_j)_j$ vanishes exactly when $(\beta_j)_j$ does.
\end{proof}

\begin{theorem}[Physical identification for the canonical tensor]
    \label{thm:ldp_corrected_canonical_physical}
    \lean{MPSTensor.IsBNTCanonicalForm.nonNormalApproxVector_oneCopyFixedPoint,
        MPSTensor.IsBNTCanonicalForm.nonNormalApproxOverlap_oneCopyFixedPoint}
    \leanok
    \uses{def:ldp_corrected_canonical_one_copy, thm:ldp_canonical_copy_coordinates,
        lem:ldp_one_copy_state}
    Let $q\geq1$, let $V$ be the polar factor of the $q$-site blocked
    canonical tensor, and let $L_j$ be its copy isometries. With
    $\ket{\Omega_j}$ denoting the product of the unembedded pairs of $\sigma_j$ and
    $K_{j,k_j}$ the pair embedding on the chosen copy,
    \begin{align}
        V^{\otimes M}\sum_j\beta'_jK_{j,k_j}^{\otimes M}\ket{\Omega_j}
         & =V^{\otimes M}\sum_j\beta_jL_j^{\otimes M}\ket{\Omega_j}.
    \end{align}
    Thus the unnormalized physical vector is independent of the chosen copies.
    Let $\ket{\phi_N}$ be the normalized periodic state of $A$ on $N=qM$ sites,
    put $\alpha_j=\beta_j/(\sum_l|\beta_l|^2)^{1/2}$, and let
    $\ket{\widetilde\phi'_N}$ and $\ket{\widetilde\phi_N}$ be the normalizations of
    $V^{\otimes M}\sum_j\alpha'_jK_{j,k_j}^{\otimes M}\ket{\Omega_j}$ and
    $V^{\otimes M}\sum_j\alpha_jL_j^{\otimes M}\ket{\Omega_j}$. Then
    \begin{align}
        \langle\widetilde\phi'_N|\phi_N\rangle
         & =\langle\widetilde\phi_N|\phi_N\rangle.
    \end{align}
    A vanishing coefficient vector or physical vector is normalized to zero.

    % Formula: V K_{j,k_j} = (mu_{j,k_j}^q / c_j) V L_j.
    % Source: corrected MSWC24 (S7), one-copy polar restriction identity.
    % Ink: V has the physical blocked output i and two full virtual inputs.
    % K and L map the two block virtual indices r,l into the full pair space.
    % Contractions: the two full virtual indices between V and K or L.
    % Boundary signature: two south virtual legs r,l and one north physical leg i.
    % The scalar has no legs; both panels keep exactly the same typed boundary.
    The local identity used in the proof is
    \begin{center}
        \tenkzkernel
        \begin{tenkzequation}
            \begin{tenkz}[rows={wire,wire}, cols=2, bonds=none]
                \tn[at=(1,1), skin=box, wide=2, name=v,
                    ports={90:physical:$i$,270@1:virtual,270@2:virtual}]{V}
                \tn[at=(2,1), skin=box, wide=2, name=k,
                    ports={90@1:virtual,90@2:virtual,
                            270@1:virtual:$r$,270@2:virtual:$l$}]{K_{j,k_j}}
                \tnwire{k.90@1}{v.270@1}
                \tnwire{k.90@2}{v.270@2}
            \end{tenkz}
            $=\frac{\mu_{j,k_j}^q}{c_j}$
            \begin{tenkz}[rows={wire,wire}, cols=2, bonds=none]
                \tn[at=(1,1), skin=box, wide=2, name=v,
                    ports={90:physical:$i$,270@1:virtual,270@2:virtual}]{V}
                \tn[at=(2,1), skin=box, wide=2, name=l,
                    ports={90@1:virtual,90@2:virtual,
                            270@1:virtual:$r$,270@2:virtual:$l$}]{L_j}
                \tnwire{l.90@1}{v.270@1}
                \tnwire{l.90@2}{v.270@2}
            \end{tenkz}
        \end{tenkzequation}
    \end{center}
\end{theorem}
\begin{proof}
    \leanok
    The relations $VK_{j,k_j}=(\mu_{j,k_j}^q/c_j)VL_j$ and
    $\beta'_j(\mu_{j,k_j}^q/c_j)^M=\beta_j$ yield the unnormalized identity. The nonzero copy weights make
    $\beta'=0$ exactly when $\beta=0$, in which case both states vanish. Otherwise,
    with $\lVert\beta\rVert=(\sum_l|\beta_l|^2)^{1/2}$ and likewise for $\beta'$,
    \begin{align*}
        V^{\otimes M}\sum_j\alpha'_jK_{j,k_j}^{\otimes M}\ket{\Omega_j}
         & =\frac{\lVert\beta\rVert}{\lVert\beta'\rVert}
        V^{\otimes M}\sum_j\alpha_jL_j^{\otimes M}\ket{\Omega_j},
    \end{align*}
    and the positive factor cancels in the normalized overlap.
\end{proof}


%% ==================================================
%% Section: Preparation with measurements
%% Sources: arXiv:2307.01696 (Malz, Styliaris, Wei, Cirac), paragraphs
%% "Preparations using measurements" and "Long-range MPS using
%% measurements"; arXiv:2103.13367 (Piroli, Styliaris, Cirac), paragraph
%% "Quantum circuits and LOCC" and "State transformations with QC and LOCC",
%% Definition "Transformations under QC and LOCC", and Example 1.
%% ==================================================

\section{Preparation with measurements}\label{sec:ldp_measurements}

Measurements followed by outcome-dependent corrections shorten
state preparation \cite{Malz2024Preparation}. In the measurement model of \cite{Piroli2021LOCC} set out
in Section~\ref{sec:qc_measurements}, Theorem~\ref{thm:qc_ghz_measurement} prepares the states
$\ket{\chi_M}=\sum_{i\in\mathbb Z_b}\alpha_i\ket{i}^{\otimes M}$ of
Definition~\ref{def:ldp_pair_family_fixed_point} in a depth
independent of $M$; their
connected correlations do not decay with distance, whereas by
Lemma~\ref{lem:qc_light_cone_correlations} those of a state prepared by a
local circuit of depth $T$ vanish beyond distance $2T$. This constant-depth preparation of long-range correlations powers both the measurement-assisted tree circuit and the preparation of tensors that are not normal.

%% ==================================================
%% Section: The tree circuit with measurements
%% Sources: arXiv:2307.01696 (Malz, Styliaris, Wei, Cirac), eq. (16) and
%% paragraph "Tree-RG circuit with measurements"; arXiv:2103.13367 (Piroli,
%% Styliaris, Cirac), paragraphs "State transformations with QC and LOCC" and
%% "Phases of matter", and Example 2.
%% ==================================================

\section{The tree circuit with measurements}\label{sec:ldp_tree_measurements}

The isometries of the tree circuit act on registers that are far apart. They
``can be teleported at neighboring registers with a constant overhead''
\cite{Malz2024Preparation}, so that every layer of the tree takes constant
depth. Theorem~\ref{thm:qc_long_range_layer} applies a layer of two-site gates
joining the sites $a$ and $a+2L+1$, on pairwise disjoint stretches
$\{a,\dots,a+2L+1\}$, with measurements in depth $5$ for every $L$. When every register is one site of the ring, composing these layers implements a binary tree of two-site gates with $k$ levels
with measurements in depth $5k$. Because the registers of a general tree circuit carry $\mathbb C^{D^2}$, Section~\ref{sec:ldp_register_tree} extends this construction to multi-site registers and establishes the resulting depth $O(\log\log(N/\varepsilon))$ for
matrix product states.

Throughout, the sites of the ring are $\{0,\dots,N-1\}$ with addition modulo
$N$, and the local labels are $\mathbb Z_d=\{0,\dots,d-1\}$ with $d\ge1$.

\begin{lemma}[Product vectors with zero at a set of sites]
    \label{lem:ldp_zero_sites_product_vector}
    \lean{MPSPreparation.isZeroOn_productVector}
    \leanok
    \uses{lem:qc_site_permutations}
    In the sense of Lemma~\ref{lem:qc_site_permutations}, a product vector whose site vectors at the sites of a set $Z$ are multiples
    of $\ket0$ has $\ket0$ at the sites of $Z$.
\end{lemma}
\begin{proof}
    \leanok
    The product vector vanishes at a configuration $x$ as soon as one factor
    $\pi_i(x_i)$ with $i\in Z$ and $x_i\neq0$ vanishes.
\end{proof}

\begin{lemma}[Binary trees of gates in depth $5k$ with measurements]
    \label{thm:ldp_tree_measurements}
    \lean{MPSPreparation.blockInterior, MPSPreparation.treeInterior,
        MPSPreparation.treeInterior_mono,
        MPSPreparation.blockInterior_subset_treeInterior, MPSPreparation.treeGate,
        MPSPreparation.levelGates, MPSPreparation.treeGate_a_val,
        MPSPreparation.treeGate_far_val, MPSPreparation.pairwise_levelGates,
        MPSPreparation.levelOp, MPSPreparation.treeOp,
        MPSPreparation.treeRounds, MPSPreparation.sum_depth_treeRounds,
        MPSPreparation.isZeroOn_levelOp_mulVec,
        MPSPreparation.interior_subset_blockInterior,
        MPSPreparation.isRoundsImplementationOn_treeRounds,
        MPSPreparation.isPreparedWithMeasurementRoundsInDepth_treeOp}
    \leanok
    \uses{thm:qc_long_range_layer, lem:qc_round_implementation_compose,
        lem:ldp_zero_sites_product_vector}
    For $i\ge0$, level $i$ of the tree applies, for every block
    $\{p2^{i+1},\dots,(p+1)2^{i+1}-1\}$ with $(p+1)2^{i+1}\le N$, a two-site
    unitary $u_{i,p}$ to the two endpoints $p2^{i+1}$ and
    $p2^{i+1}+2^{i+1}-1$ of the block. Let $T_k$ be the product of the levels
    $0,\dots,k-1$, the level $k-1$ applied first. Let $I_B$ be the set of
    sites $x$ whose residue modulo $B$ lies in $\{1,\dots,B-2\}$ and whose
    block $\{qB,\dots,(q+1)B-1\}$, $q=\lfloor x/B\rfloor$, satisfies
    $(q+1)B\le N$, and let $J_k=I_2\cup I_4\cup\dots\cup I_{2^k}$. Then $2k$
    rounds of total depth $5k$ implement $T_k$ on the vectors with $\ket0$ at
    the sites of $J_k$. In particular, for a nonzero product vector $\ket\pi$
    whose site vectors on $J_k$ are multiples of $\ket0$, the vector
    $T_k\ket\pi$ is prepared with measurement rounds in depth $5k$.
\end{lemma}
\begin{proof}
    \leanok
    The gate of the block $p$ of level $i$ is a distant gate with $a=p2^{i+1}$
    and $L=2^i-1$. Its stretch is the block, so the stretches of a level are
    pairwise disjoint, and its interior lies in $I_{2^{i+1}}\subseteq J_{i+1}$.
    By Theorem~\ref{thm:qc_long_range_layer} two rounds of depth $5$ implement
    level $i$ on the vectors with $\ket0$ at $J_{i+1}$. The gates of level $i$
    act on sites $x$ with $x\equiv0$ or $x+1\equiv0$ modulo $2^{i+1}$, hence
    modulo $2^{j+1}$ for every $j<i$, so on no site of $J_i$; and
    $J_i\subseteq J_{i+1}$. Level $i$ therefore maps the vectors with $\ket0$ at
    $J_{i+1}$ to vectors with $\ket0$ at $J_i$, and since $J_0$ is empty,
    Lemma~\ref{lem:qc_round_implementation_compose} composes the levels by
    induction over $k$.
\end{proof}

\begin{remark}[The tree circuit of matrix product states]
    \label{rem:ldp_tree_measurements_mps}
    \uses{thm:ldp_tree_measurements, thm:ldp_loglog_depth,
        thm:ldp_loglog_depth_every_length}
    In the tree circuit (\ref{eq:ldp_tree_factorization}) a register above
    the lowest level carries $\mathbb C^{D^2}$, which is a group of sites once
    $D^2>d$. Lemma~\ref{thm:ldp_loglog_depth} takes registers of $s$
    consecutive sites, $s$ a length at which the blocked tensor is injective,
    and obtains the depth $O(\log\log(N/\varepsilon))$ of
    \cite{Malz2024Preparation} for normal tensors and the chain lengths
    divisible by $s2^{k+1}$ for some $k$ with
    $a\log(N/\varepsilon)+b\le s2^{k+1}\le2(a\log(N/\varepsilon)+b)$; as recorded in
    \cite{gap:mswc24_tree_measurement_scope},
    Lemma~\ref{thm:ldp_loglog_depth_every_length} extends this depth bound to every chain
    length using trees on leaves of unequal widths.
\end{remark}

%% ==================================================
%% Section: Registers of several sites and the depth O(log log(N/eps))
%% Sources: arXiv:2307.01696 (Malz, Styliaris, Wei, Cirac), eqs. (5),
%% (10)-(12) and (16), Lemma 1'(i), and paragraph "Tree-RG circuit with
%% measurements".
%% ==================================================

\section{Registers of several sites and the depth \texorpdfstring{$O(\log\log(N/\varepsilon))$}{O(log log(N/epsilon))}}
\label{sec:ldp_register_tree}

In the tree circuit (\ref{eq:ldp_tree_factorization}) a register above the
lowest level carries $\C^{D^2}$, which is a group of sites once $D^2>d$. This
section groups the sites into registers of $s$ consecutive sites.
Lemma~\ref{thm:qc_register_layer} teleports registers site by site and
applies a layer of unitaries on pairs of distant registers in a depth
depending only on $s$ and on the number of nearest-neighbour gates of each
unitary. With $s$ a length at which the blocked
tensor is injective, the trees of all the blocks prepare the approximating state
with measurements in depth $O(k)$ for block length $s2^{k+1}$, and Lemma~1'(i)
of \cite{Malz2024Preparation} gives the depth $O(\log\log(N/\varepsilon))$ of
the paragraph ``Tree-RG circuit with measurements'' of the source, for the chain
lengths divisible by $s2^{k+1}$ for some $k$ with
$a\log(N/\varepsilon)+b\le s2^{k+1}\le2(a\log(N/\varepsilon)+b)$.
Section~\ref{sec:ldp_unequal_trees} places trees on blocks of any length and
obtains this depth for normal tensors $A$ and every chain length $N\ge2$ with
$\ket{\phi_N(A)}\neq0$, which holds for all large $N$
\cite{gap:mswc24_tree_measurement_scope}.

Throughout, the sites of the ring are $\{0,\dots,N-1\}$ with addition modulo
$N$, the local labels are $\mathbb Z_d$ with $d\ge1$, and $s\ge1$.

\begin{lemma}[Trees of register gates on the blocks]
    \label{thm:ldp_register_tree_rounds}
    \lean{MPSPreparation.blockSite_add_natCast,
        MPSPreparation.blockLayerOp_mul, MPSPreparation.blockLayerOp_one,
        MPSPreparation.blockLayerOp_eq_list_prod, MPSPreparation.regSpan,
        MPSPreparation.regOffset, MPSPreparation.mul_sub_two_add_two_mul,
        MPSPreparation.block_pos, MPSPreparation.regWindow,
        MPSPreparation.regSpan_eq, MPSPreparation.two_le_regSpan,
        MPSPreparation.succ_mul_regSpan_le, MPSPreparation.regOffset_lt,
        MPSPreparation.mul_mul_regSpan_add_lt, MPSPreparation.regWindow_val,
        MPSPreparation.regWindow_injective,
        MPSPreparation.mul_mul_add_ne_of_ne,
        MPSPreparation.regWindow_val_ne_of_interior,
        MPSPreparation.regWindow_injective₂, MPSPreparation.regLevelOp,
        MPSPreparation.regTreeOp, MPSPreparation.block_le,
        MPSPreparation.regGate, MPSPreparation.regLevelGates,
        MPSPreparation.two_mul_regGate_L, MPSPreparation.regGate_a_add,
        MPSPreparation.regGate_sites, MPSPreparation.mem_span_regGate,
        MPSPreparation.mem_interior_regGate,
        MPSPreparation.blockSite_ne_of_subBlock_ne,
        MPSPreparation.mem_regLevelGates,
        MPSPreparation.pairwise_regLevelGates,
        MPSPreparation.list_prod_map_op_regLevelGates,
        MPSPreparation.regWindow_val_ne, MPSPreparation.regInteriorFrom,
        MPSPreparation.regInteriorFrom_succ_subset,
        MPSPreparation.interior_subset_regInteriorFrom,
        MPSPreparation.disjoint_range_sites_interior,
        MPSPreparation.isZeroOn_regInteriorFrom_mulVec,
        MPSPreparation.regCentralSites,
        MPSPreparation.regInteriorFrom_zero_subset,
        MPSPreparation.exists_rounds_regTreeOp_aux,
        MPSPreparation.exists_rounds_blockLayerOp_regTreeOp}
    \leanok
    \uses{thm:qc_register_layer, lem:qc_round_implementation_compose,
        lem:qc_placed_operators}
    Let $k\ge0$, $q=s2^{k+1}$, and cut the ring of $N=Mq$ sites into $M\ge1$
    blocks of $q$ sites, each grouped into $2^{k+1}$ registers of $s$
    consecutive sites. For $0\le j\le k$ and $0\le p<2^j$ the sub-block $p$ of
    depth $j$ consists of the registers $p2^{k+1-j},\dots,(p+1)2^{k+1-j}-1$, and
    a unitary $X_{j,p}$ on $2s$ sites acts on its first and its last register.
    Let $T$ be the product of the depths $k,\dots,0$ of a block, the depth $0$
    applied first. If every $X_{j,p}$ is a product of at most $K$ unitaries on
    neighbouring sites, measurement rounds of total depth $(k+1)(4s+K+2)$
    implement $T$ on every block, $T^{\otimes M}$, on the vectors with $\ket0$
    at every site of every block other than its first and its last $s$ sites.
\end{lemma}
\begin{proof}
    \leanok
    The unitaries of depth $j$ of all the blocks are register gates whose
    stretches are their sub-blocks, which are pairwise disjoint, so
    Lemma~\ref{thm:qc_register_layer} applies the depth $j$ in depth
    $4s+K+2$ on the vectors with $\ket0$ at the sites strictly inside the
    sub-blocks of depth $j$ between their first and last registers. A unitary
    of depth $j'<j$ acts on the first and the last register of a sub-block of
    depth $j'$, which are the first or the last register of a sub-block of
    depth $j$, so on no site where depth $j$ needs $\ket0$. Induction on the
    depth and Lemma~\ref{lem:qc_round_implementation_compose} give the claim.
\end{proof}

\begin{lemma}[Amplitudes of trees]
    \label{lem:ldp_tree_amplitudes}
    \lean{MPSTensor.unpair, MPSTensor.two_mul_add_lt_two_pow,
        MPSTensor.unpair_apply, MPSTensor.unpair_injective,
        MPSTensor.unpair_surjective, MPSTensor.unpair_bijective,
        MPSTensor.decodeBlock_blockIndexOfList, MPSTensor.decodeBlock_finCongr,
        MPSTensor.decodeBlock_pairRegroupEquiv, MPSTensor.treeAmp,
        MPSTensor.binaryTreeMatrix_apply_unpair}
    \leanok
    \uses{def:ldp_binary_mera}
    The tree map (\ref{eq:ldp_binary_mera_tree}) with finest layer $W$ and
    coarser layers $V_0,V_1,\dots$, at $2^{j+1}$ leaves grouped into
    neighbouring pairs $e_0,\dots,e_{2^j-1}$, is $\sum_{e'}\prod_pW(e_p,e'_p)$
    times the tree map of the coarser layers at the outputs $e'$ grouped into
    pairs in their turn, the sum running over these outputs.
\end{lemma}
\begin{proof}
    \leanok
    Unfolding the recursion that defines the tree map and writing out the sum
    of a matrix product gives the claimed expression.
\end{proof}

\begin{lemma}[The tree of registers implements the isometry of a block]
    \label{thm:ldp_register_tree_isometry}
    \lean{Matrix.polarIso_submatrix_equiv,
        MPSPreparation.exists_unitary_apply_eq_extend,
        MPSPreparation.extUnitary, MPSPreparation.extUnitary_mem_unitary,
        MPSPreparation.extUnitary_apply,
        MPSPreparation.regSpan_eq_two_mul, MPSPreparation.regWindow_two_mul,
        MPSPreparation.regWindow_two_mul_add_one,
        MPSPreparation.regWindow_two_mul_notMem,
        MPSPreparation.regWindow_two_mul_add_one_notMem,
        MPSPreparation.exists_regWindow_eq,
        MPSPreparation.polarIsoMatrix_blockTensor_blockTensor,
        MPSPreparation.iteratedBlockIndex_decodeBlockEquiv_symm,
        MPSPreparation.regInput, MPSPreparation.coarseOutput,
        MPSPreparation.fineOutput, MPSPreparation.regInput_injective,
        MPSPreparation.coarseOutput_injective,
        MPSPreparation.fineOutput_bijective, MPSPreparation.regLayer,
        MPSPreparation.regTreeGate, MPSPreparation.regTreeGate_mem_unitary,
        MPSPreparation.regTreeGate_apply_of_lt,
        MPSPreparation.regTreeGate_apply_of_le,
        MPSPreparation.regRegisterCfg, MPSPreparation.regLeafCfg,
        MPSPreparation.regRegisterCfg_injective,
        MPSPreparation.regLeafCfg_bijective, MPSPreparation.regLeafCfg_apply,
        MPSPreparation.regRegisterCfg_eq_placeCfg,
        MPSPreparation.placeCfg_regInput_zero,
        MPSPreparation.regTreeOp_apply_regRegisterCfg,
        MPSPreparation.regTreeOp_apply_regLeafCfg,
        MPSPreparation.regTreeOp_apply_blockInputCfg}
    \leanok
    \uses{thm:ldp_tree_factorization, lem:ldp_tree_layers_isometries,
        thm:ldp_unitary_implementing_isometry, lem:ldp_tree_amplitudes,
        lem:qc_placed_configurations, thm:ldp_register_tree_rounds,
        thm:ldpolar_tensor_decomposition}
    Let $A$ be a tensor whose blocked tensor $\mathcal B$ over $s$ sites has an
    injective two-site blocked tensor, let $k\ge0$ and $q=s2^{k+1}$, and let
    $\mathcal B_q=VP$ be the polar decomposition of $A$ blocked over $q$ sites.
    Let a register of $s$ sites carry the bond legs $\C^D$ through an injective
    map of $\{0,\dots,D-1\}$ into its configurations, and $\C^{D^2}$ through an
    injective map of $\{0,\dots,D^2-1\}$ into them. Every layer of the tree
    circuit (\ref{eq:ldp_tree_factorization}) of $\mathcal B$ is an isometry
    from one register to two, and extends to a unitary on the two registers of
    each sub-block of Lemma~\ref{thm:ldp_register_tree_rounds}. The tree $T$
    of these unitaries satisfies
    \begin{align*}
        \bra\tau T\ket{l,0,\dots,0,r}=\bra\tau V\ket{l,r}
    \end{align*}
    for every configuration $\tau$ of the block and all bond indices $l,r$,
    where $\ket{l,0,\dots,0,r}$ carries $l$ on the first and $r$ on the last
    register.
\end{lemma}
\begin{proof}
    \leanok
    By Lemma~\ref{lem:ldp_tree_layers_isometries} the layers of the tree circuit
    of $\mathcal B$ are isometries, and they extend to unitaries
    (Lemma~\ref{thm:ldp_unitary_implementing_isometry}) from their input
    register, the other register in $\ket0$, to their two output registers.
    The unitary of the root reads $l$ and $r$. After the depths $0,\dots,j$ the
    registers of the sub-blocks of depth $j+1$ carry the outputs of these
    layers, with the amplitude of the tree of the coarser layers, by
    Lemma~\ref{lem:qc_placed_configurations}: the registers written by
    a unitary of depth $j$ are the first and the last register of the two
    halves of its sub-block, on which the two unitaries below it act. At the
    finest depth the block carries the leaves, with the amplitude of the whole
    tree map, which by Lemma~\ref{lem:ldp_tree_amplitudes} and
    Lemma~\ref{thm:ldp_tree_factorization} is the isometry of the polar
    decomposition of $\mathcal B$ blocked over $2^{k+1}$ sites, that is, of
    $A$ blocked over $q$ sites.
\end{proof}

\begin{lemma}[The approximating state with measurements in depth $O(k)$]
    \label{thm:ldp_register_tree_approximating_state}
    \lean{MPSPreparation.isZeroOn_regCentralSites_pairLayerOp_mulVec,
        MPSPreparation.exists_isPreparedWithMeasurementRoundsInDepth_approximatingMPVState}
    \leanok
    \uses{thm:ldp_register_tree_isometry, thm:ldp_register_tree_rounds,
        thm:ldp_approximating_state_depth, def:ldp_approximating_state,
        thm:qc_bounded_unitary_two_site_gates, thm:ldp_chain_depth,
        def:qc_measurement_rounds}
    For every $s\ge1$ there is $C$, depending only on $d$ and $s$, with the
    following property. Let $A$ be a tensor whose blocked tensor over $s$ sites
    is injective, and let $\sigma\ge0$ with $\Tr\sigma=1$. For every $k\ge0$
    and $M\ge1$, the approximating state of
    Definition~\ref{def:ldp_approximating_state} with blocks of $q=s2^{k+1}$
    sites, on $N=Mq$ sites, is prepared with measurement rounds in depth at most
    $C(k+1)$. No outcome is post-selected.
\end{lemma}
\begin{proof}
    \leanok
    Injectivity gives $D^2\le d^s$, so a register of $s$ sites carries
    $\C^{D^2}$ and $\C^D$, and the two-site blocked tensor of the injective
    blocked tensor is injective. By
    Theorem~\ref{thm:qc_bounded_unitary_two_site_gates} there is $K$,
    depending on $d$ and $s$, such that every unitary on $2s$ sites is a
    product of at most $K$ unitaries on neighbouring sites. A local circuit of
    depth $K$, applied in one round of depth $K+2$, prepares the pairs
    $\ket\omega$ of the fixed point on the last register of every block and the
    first of the next, from the all-$\ket0$ product state; the other sites keep
    $\ket0$. Lemma~\ref{thm:ldp_register_tree_rounds}, with the unitaries of
    Lemma~\ref{thm:ldp_register_tree_isometry}, then applies the trees of all
    the blocks in depth $(k+1)(4s+K+2)$. The output is
    $T^{\otimes M}\bigotimes\ket\omega=V^{\otimes M}\bigotimes\ket\omega$, the
    approximating state, as in Lemma~\ref{thm:ldp_approximating_state_depth};
    the total depth is at most $(2K+4s+4)(k+1)$.
\end{proof}

\begin{lemma}[Preparation with measurements in depth $O(\log\log(N/\varepsilon))$]
    \label{thm:ldp_loglog_depth}
    \lean{MPSPreparation.exists_isPreparedWithMeasurementRoundsInDepth_approximationError_le,
        MPSPreparation.exists_isPreparedWithMeasurementRoundsInDepth_le_logb_log}
    \leanok
    \uses{thm:ldp_register_tree_approximating_state,
        lem:ldp_approximation_error_slope_threshold,
        thm:normalized_primitive_gauge, lem:ldp_normal_gap,
        lem:ldp_approximating_state_unit_norm,
        lem:ldp_normalized_state_rescaling}
    Let $A$ be normal with $D\ge1$. There are $s\ge1$, $C$, $a>0$ and $b\ge1$,
    depending only on $A$, with the following property. Let $\varepsilon>0$,
    $N\ge1$ and $k\ge0$ be such that $s2^{k+1}$ divides $N$ and
    $s2^{k+1}\ge a\log(N/\varepsilon)+b$. Then a unit vector $\ket\psi$ with
    $\varepsilon(\psi,\phi_N)\le\varepsilon$ is prepared with measurement
    rounds in depth at most $C(k+1)$. If
    $s2^{k+1}\le2(a\log(N/\varepsilon)+b)$, then
    $k\le\log_2(a\log(N/\varepsilon)+b)$, so the depth is at most
    $C(\log_2(a\log(N/\varepsilon)+b)+1)=O(\log\log(N/\varepsilon))$. For
    normal tensors, this matches the depth in the paragraph ``Tree-RG circuit with
    measurements'' of~\cite{Malz2024Preparation} at chain lengths divisible by
    $s2^{k+1}$ for some $k$ in this window, for instance $N=Ms2^K$ with $K$
    large; the remaining chain lengths are treated in
    Lemma~\ref{thm:ldp_loglog_depth_every_length}.
\end{lemma}
\begin{proof}
    \leanok
    Bringing $A$ into the gauge~(\ref{eq:ldp_normal_gauge}) as in the proof of
    Lemma~\ref{thm:ldp_depth_upper_bound_divisible} multiplies
    $\ket{\phi_N}$ by a nonzero scalar; let $\xi$ be the correlation length of a
    bound $t<1$ on the moduli of the subleading eigenvalues, and let $s$ be a
    length at which the blocked tensor in this gauge is injective. With $a=\xi$
    and $b=b_0+1$ for $b_0$ from
    Lemma~\ref{lem:ldp_approximation_error_slope_threshold}, and with $N=Mq$ and
    $q=s2^{k+1}$, the approximating state with blocks of $q$ sites is a unit
    vector with error at most $\varepsilon$, prepared in depth at most
    $C(k+1)$ by Lemma~\ref{thm:ldp_register_tree_approximating_state}. If
    $s2^{k+1}\le2(a\log(N/\varepsilon)+b)$ then $2^k\le
        a\log(N/\varepsilon)+b$ since $s\ge1$.
\end{proof}

%% ==================================================
%% Section: Trees on blocks of any length and the depth O(log log(N/eps))
%% for every chain length
%% Sources: arXiv:2307.01696 (Malz, Styliaris, Wei, Cirac), eqs. (5),
%% (10)-(12) and (16), Lemma 1 and Lemma 1'(i), paragraph "Tree-RG circuit
%% with measurements", and Supplemental Material, proof of Theorem 1.
%% ==================================================

\section{Trees on blocks of any length}
\label{sec:ldp_unequal_trees}

The trees of Section~\ref{sec:ldp_register_tree} cover blocks of $s2^{k+1}$
sites, so they prepare the approximating state only for chain lengths divisible
by $s2^{k+1}$. The Supplemental Material of~\cite{Malz2024Preparation}, proof
of Theorem~1, cuts a chain of any length into blocks ``all of the same size,
$q_N$, except for the last one, which may be larger''. This section places trees
on blocks of any length, with leaves of unequal widths, and obtains the depth
$O(\log\log(N/\varepsilon))$ of the paragraph ``Tree-RG circuit with
measurements'' of the source for normal tensors $A$ and every chain length
$N\ge2$ with $\ket{\phi_N(A)}\neq0$, which holds for all large $N$
\cite{gap:mswc24_tree_measurement_scope}.

Throughout, $d\ge1$, $s\ge1$, the sites of a block of $n$ sites are
$\{0,\dots,n-1\}$, and for a tensor $A$ we write $\mathcal B_n=V_nP_n$ for the
polar decomposition of $A$ blocked over $n$ sites, with $V_n$ read on the
configurations of the $n$ sites.

\begin{lemma}[The partial isometry of a block of unequal halves]
    \label{lem:ldp_polar_merge}
    \lean{MPSPreparation.pairProductTensor, MPSPreparation.cfgPolarIso,
        MPSPreparation.mergeIso, MPSPreparation.isIsometry_cfgPolarIso,
        MPSPreparation.blockTensor_add_apply,
        MPSPreparation.cfgPolarIso_append, MPSPreparation.cfgPolarIso_split,
        MPSPreparation.polarPosTensor_blockTensor_add,
        MPSPreparation.pairProductTensor_self, MPSPreparation.mergeIso_self,
        MPSPreparation.isInjective_pairProductTensor,
        MPSPreparation.isIsometry_mergeIso,
        MPSPreparation.polarPosTensor_blockTensor_two_pow,
        MPSPreparation.mergeIso_two_pow}
    \leanok
    \uses{thm:ldpolar_tensor_decomposition, thm:ldp_tree_factorization}
    Let $A$ be a tensor, $n_1,n_2\ge0$, and let
    $W:\C^{D^2}\to\C^{D^2}\otimes\C^{D^2}$ be the partial isometry of the polar
    decomposition of the tensor $(a,b)\mapsto P_{n_1}^aP_{n_2}^b$ of two
    neighbouring sites. Then
    \begin{align}
        V_{n_1+n_2}=(V_{n_1}\otimes V_{n_2})W,
        \label{eq:ldp_polar_merge}
    \end{align}
    and $P_{n_1+n_2}$ is the positive part of the tensor
    $(a,b)\mapsto P_{n_1}^aP_{n_2}^b$. No injectivity is assumed, as in the
    footnote of~\cite[paragraph ``The sequential-RG circuit'']{Malz2024Preparation}.
    For $n_1=n_2=n$ the tensor $(a,b)\mapsto P_n^aP_n^b$ is the two-site blocked
    tensor of $P_n$, and $W$ is the partial isometry of its polar decomposition.
    For $n=2^j$, $j\ge1$, the positive part $P_n$ is the tensor $T_j$ of the
    tree of \cite[eq.~(16)]{Malz2024Preparation}, so $W$ is its layer
    $V^{(j+1)}$. If the blocked tensor over $n_1+n_2$
    sites is injective, so is the tensor $(a,b)\mapsto P_{n_1}^aP_{n_2}^b$, and
    $W$ is an isometry.
\end{lemma}
\begin{proof}
    \leanok
    \uses{thm:ldpolar_tensor_rotate, thm:ldpolar_tensor_injective}
    Write $\mathcal B_{n_i}^{\tau}=\sum_y\bra\tau V_{n_i}\ket yP_{n_i}^y$. Then
    $\mathcal B_{n_1+n_2}^{\tau_1\tau_2}
        =\mathcal B_{n_1}^{\tau_1}\mathcal B_{n_2}^{\tau_2}$, so
    $\mathcal B_{n_1+n_2}$ is the tensor $C$ of the positive parts rotated by
    $K=V_{n_1}\otimes V_{n_2}$. Here
    $K^\dagger K=\Pi_{n_1}\otimes\Pi_{n_2}$, and $\Pi_{n_i}$ fixes the tensor of
    $P_{n_i}$, so $K^\dagger K$ fixes $C$. By
    Theorem~\ref{thm:ldpolar_tensor_rotate}, $V_{n_1+n_2}=KW$ and $C$ has the
    positive part of $\mathcal B_{n_1+n_2}$. The matrices
    $\mathcal B_{n_1+n_2}^{\tau}$ are combinations of the products
    $P_{n_1}^aP_{n_2}^b$, so these span the full matrix algebra when the former
    do. For $n=2^j$, the equal-halves case of the identity of positive parts
    gives $P_{2^{j+1}}$ as the positive part of the two-site blocked tensor of
    $P_{2^j}$, which is the recursion defining $T_{j+1}$ from $T_j$; induction
    on $j$ from $P_2=T_1$ gives $P_{2^j}=T_j$.
\end{proof}

\begin{theorem}[Layers of register gates one after the other]
    \label{thm:ldp_register_layers}
    \lean{MPSPreparation.registerLayersOp,
        MPSPreparation.exists_rounds_registerLayersOp}
    \leanok
    \uses{thm:qc_register_layer, lem:qc_round_implementation_compose}
    Let $G_0,\dots,G_{m-1}$ be lists of register gates, the gates of each list
    with pairwise disjoint stretches, such that no register of a gate of $G_{j'}$
    lies in the interior of a gate of $G_j$ for $j'<j$, and every matrix of a
    gate a product of at most $K$ unitaries on neighbouring sites. Then
    measurement rounds of total depth $m(4s+K+2)$ implement the product of the
    gates of $G_{m-1},\dots,G_0$, the list $G_0$ applied first, on the vectors
    with $\ket0$ at the interiors of all the gates.
\end{theorem}
\begin{proof}
    \leanok
    By Theorem~\ref{thm:qc_register_layer} each list is applied in depth
    $4s+K+2$ on the vectors with $\ket0$ at the interiors of its gates. The
    gates of $G_{j'}$ act on their registers only, which lie outside the
    interiors of the later lists, so they keep $\ket0$ there. Induction on the
    number of lists and Lemma~\ref{lem:qc_round_implementation_compose}.
\end{proof}

\begin{definition}[Tree layout of a block]
    \label{def:ldp_tree_layout}
    \lean{MPSPreparation.leafOffset, MPSPreparation.leafOffset_succ,
        MPSPreparation.leafOffset_mono, MPSPreparation.leafOffset_add_le,
        MPSPreparation.even_leafOffset_sub, MPSPreparation.IsTreeLayout,
        MPSPreparation.IsTreeLayout.pos, MPSPreparation.nodeLeaves,
        MPSPreparation.nodeStart, MPSPreparation.nodeStop,
        MPSPreparation.nodeLeaves_pos, MPSPreparation.nodeLeaves_eq_two_mul,
        MPSPreparation.succ_mul_nodeLeaves_le, MPSPreparation.nodeStart_add_le,
        MPSPreparation.nodeStop_le, MPSPreparation.nodeStop_le_n,
        MPSPreparation.even_nodeStop_sub, MPSPreparation.nodeStop_le_nodeStart,
        MPSPreparation.nodeStart_two_mul,
        MPSPreparation.nodeStop_two_mul_add_one, MPSPreparation.nodeStop_two_mul,
        MPSPreparation.nodeOffset, MPSPreparation.nodeOffset_lt,
        MPSPreparation.nodeStart_le_nodeOffset,
        MPSPreparation.exists_nodeOffset_eq,
        MPSPreparation.nodeOffset_notMem_interior, MPSPreparation.nodeWindow,
        MPSPreparation.nodeWindow_val, MPSPreparation.nodeWindow_injective₂,
        MPSPreparation.nodeLen, MPSPreparation.leafLen,
        MPSPreparation.leafLen_eq, MPSPreparation.leafOffset_add_leafLen,
        MPSPreparation.two_mul_le_leafLen, MPSPreparation.leafWindow,
        MPSPreparation.leafWindow_val, MPSPreparation.leafWindow_injective,
        MPSPreparation.disjoint_range_leafWindow,
        MPSPreparation.exists_leafWindow_eq, MPSPreparation.leafWindow_succ}
    \leanok
    A \emph{tree layout} of a block of $n$ sites with $h+1$ coarse depths is a
    sequence of $2^{h+1}$ even widths $w_0,w_1,\dots$, each at least $2s$, with
    $o_{2^{h+1}}\le n$, where $o_m=w_0+\cdots+w_{m-1}$. The leaf $i$ covers the
    sites $o_i,\dots,o_{i+1}-1$, the last leaf also the sites $o_{2^{h+1}},\dots,
        n-1$. For $0\le j\le h+1$ and $0\le p<2^j$, the node $p$ of depth $j$
    covers the leaves $p2^{h+1-j},\dots,(p+1)2^{h+1-j}-1$; its two registers are
    the sites $o_{p2^{h+1-j}},\dots,o_{p2^{h+1-j}}+s-1$ and the $s$ sites before
    $o_{(p+1)2^{h+1-j}}$. The registers of a node are an even number of sites
    apart, and they are the first register of its first half and the last
    register of its second half. A register of a node of depth $j'$ is a
    register of a node of every depth $j\ge j'$, and no register of a node of
    depth $j$ lies strictly between the two registers of a node of depth $j$.
\end{definition}

\begin{theorem}[Trees on blocks of any length with measurements]
    \label{thm:ldp_unequal_tree_rounds}
    \lean{MPSPreparation.treeLevelOp, MPSPreparation.treeLevelsOp,
        MPSPreparation.treeLeafOp, MPSPreparation.treeLevelsOp_succ',
        MPSPreparation.treeRegGate,
        MPSPreparation.two_mul_treeRegGate_L,
        MPSPreparation.treeRegGate_a_add, MPSPreparation.treeRegGate_sites,
        MPSPreparation.mem_span_treeRegGate,
        MPSPreparation.mem_interior_treeRegGate,
        MPSPreparation.treeLevelGates, MPSPreparation.mem_treeLevelGates,
        MPSPreparation.pairwise_treeLevelGates,
        MPSPreparation.disjoint_range_sites_interior_treeLevelGates,
        MPSPreparation.list_prod_map_op_treeLevelGates,
        MPSPreparation.registerLayersOp_treeLevelGates,
        MPSPreparation.treeCentralSites,
        MPSPreparation.exists_rounds_blockLayerOp_treeLevelsOp,
        MPSPreparation.isCircuitOn_blockLayerOp_treeLeafOp}
    \leanok
    \uses{def:ldp_tree_layout, thm:ldp_register_layers,
        thm:ldp_register_tree_rounds, def:qc_local_circuit,
        lem:qc_parallel_circuits}
    Cut a ring into blocks, each with a tree layout with $h+1$ coarse depths, and
    let a unitary $X_{j,p}$ on $2s$ sites act on the two registers of the node
    $p$ of depth $j\le h$ of every block. If every $X_{j,p}$ is a product of at
    most $K$ unitaries on neighbouring sites, measurement rounds of total depth
    $(h+1)(4s+K+2)$ implement the depths $0,\dots,h$ of the trees of all the
    blocks, the depth $0$ applied first, on the vectors with $\ket0$ at the
    sites $s,\dots,o_{2^{h+1}}-s-1$ of every block, $o_{2^{h+1}}$ the number of
    sites of its leaves. If every leaf carries a product of at most $K'$
    unitaries on neighbouring sites of the leaf, the leaves of all the blocks
    form a local circuit of depth $K'$.
\end{theorem}
\begin{proof}
    \leanok
    The unitary of a node is a register gate whose stretch is the node, without
    the sites after the leaves; the nodes of one depth are pairwise disjoint.
    By Definition~\ref{def:ldp_tree_layout} a register of a coarser depth is
    never strictly between the two registers of a finer node, so
    Theorem~\ref{thm:ldp_register_layers} applies the depths one after the
    other. The leaves are pairwise disjoint windows of consecutive sites, so
    their unitaries act in parallel.
\end{proof}

\begin{theorem}[The tree of a block of any length implements its isometry]
    \label{thm:ldp_unequal_tree_isometry}
    \lean{MPSPreparation.treeInput, MPSPreparation.treeInput_injective,
        MPSPreparation.nodeCfg, MPSPreparation.segCfg,
        MPSPreparation.leafRegWindow, MPSPreparation.leafRegWindow_val,
        MPSPreparation.leafRegWindow_injective₂, MPSPreparation.leafInput,
        MPSPreparation.leafInput_injective, MPSPreparation.treeNodeGate,
        MPSPreparation.treeLeafGate, MPSPreparation.treeBlockOp,
        MPSPreparation.treeNodeGate_mem_unitary,
        MPSPreparation.treeLeafGate_mem_unitary, MPSPreparation.treeBlockOp_mem_unitary,
        MPSPreparation.nodeLen_two_mul_add, MPSPreparation.two_mul_le_nodeLen,
        MPSPreparation.placeCfg_coarseOutput_eq_nodeCfg,
        MPSPreparation.nodeCfg_comp_leafWindow,
        MPSPreparation.treeLeafOp_apply_nodeCfg,
        MPSPreparation.treeLevelsOp_apply_nodeCfg,
        MPSPreparation.treeBlockOp_apply}
    \leanok
    \uses{def:ldp_tree_layout, thm:ldp_unitary_implementing_isometry}
    Let $V_m\colon\C^\chi\to(\C^d)^{\otimes m}$ and
    $W_{l,r}\colon\C^\chi\to\C^\chi\otimes\C^\chi$ be isometries for
    $m,l,r\ge s$, satisfying $V_{l+r}=(V_l\otimes V_r)W_{l,r}$.
    Let a register of $s$ sites carry $\C^\chi$ through an injective map of
    $\{0,\dots,\chi-1\}$ into its configurations, and let a block of $n$ sites
    carry a tree layout. Let the node $p$ of depth $j$ cover $\ell_{j,p}$ sites.
    The unitary of a node of depth $j\le h$ extends $W_{l,r}$ for its
    two child lengths, from its input to
    its two registers; the input of the root is placed on its two registers by
    a given injective map, and below the root the input is a register on the
    first register of an even node and on the last register of an odd one, the
    other in $\ket0$. The unitary of a leaf of $\ell$ sites extends $V_\ell$ from
    its input to all its configurations. The tree $T$ of these unitaries, the
    depths $0,\dots,h$ applied first and the leaves last, satisfies
    \begin{align*}
        \bra\tau T\ket{\iota_0(x)}=\bra\tau V_n\ket x
    \end{align*}
    for every configuration $\tau$ of the block, with $\iota_0(x)$ the input
    $x$ on the two registers of the root and $\ket0$ elsewhere.
\end{theorem}
\begin{proof}
    \leanok
    \uses{thm:ldp_unitary_implementing_isometry, lem:ldp_tree_amplitudes}
    By induction from the leaves: after the leaves and the depths $h,\dots,j$,
    the amplitude of $\tau$ on the inputs $x_p$ of the nodes of depth $j$ is
    $\prod_p\bra{\tau_{j,p}}V_{\ell_{j,p}}\ket{x_p}$, with $\tau_{j,p}$ the part
    of $\tau$ on the node. At the leaves this is the definition of their
    unitaries, since the leaves cover the block. A unitary of depth $j$ maps its
    input to its two registers with the amplitudes of $W$, and the registers it
    writes are the inputs of the two nodes below it, by the second item of
    Lemma~\ref{lem:ldp_tree_amplitudes}; summing over these inputs and
    applying $V_{l+r}=(V_l\otimes V_r)W_{l,r}$ at every node gives the claim for
    depth $j$. The root covers the block.
\end{proof}

\begin{theorem}[Coherent measurement implementation of bounded-leaf trees]
    \label{thm:ldp_supported_tree_rounds}
    \lean{QuantumCircuit.MeasurementRound.OutcomeHistory,
        QuantumCircuit.MeasurementRound.historyKraus,
        QuantumCircuit.MeasurementRound.historyKraus_mulVec_mem_outputs,
        QuantumCircuit.MeasurementRound.IsRoundsImplementationOn.exists_history_scalar,
        QuantumCircuit.zeroOnSubmodule, QuantumCircuit.mem_zeroOnSubmodule,
        MPSPreparation.exists_rounds_blockLayerOp_treeBlockOp,
        MPSPreparation.treeBlockOp_apply_supportedCfgPolarIso,
        MPSPreparation.exists_rounds_blockLayerOp_supportedPolarTree}
    \leanok
    \uses{thm:ldp_unequal_tree_isometry, thm:ldp_supported_polar_tree,
        def:qc_measurement_rounds, def:qc_round_implementation}
    Fix $d,s\ge1$ and a physical leaf bound $c$.
    There is $C$, depending only on $d,s,c$, such that any ring partitioned
    into blocks with tree layouts of $h+1$ coarse depths and leaves of at most
    $c$ sites has the following property. For the register tree $T_b$ of each
    block from Theorem~\ref{thm:ldp_unequal_tree_isometry}, measurement rounds
    implement $T=\bigotimes_bT_b$ in depth at most $C(h+1)$ on the subspace
    $E$ with every interior scratch site in $\ket0$.
    For every fixed outcome history $\mathbf m$, there is a scalar
    $a_{\mathbf m}$, independent of $v\in E$, such that
    \begin{align}
        K_{\mathbf m}v & =a_{\mathbf m}Tv.
        \label{eq:ldp_tree_branch_scalar}
    \end{align}
    The scalar may be zero. The equality includes every physical output site.
    In particular, substituting the supported polar maps of
    Theorem~\ref{thm:ldp_supported_polar_tree} gives one physical
    implementation on all supported labels and their coherent superpositions.
    Only nearest-neighbour unitary layers contribute to the depth;
    measurements, classical communication and on-site corrections are free.
    The number of rounds is allowed to grow with $h$.
\end{theorem}
\begin{proof}
    \leanok
    \uses{thm:ldp_unequal_tree_rounds, thm:qc_bounded_unitary_two_site_gates}
    Each node unitary acts on $2s$ sites, so a fixed number $K$ of
    nearest-neighbour gates suffices. Register teleportation applies every
    coarse depth in depth $4s+K+2$, independently of register separation.
    The finitely many leaf widths between $2$ and $c$ have a common exact
    synthesis bound, and the leaf windows are disjoint.
    Concatenating the coarse rounds and the leaf circuit gives $C(h+1)$.
    The gates extend the same isometries on every supported input.

    To verify the independence in (\ref{eq:ldp_tree_branch_scalar}), fix a
    history and its linear Kraus map $K$. The all-input implementation says
    $Kv\in\C Tv$ for every $v\in E$. Since $T$ is unitary,
    $T^\dagger K$ restricts to an endomorphism of $E$ preserving every
    one-dimensional subspace. Such an endomorphism is $aI$, including when
    $\dim E$ is zero or one. Thus $Kv=aTv$ with one scalar throughout $E$.
\end{proof}

\begin{lemma}[Leaves of a block of any length]
    \label{lem:ldp_balanced_widths}
    \lean{MPSPreparation.balancedWidths,
        MPSPreparation.leafOffset_balancedWidths,
        MPSPreparation.leafOffset_balancedWidths_self,
        MPSPreparation.isTreeLayout_balancedWidths,
        MPSPreparation.le_leafOffset_balancedWidths_add_one,
        MPSPreparation.leafLen_balancedWidths_le}
    \leanok
    \uses{def:ldp_tree_layout}
    Let $L\ge2^{h+2}s$. The widths $w_i=2\lfloor\lfloor L/2\rfloor/2^{h+1}\rfloor$,
    increased by $2$ for the first $\lfloor L/2\rfloor\bmod2^{h+1}$ leaves, are a
    tree layout of a block of $L$ sites with $h+1$ coarse depths, whose leaves
    cover all the sites but at most one. If $L\le2^{h+1}c$, every leaf has at most
    $c+3$ sites.
\end{lemma}
\begin{proof}
    \leanok
    The widths sum to $2\lfloor L/2\rfloor$, and
    $\lfloor\lfloor L/2\rfloor/2^{h+1}\rfloor\ge s$. If $L\le2^{h+1}c$, then
    $2\lfloor\lfloor L/2\rfloor/2^{h+1}\rfloor\le c$; the last leaf has one more
    site when $L$ is odd.
\end{proof}

\begin{theorem}[The state of blocks of unequal lengths with measurements]
    \label{thm:ldp_unequal_blocks_measurements}
    \lean{MPSPreparation.isZeroOn_pairLayerOp_mulVec_productVector,
        MPSPreparation.blockInputCfg_eq_placeCfg,
        MPSPreparation.mul_self_le_pow_of_isInjective_blockTensor,
        MPSPreparation.le_pow_sub_one_of_mul_self_le,
        MPSPreparation.exists_isPreparedWithMeasurementRoundsInDepth_blockIsometryState}
    \leanok
    \uses{thm:ldp_unequal_tree_isometry, thm:ldp_unequal_tree_rounds,
        lem:ldp_balanced_widths, def:ldp_block_isometry_state,
        thm:ldp_approximating_state_depth,
        thm:qc_bounded_unitary_two_site_gates, def:qc_measurement_rounds,
        lem:qc_round_implementation_compose}
    For every $s\ge2$ and $c$ there is $C$, depending only on $d$, $s$ and $c$,
    with the following property. Let $A$ be a tensor whose blocked tensors over
    $m\ge s$ sites are injective, and $\omega$ a unit vector on the pair space.
    For every $h\ge0$ and every cutting of a ring into $M\ge1$ blocks of lengths
    $2^{h+2}s\le\ell_k\le2^{h+1}c$, the state~(\ref{eq:ldp_block_isometry_state})
    of these blocks with the pair $\omega$ is prepared with measurement rounds in
    depth at most $C(h+1)$. No outcome is post-selected.
\end{theorem}
\begin{proof}
    \leanok
    Injectivity over $s$ sites gives $D^2\le d^s$, hence $D\le d^{s-1}$ since
    $s\ge2$, so a leg of a pair is written on the first $s-1$ sites of a
    register, the last site in $\ket0$, and a register carries $\C^{D^2}$. Cut
    every block into leaves as in Lemma~\ref{lem:ldp_balanced_widths}. The input
    $\ket{l,0,\dots,0,r}$ of a block, with $l$ on its first and $r$ on its last
    $s$ sites, lies on the two registers of its root: when the length is odd,
    the last site of the block lies after the leaves, and it carries the last
    site of the leg $r$, which is $\ket0$. A local circuit of bounded depth,
    applied in one round, prepares the pairs from the all-$\ket0$ product
    state; Theorem~\ref{thm:ldp_unequal_tree_rounds} applies the depths
    $0,\dots,h$ of the trees with the unitaries of
    Theorem~\ref{thm:ldp_unequal_tree_isometry}; and a local circuit of
    bounded depth, by Theorem~\ref{thm:qc_bounded_unitary_two_site_gates} for
    the boundedly many lengths of the leaves, applies the leaves in one round.
    The output is the state of the blocks as in
    Theorem~\ref{thm:ldp_approximating_state_depth}.
\end{proof}

\begin{theorem}[Preparation with measurements in depth $O(\log\log(N/\varepsilon))$, every chain length]
    \label{thm:ldp_loglog_depth_every_length}
    \lean{MPSPreparation.exists_isPreparedWithMeasurementRoundsInDepth_le_log_log_of_mpvState_ne_zero,
        MPSPreparation.exists_isPreparedWithMeasurementRoundsInDepth_le_log_log}
    \leanok
    \uses{def:qc_measurement_rounds}
    Let $A$ be normal. There is a constant $c$ depending only on $A$ such that,
    for every $N\ge2$ and $0<\varepsilon\le1$ for which the periodic state
    $\ket{\phi_N(A)}$ is nonzero, a unit vector $\ket\psi$ with
    $\varepsilon(\psi,\phi_N)\le\varepsilon$ is prepared with measurement rounds
    in depth at most $c\log(\log(N/\varepsilon)+1)$, which is
    $O(\log\log(N/\varepsilon))$. No outcome is post-selected. If $D\ge1$, there
    is $N_0$ such that the condition $\ket{\phi_N(A)}\neq0$ holds for every
    $N\ge N_0$. This is the depth of the paragraph ``Tree-RG circuit with
    measurements'' of~\cite{Malz2024Preparation} for normal tensors; the source
    normalizes $\ket{\phi_N}$, which presupposes $\ket{\phi_N(A)}\neq0$
    \cite{gap:mswc24_depth_upper_bound_nonzero_state}.
\end{theorem}
\begin{proof}
    \leanok
    \uses{thm:ldp_unequal_blocks_measurements,
        thm:ldp_block_approximation_error, thm:ldp_short_chain_preparation,
        thm:normalized_primitive_gauge, lem:ldp_normalized_state_rescaling,
        lem:ldp_periodic_state_nonzero,
        lem:ldp_approximating_state_unit_norm,
        lem:qc_measurement_rounds_outputs, lem:qc_measurement_unitary_case}
    Bring $A$ into the gauge~(\ref{eq:ldp_normal_gauge}), which multiplies
    $\ket{\phi_N}$ by a nonzero scalar, with a bound $t<1$ on the moduli of the
    subleading eigenvalues, $r=-\log t$, and an injectivity length $L$;
    put $s=L+2$, $a=1/r$, $b=\max(\log K,0)/r+4s+1$ with $K$ the constant of
    Theorem~\ref{thm:ldp_block_approximation_error} at $\gamma=1/2$, and
    $q=\lceil a\log(N/\varepsilon)+b\rceil\ge4s$. If $q\le N$, write
    $N=(M-1)q+q'$ with $q\le q'<2q$; the state of these blocks with the pair of
    the fixed point has error at most $KMe^{-rq}\le\varepsilon$ by
    Theorem~\ref{thm:ldp_block_approximation_error}, and with
    $2^{h+2}s\le q<2^{h+3}s$ it is prepared in depth $C(h+1)$ by
    Theorem~\ref{thm:ldp_unequal_blocks_measurements} with $c=8s$. If
    $4s\le N<q$, the chain is one block, and by
    Theorem~\ref{thm:ldp_short_chain_preparation} $\ket{\phi_N}$ is the state of
    this block with the pair $\omega\propto P\ket{\mathbb 1}$, prepared in depth
    $C(h+1)$ with $2^{h+2}s\le N<2^{h+3}s$. If $N<4s$, $\ket{\phi_N}$ is prepared
    by a local circuit of bounded depth. In each case $2^h\le
        a\log(N/\varepsilon)+b+1\le(a+b+1)(\log(N/\varepsilon)+1)$, and
    $\log(\log(N/\varepsilon)+1)\ge\log(1+\log2)>0$, so the depth is at most a
    constant times $\log(\log(N/\varepsilon)+1)$.
\end{proof}


%% ==================================================
%% Subsection: Long-range MPS using measurements
%% Sources: arXiv:2307.01696 (Malz, Styliaris, Wei, Cirac), paragraph
%% "Long-range MPS using measurements", eq. (19) and the paragraph after it;
%% arXiv:2103.13367 (Piroli, Styliaris, Cirac), Example 1.
%% Scope restrictions: orthogonal blocks, block lengths dividing N;
%% docs/paper-gaps/mswc24_block_form_mixed_overlap.tex,
%% docs/paper-gaps/mswc24_measurement_preparation_scope.tex.
%% ==================================================

\subsection{Tensors that are not normal, with measurements}
\label{sec:ldp_long_range_measurements}

For a tensor that is not normal, \cite{Malz2024Preparation} prepares the
fixed point $\ket{\Omega'}=W^{\otimes N/q}\ket{\chi_{N/q}}$ of
Lemma~\ref{thm:ldp_ghz_form} by first preparing $\ket{\chi_{N/q}}$ with
measurements and then applying the isometries
$W\colon\ket j\mapsto\ket{\omega_j}$ and the isometries of the blocked tensor
as circuits. With the block length $q\propto\log(N/\varepsilon)$, ``If instead
measurements are only used for the preparation of $\ket{\chi_{N/q}}$, the depth
is $O(\log(N/\varepsilon))$''. This subsection proves this for direct sums of
normal blocks whose $q$-site states are orthogonal and for block lengths
dividing $N$ \cite{gap:mswc24_measurement_preparation_scope};
Section~\ref{sec:ldp_overlapping_measurements} removes both restrictions. On the chain of
$N$ sites the label of each block is stored in the last sites of the block, and
the one-round construction prepares the GHZ-type state in depth $O(q)$,
while Lemma~\ref{thm:ldp_sparse_window_ghz} prepares it in constant depth
using additional measurement rounds. The sequential block isometries take
depth $O(q)$.

\begin{definition}[Registers and the GHZ-type state on them]
    \label{def:ldp_window_ghz_state}
    \lean{MPSPreparation.registerSite, MPSPreparation.ancillaSite,
        MPSPreparation.registerSite_eq, MPSPreparation.ancillaSite_eq,
        MPSPreparation.pairSite_eq_pairSite_iff, MPSPreparation.registerSite_inj,
        MPSPreparation.ancillaSite_inj, MPSPreparation.registerSite_ne_ancillaSite,
        MPSPreparation.registerSite_injective,
        MPSPreparation.blockSite_eq_ancillaSite, MPSPreparation.windowGHZState,
        MPSPreparation.registerCfg, MPSPreparation.registerCfg_registerSite,
        MPSPreparation.registerCfg_of_forall_ne, MPSPreparation.windowInput,
        MPSPreparation.windowInput_injective,
        MPSPreparation.registerCfg_comp_pairSite,
        MPSPreparation.registerCfg_of_forall_pairSite_ne,
        MPSPreparation.windowGHZState_extend}
    \leanok
    \uses{thm:ldp_approximating_state_depth, def:ldp_pair_family_fixed_point}
    Cut the ring of $N$ sites into $M\geq1$ blocks of lengths
    $\ell_k\geq2r_1$. The \emph{register} $R_k$ consists of the last $r_1$
    sites of block $k$, and its \emph{ancilla} of the first $r_1$ sites of
    block $k+1$, cyclically; together they form the pair window $k$. For a
    vector $\alpha$ on $(\C^d)^{\otimes r_1}$, the \emph{GHZ-type state on the
        registers} has the amplitude $\alpha(u)$ at a configuration carrying the
    same configuration $u$ on every register and $0$ at every other site, and
    the amplitude $0$ at every other configuration. If the labels $j$ of $b$
    blocks are encoded injectively as configurations $\delta(j)$ of $r_1$ sites
    and $\alpha(\delta(j))=\alpha_j$, $\alpha(u)=0$ for every other $u$, it is
    $\sum_j\alpha_j\ket{x_j}$, where $x_j$ carries $\delta(j)$ on every register
    and $0$ elsewhere: the state $\ket{\chi_M}$ of
    Definition~\ref{def:ldp_pair_family_fixed_point}, with the label of each
    block stored in its register. The configuration $x_j$ carries
    $(\delta(j),0,\dots,0)$ on every pair window and $0$ outside the windows.
\end{definition}

\begin{lemma}[Block and window layers acting on configurations]
    \label{lem:ldp_block_layers_configurations}
    \lean{MPSPreparation.blockLayerOp_mulVec_eq_comp,
        MPSPreparation.pairLayerOp_mulVec_eq_comp}
    \leanok
    \uses{lem:qc_configuration_layers, def:ldp_window_ghz_state,
        thm:ldp_approximating_state_depth}
    Cut the ring of $N$ sites into $M$ blocks of lengths $\ell_k$. If operators
    $U_k$ on the $\ell_k$ sites of the blocks act by maps $f_k$ of the
    configurations of the blocks, $U_k(v)=v\circ f_k$, then the layer of the
    $U_k$ on the blocks acts by the map changing the configuration of every
    block $k$ by $f_k$; the blocks cover the ring. If $\ell_k\geq2r_1$ and
    operators $W_k$ on $2r_1$ sites act by maps $f_k$, the layer of the $W_k$
    on the pair windows of Definition~\ref{def:ldp_window_ghz_state} acts by
    the map changing the configuration of every pair window $k$ by $f_k$ and
    keeping the other sites.
\end{lemma}
\begin{proof}
    \leanok
    The blocks are pairwise disjoint, and so are the pair windows; this is
    Lemma~\ref{lem:qc_configuration_layers}.
\end{proof}

\begin{lemma}[A controlled shift between the ends of a block]
    \label{lem:ldp_block_shift}
    \lean{MPSPreparation.copyShift, MPSPreparation.blockShift,
        MPSPreparation.tailShift, MPSPreparation.exists_blockShiftUnitary}
    \leanok
    \uses{lem:qc_configuration_layers, thm:qc_bounded_unitary_two_site_gates,
        thm:ldp_approximating_state_depth}
    Let $r_1\geq1$ and values be taken in $\mathbb Z_d$. There is $C$ such
    that for every $q\geq3r_1$ some unitary on $q$ sites, a product of at most
    $Cq$ gates on neighbouring sites, acts by the map of configurations that
    adds the last $r_1$ sites to the first $r_1$, site by site, and leaves the
    other sites unchanged.
\end{lemma}
\begin{proof}
    \leanok
    SWAP gates move the first $r_1$ sites next to the last $r_1$ in at most
    $2qr_1$ gates; one unitary on the last $2r_1$ sites, of bounded depth by
    Theorem~\ref{thm:qc_bounded_unitary_two_site_gates}, adds the second half
    to the first; and the same SWAP gates move the sites back. By
    Lemma~\ref{lem:qc_configuration_layers} the composite acts by the stated
    map.
\end{proof}

\begin{lemma}[The GHZ-type state on the registers with measurements]
    \label{thm:ldp_window_ghz_measurement}
    \lean{MPSPreparation.exists_isPreparedWithMeasurementsInDepth_windowGHZState}
    \leanok
    \uses{def:ldp_window_ghz_state, def:qc_measurement_preparation,
        lem:qc_configuration_layers, lem:ldp_block_layers_configurations,
        lem:ldp_block_shift, thm:qc_ghz_measurement}
    Let $r_1\geq2$. There is $C$ such that for every ring of $N$ sites cut into
    $M\geq1$ blocks of lengths $3r_1\leq\ell_k\leq L$ and every unit vector
    $\alpha$ on $(\C^d)^{\otimes r_1}$, the GHZ-type state on the registers of
    Definition~\ref{def:ldp_window_ghz_state} is prepared with measurements in
    depth at most $CL$. Following the construction of
    \cite[Example~1]{Piroli2021LOCC} for registers of $r_1$ sites, this prepares
    $\ket{\chi_{N/q}}$ in the paragraph ``Long-range MPS using measurements'' of
    \cite{Malz2024Preparation}; as recorded in
    \cite{gap:mswc24_measurement_preparation_scope}, the depth is $O(L)$ rather
    than constant because the register of a block and the ancilla at its start
    are $\ell_k-r_1$ sites apart.
\end{lemma}
\begin{proof}
    \leanok
    Start from the product vector with the all-ones vector on the sites of the
    registers $R_k$, $k\neq0$, and $\ket0$ elsewhere. One unitary on the sites
    of $R_0$, of bounded depth, sends $\ket{0\cdots0}$ to $\sum_u\alpha(u)\ket u$.
    One layer of gates on the pair windows adds each register to its ancilla,
    and on every block the unitary of Lemma~\ref{lem:ldp_block_shift} subtracts
    the register of the block from the ancilla at its start, which belongs to the
    register of the previous block. The circuit has depth $O(L)$. Measuring every
    ancilla, with outcomes $m_k\in\mathbb Z_d^{r_1}$ and partial sums
    $p_k=m_0+\cdots+m_{k-1}$, and correcting by adding $m_k$ on the ancilla of $R_k$ and
    $-p_k$ on $R_k$ site by site gives, by Lemma~\ref{lem:qc_configuration_layers},
    a corrected vector with amplitude $\alpha(x|_{R_0})$ at a
    configuration $x$ that is $0$ outside the registers and satisfies
    $m_k+(x|_{R_{k+1}}-p_{k+1})-(x|_{R_k}-p_k)=0$ for every $k$, cyclically, and
    amplitude $0$ elsewhere. These conditions hold exactly when
    $\sum_km_k=0$ and all the registers of $x$ agree. So every outcome with
    $\sum_km_k=0$ gives the GHZ-type state, and every other outcome gives the
    zero vector, which has probability zero.
\end{proof}

\begin{lemma}[Constant-depth GHZ preparation on sparse registers]
    \label{thm:ldp_sparse_window_ghz}
    \lean{MPSPreparation.exists_isPreparedWithMeasurementRoundsInDepth_windowGHZState,
        MPSPreparation.exists_rounds_sparseBlockShift,
        MPSPreparation.windowCentralSites,
        MPSPreparation.windowGHZInitial,
        MPSPreparation.exists_windowGHZSeedUnitary,
        MPSPreparation.windowGHZDifference,
        MPSPreparation.windowGHZDifference_eq_mulVec,
        MPSPreparation.exists_windowGHZCorrectionRound,
        MPSPreparation.SparseRegister.prefixEmbedding,
        MPSPreparation.SparseRegister.padding,
        MPSPreparation.SparseRegister.padding_first,
        MPSPreparation.SparseRegister.padding_of_ge,
        MPSPreparation.SparseRegister.exists_isPairProduct_padding,
        MPSPreparation.SparseRegister.paddingAll,
        MPSPreparation.SparseRegister.paddingAll_blockSite,
        MPSPreparation.SparseRegister.paddingAll_symm_blockSite,
        MPSPreparation.SparseRegister.exists_isCircuitOn_paddingAll,
        MPSPreparation.SparseRegister.ends,
        MPSPreparation.SparseRegister.ends_injective,
        MPSPreparation.SparseRegister.ends_shift_mulVec,
        MPSPreparation.SparseRegister.gate,
        MPSPreparation.SparseRegister.gate_span_subset,
        MPSPreparation.SparseRegister.pairwise_gates,
        MPSPreparation.SparseRegister.gate_sites,
        MPSPreparation.SparseRegister.gate_interior,
        MPSPreparation.SparseRegister.padding_conj_gates}
    \leanok
    \uses{thm:ldp_window_ghz_measurement, def:qc_measurement_rounds,
        thm:qc_register_layer}
    Fix $d\ge1$ and $r_1\ge2$. There is $C$ such that, for every cyclic
    partition into $M\ge1$ blocks of lengths $\ell_k\ge3r_1$ and every unit
    vector $\alpha$ on $(\C^d)^{\otimes r_1}$, the GHZ-type state of
    Definition~\ref{def:ldp_window_ghz_state} is prepared by measurement rounds
    of total quantum depth at most $C$. There is no upper bound on $\ell_k$.
    Every site outside the specified registers ends in $\ket0$. The statement
    includes $M=1$ and every combination of even and odd block lengths.
\end{lemma}
\begin{proof}
    \leanok
    Starting from the seed and copying gates of
    Lemma~\ref{thm:ldp_window_ghz_measurement}, in a block of length $q$ let
    $e=q\bmod2$ and $L=\lfloor(q-2r_1)/2\rfloor$, so
    $2L+2r_1+e=q$. When $e=1$, rotating the first $r_1+1$ sites sends
    $\ket{a_1\cdots a_{r_1},0}$ to $\ket{0,a_1\cdots a_{r_1}}$; for $e=0$
    no rotation is needed. The registers now occupy offsets
    $e,\ldots,e+r_1-1$ and $q-r_1,\ldots,q-1$, separated by $2L$ zero sites, and
    their teleportation stretches lie inside disjoint blocks. Applying the
    controlled subtraction of Lemma~\ref{thm:qc_register_layer} and undoing
    the rotations yields the controlled
    subtraction between the original end registers, with all central sites
    fixed. The padding rotations use only $r_1+1$ sites each, so both have
    constant depth; the register gates have depth $4r_1+K+2$ for a fixed
    $2r_1$-site gate bound $K$.

    Measuring and correcting the ancillas as in
    Lemma~\ref{thm:ldp_window_ghz_measurement} enforces the cyclic consistency
    condition $\sum_k m_k=0$; every other outcome has zero probability.
    The same correction acts on every label amplitude, preserving coherent
    superpositions. Quantum layers on neighbouring sites are counted;
    on-site measurements, outcome-dependent on-site corrections, and global
    classical communication are free. No postselection, additional physical
    sites, or free unitary on an entire block is used.
\end{proof}

\begin{lemma}[Sequential factorization on the first inputs]
    \label{lem:ldp_sequential_factorization_first_inputs}
    \lean{MPSPreparation.exists_isometric_chain_of_eq_mul_of_le,
        MPSPreparation.polarIsoMatrix_chainBlockTensor_eq_sum,
        MPSPreparation.polarIsoMatrix_blockTensor_eq_sum}
    \leanok
    \uses{thm:ldp_sequential_factorization_general}
    In Lemma~\ref{thm:ldp_sequential_factorization_general}, let the matrix
    elements of $V$ have the stated form only for the first $r$ inputs $x$, and let
    these $r\leq D^2$ columns of $V$ be orthonormal. Then the same conclusion
    holds with $b_q=r$, for these $r$ inputs. For an injective blocked
    tensor $B=VP$ of tensors $A_1,\dots,A_q$, one on each site,
    $V=BP^{-1}$ has the matrix elements
    $\bra{i_1\cdots i_q}V\ket x=\sum_{\alpha,\beta}
        (A_1^{i_1}\cdots A_q^{i_q})_{\alpha\beta}(P^{-1})_{(\alpha,\beta),x}$;
    for equal tensors $A_p=A$ this reads
    $\sum_{\alpha,\beta}(A^{i_1}\cdots A^{i_q})_{\alpha\beta}
        (P^{-1})_{(\alpha,\beta),x}$ with $B=\mathcal B_q$.
\end{lemma}
\begin{proof}
    \leanok
    The proof of Lemma~\ref{thm:ldp_sequential_factorization_general} applies
    to the first $r$ columns of $G$, the others being replaced by zero: the
    remainder of the sweep has orthonormal first $r$ columns because those of
    $V$ are, and absorbing it into the last site keeps that site isometric on
    its first $r$ levels.
\end{proof}

\begin{lemma}[One block unitary for blocks with orthogonal states]
    \label{thm:ldp_block_sum_unitary}
    \lean{MPSPreparation.exists_blockUnitary_of_equiv, MPSPreparation.flatCoord,
        MPSPreparation.flatCoord_injective, MPSPreparation.flatCoord_ne,
        MPSPreparation.evalWord_blockSum_flatCoord,
        MPSPreparation.sum_star_polarIsoMatrix_mul,
        MPSPreparation.sum_mul_self_le_pow, MPSPreparation.exists_blockSumUnitary}
    \leanok
    \uses{lem:ldp_sequential_factorization_first_inputs,
        thm:ldp_approximating_state_depth, lem:ldp_orthogonal_sum_polar}
    Let $A_1,\dots,A_b$ be tensors of bond dimensions $D_j$, and place their bond
    indices side by side, so that the bond index $a$ of block $j$ becomes the
    index $\iota_j(a)$ among $D_1+\cdots+D_b\geq1$ indices; encode these
    indices injectively in $r_1\geq1$ sites. There is $C$ such that for every
    block length $q\geq3r_1$ at which every blocked tensor $B_j=V_jP_j$ is
    injective and $B_j^\dagger B_{j'}=0$ for $j\neq j'$, some unitary $U$ on $q$
    sites, a product of at most $Cq$ gates on neighbouring sites, satisfies
    $U\ket{\iota_j(l),0,\iota_j(r)}=V_j\ket{l,r}$ for all $j$, $l$ and $r$. The
    columns $V_j\ket{l,r}$ are jointly orthonormal, so
    $\sum_jD_j^2\leq d^q$.
\end{lemma}
\begin{proof}
    \leanok
    The isometries of distinct blocks have orthogonal ranges because
    $B_j^\dagger B_{j'}=0$, so the map $(j,l,r)\mapsto V_j\ket{l,r}$ has
    orthonormal columns, and the rank of the matrix of these columns gives
    $\sum_jD_j^2\leq d^q$. Its matrix elements are those of the direct sum
    $\bigoplus_jA_j$ with the matrices $P_j^{-1}$ absorbed on the right, by
    Lemma~\ref{lem:ldp_sequential_factorization_first_inputs}. Ordering the
    bond pairs so that the pairs $(\iota_j(l),\iota_j(r))$ come first, the
    lemma factors this map, and the staircase of
    Lemma~\ref{thm:ldp_approximating_state_depth}, for this ordering, implements
    it in $O(q)$ gates.
\end{proof}

\begin{lemma}[The fixed point of orthogonal blocks with measurements]
    \label{thm:ldp_orthogonal_blocks_measurement}
    \lean{MPSPreparation.exists_isPreparedWithMeasurementsAndCircuitInDepth_sum_blockIsometryState_of_encoding,
        MPSPreparation.exists_isPreparedWithMeasurementsAndCircuitInDepth_sum_blockIsometryState}
    \leanok
    \uses{def:qc_measurement_circuit_preparation, thm:ldp_window_ghz_measurement,
        thm:ldp_block_sum_unitary, def:ldp_block_isometry_state,
        thm:ldp_approximating_state_depth}
    For every family of bond dimensions $D_1,\dots,D_b$ there are $C$ and $L_0$
    such that the following holds. Let $A_j$ be tensors of bond dimensions $D_j$,
    $\omega_j$ unit vectors on $\C^{D_j}\otimes\C^{D_j}$ and $\alpha$ a unit
    vector in $\C^b$. Cut the ring of $N$ sites into $M\geq1$ blocks of lengths
    $L_0\leq\ell_k\leq L$ such that every blocked tensor of every $A_j$ is
    injective and $B_j^\dagger B_{j'}=0$ for $j\neq j'$ on every block. For
    blocks with orthogonal states, this prepares
    $V^{\otimes N/q}W^{\otimes N/q}\ket{\chi_{N/q}}$ in the paragraph
    ``Long-range MPS using measurements'' of \cite{Malz2024Preparation}:
    $\sum_j\alpha_j\ket{\psi_{\ell,\omega_j}(A_j)}$, with the states of
    Definition~\ref{def:ldp_block_isometry_state}, is prepared with
    measurements and a circuit in depth at most $CL$.
\end{lemma}
\begin{proof}
    \leanok
    Encode the labels and the bond indices in $r_1=b+\sum_jD_j+2$ sites; for
    $d\geq2$ this is possible since $r_1<d^{r_1}$, and for $d=1$
    Lemma~\ref{thm:ldp_block_sum_unitary} gives $\sum_jD_j^2\leq1$, so $b\leq1$
    and $\sum_jD_j\leq1$. Lemma~\ref{thm:ldp_window_ghz_measurement} prepares
    the GHZ-type state on the registers with $\alpha(\delta(j))=\alpha_j$.
    The columns $\ket{\omega_j}$, placed on the bond indices of block $j$ and
    encoded on a pair window, are orthonormal, so one unitary $W$ on a pair
    window sends $\ket{\delta(j),0}$ to them. Applying $W$ on every window and the
    unitary of Lemma~\ref{thm:ldp_block_sum_unitary} on every block acts on the
    GHZ-type state $\sum_j\alpha_j\ket{x_j}$ of
    Definition~\ref{def:ldp_window_ghz_state}: the window layer sends $\ket{x_j}$ to the state of
    a window layer preparing $\ket{\omega_j}$ from $\ket{0\cdots0}$, and the
    computation of Lemma~\ref{thm:ldp_approximating_state_depth} gives
    $\ket{\psi_{\ell,\omega_j}(A_j)}$.
\end{proof}

\begin{lemma}[The corrected approximating state with measurements]
    \label{thm:ldp_repeated_block_measurement}
    \lean{MPSPreparation.copyApproxVector_eq_sum,
        MPSPreparation.sum_norm_sq_copyApproxVector,
        MPSPreparation.exists_isPreparedWithMeasurementsAndCircuitInDepth_copyApproxVector}
    \leanok
    \uses{thm:ldp_orthogonal_blocks_measurement, def:ldp_repeated_block_state,
        thm:ldp_repeated_block_polar, thm:ldp_approximating_state_depth}
    There are $C$ and $L_0$, depending only on $d$ and the bond dimensions
    $D_1,\dots,D_b$, such that the following holds. Let $A$ be the direct sum
    with multiplicities of Definition~\ref{def:ldp_repeated_block_sum}, let
    $\sigma_j\geq0$ with $\Tr\sigma_j=1$, let $q\geq L_0$ be a block length at
    which every blocked tensor of every block is injective and
    $B_j^\dagger B_{j'}=0$ for $j\neq j'$, let $M\geq1$, and let $\alpha$ be a
    unit vector. Then the corrected approximating state
    $V^{\otimes M}\sum_j\alpha_jL_j^{\otimes M}\ket{\Omega_j}$ of
    Definition~\ref{def:ldp_repeated_block_state}, on $N=Mq$ sites, equals
    $\sum_j\alpha_j\ket{\psi_{q,\omega_j}(A_j)}$ with $\omega_j$ the pair of
    $\sigma_j$, is a unit vector, and is prepared with measurements and a
    circuit in depth at most $Cq$. Under the scope restriction to blocks whose
    $q$-site states are orthogonal recorded in
    \cite{gap:mswc24_block_form_mixed_overlap,gap:mswc24_repeated_block_corrected_state},
    this establishes the preparation in the paragraph ``Long-range MPS using
    measurements'' of \cite{Malz2024Preparation} for the corrected
    approximating state.
\end{lemma}
\begin{proof}
    \leanok
    By Lemma~\ref{thm:ldp_repeated_block_polar} the corrected approximating
    state is $\sum_j\alpha_j\ket{\phi_M(V_jP_{j,\infty})}$, and each term is the
    approximating state of the block $A_j$, which is
    $\ket{\psi_{q,\omega_j}(A_j)}$ by
    Lemma~\ref{thm:ldp_approximating_state_depth}. With
    $\ket{\psi_j}=\ket{\psi_{q,\omega_j}(A_j)}$, each $\ket{\psi_j}$ is a unit
    vector, and for $j\neq j'$ the hypothesis $B_j^\dagger B_{j'}=0$ makes the
    physical matrices of the approximating tensors of $A_j$ and $A_{j'}$
    orthogonal, so that the overlap of the two states on $M$ blocks vanishes.
    Hence $\braket{\psi_j}{\psi_{j'}}=\delta_{jj'}$ and
    \begin{align}
        \Bigl\lVert\sum_j\alpha_j\ket{\psi_j}\Bigr\rVert^2
        =\sum_{j,j'}\overline{\alpha_j}\alpha_{j'}\braket{\psi_j}{\psi_{j'}}
        =\sum_j\lvert\alpha_j\rvert^2=1.
        \notag
    \end{align}
    Lemma~\ref{thm:ldp_orthogonal_blocks_measurement} with $\ell_k=q$
    prepares the sum.
\end{proof}

\begin{lemma}[A common gap for normal blocks]
    \label{lem:ldp_common_gap}
    \lean{MPSPreparation.exists_forall_eigenvalue_norm_le}
    \leanok
    \uses{lem:ldp_normal_gap}
    For finitely many normal left-canonical tensors $A_1,\dots,A_b$, $b\geq1$, of
    positive bond dimensions, there is $0<t<1$ bounding the moduli of the
    eigenvalues other than $1$ of every transfer map $\E_{A_j}$.
\end{lemma}
\begin{proof}
    \leanok
    Lemma~\ref{lem:ldp_normal_gap} gives such a bound $t_j$ for every block;
    taking $t=\max_jt_j$ bounds them all.
\end{proof}

\begin{lemma}[Tensors that are not normal, with measurements, in depth
        $O(\log(N/\varepsilon))$]
    \label{thm:ldp_long_range_measurement_depth}
    \lean{MPSPreparation.exists_isPreparedWithMeasurementsAndCircuitInDepth_le_log_repeatedBlockSum}
    \leanok
    \uses{thm:ldp_repeated_block_measurement,
        thm:ldp_approximation_error_repeated_orthogonal_blocks, lem:ldp_common_gap,
        thm:ldp_depth_upper_bound_divisible, lem:ldp_depth_log_window,
        lem:ldp_block_length_threshold_error}
    Let $A$ be the direct sum with multiplicities of
    Definition~\ref{def:ldp_repeated_block_sum}, with $b\geq1$ blocks, every
    block $A_j$ normal in the gauge (\ref{eq:ldp_normal_gauge}) with a positive
    definite fixed point $\sigma_j$ of trace one. There are $a>0$, $b'\geq1$ and
    $c$, depending only on $A$, with the following property. Let $N\geq2$,
    $0<\varepsilon\leq1$, and let $q$ be a block length dividing $N$ with
    $a\log(N/\varepsilon)+b'\leq q\leq2(a\log(N/\varepsilon)+b')$, at which
    $B_j^\dagger B_{j'}=0$ for $j\neq j'$, such that the weights
    $\beta_j=\sum_k\mu_{j,k}^N$ are not all zero. Then a unit vector $\ket\psi$
    with $\varepsilon(\psi,\phi_N)\leq\varepsilon$ is prepared with measurements
    and a circuit in depth at most $c\log(N/\varepsilon)$. As recorded in
    \cite{gap:mswc24_measurement_preparation_scope}, this establishes the depth
    $O(\log(N/\varepsilon))$ of the paragraph ``Long-range MPS using
    measurements'' of \cite{Malz2024Preparation}, with measurements used only
    for the GHZ-type state, for the corrected approximating state, restricted to
    blocks whose $q$-site states are orthogonal and to block lengths dividing $N$.
    Because $q$ must divide $N$
    and lie in the window above, the statement applies only to chain lengths
    with a divisor in the window and is vacuous for the others, for instance for
    every large prime $N$.
\end{lemma}
\begin{proof}
    \leanok
    Let $t$ be the common gap of Lemma~\ref{lem:ldp_common_gap}, $r=-\tfrac14\log t$,
    and $K$ the constant of
    Lemma~\ref{thm:ldp_approximation_error_repeated_orthogonal_blocks} at
    $\gamma=1/4$, so the error of the corrected approximating state with the
    amplitudes $\alpha_j=\beta_j/(\sum_l\lvert\beta_l\rvert^2)^{1/2}$ is at most
    $KMe^{-rq}$. With $a=1/r$ and $b'=a\max(\log K,0)+L_0+L_1+1$, where $L_0$ is the
    length of Lemma~\ref{thm:ldp_repeated_block_measurement} and $L_1$ is the sum
    of the injectivity lengths of the blocks, one has $q\geq L_0$, every blocked
    tensor is injective, and the state is prepared in depth at most $Cq$. Since
    $N=Mq$, $rq=q/a$ and $a\log(N/\varepsilon)+b'\leq q$, the error is at most
    $\varepsilon$ by Lemma~\ref{lem:ldp_block_length_threshold_error}, and
    $Cq\leq c\log(N/\varepsilon)$ with $c=C(2a+2b'/\log2)$ by
    Lemma~\ref{lem:ldp_depth_log_window}.
\end{proof}

\begin{theorem}[All translation-invariant MPS with measurements]
    \label{thm:ldp_all_ti_measurement_source}
    \uses{thm:ldp_long_range_measurement_depth, thm:ldp_overlapping_measurement_depth}
    Every translation-invariant MPS of constant bond dimension, short- or
    long-range correlated, is prepared with error $\varepsilon$ in depth
    $O(\log(N/\varepsilon))$ by local circuits assisted by measurements, the
    measurements being used only for the preparation of the GHZ-type state
    \cite[paragraph ``Long-range MPS using measurements'']{Malz2024Preparation}.
    Lemma~\ref{thm:ldp_long_range_measurement_depth} proves this for direct
    sums of normal blocks whose $q$-site states are orthogonal and for block
    lengths dividing $N$. Lemma~\ref{thm:ldp_overlapping_measurement_depth}
    proves it for direct sums of normal blocks in the canonical form
    \cite[eq.~(S2)]{Malz2024Preparation} whose $q$-site states may overlap,
    with every multiplicity and every nonzero complex weight, and for every
    chain length $N\geq2$ at which $\ket{\phi_N(A)}\neq0$, with the mixed
    transfer maps of distinct blocks of spectral radius below one, the
    isometries of the direct sum with weights one, and the GHZ-type state
    prepared in depth $O(q)$ rather than in constant depth
    \cite{gap:mswc24_measurement_preparation_scope}. The same paragraph claims
    the depth $O(\log\log(N/\varepsilon))$ when the tree-RG circuit is also
    assisted by measurements; for normal tensors and the chain lengths
    divisible by $s2^{k+1}$ for some $k$ with $a\log(N/\varepsilon)+b\le s2^{k+1}\le2(a\log(N/\varepsilon)+b)$ this is Lemma~\ref{thm:ldp_loglog_depth}, and for
    tensors that are not normal it is not formalized.
\end{theorem}

%% ==================================================
%% Subsection: Overlapping blocks with measurements, for every chain length
%% Sources: arXiv:2307.01696 (Malz, Styliaris, Wei, Cirac), paragraph
%% "Long-range MPS using measurements", Supplemental Material, eqs. (S2)-(S7),
%% Lemma 1'(ii) and the proof of Theorem 1; arXiv:2103.13367 (Piroli,
%% Styliaris, Cirac), Example 1 and Supplemental Material, eqs. final_eq to
%% finished.
%% Local fixes: isometry of the unweighted direct sum, rate of the overlapping
%% blocks, nonvanishing periodic state;
%% docs/paper-gaps/mswc24_measurement_preparation_scope.tex,
%% docs/paper-gaps/mswc24_block_form_mixed_overlap.tex,
%% docs/paper-gaps/mswc24_depth_upper_bound_nonzero_state.tex.
%% Scope restriction: depth O(q) of the GHZ-type state;
%% docs/paper-gaps/mswc24_measurement_preparation_scope.tex.
%% ==================================================

\subsection{Overlapping blocks with measurements, for every chain length}
\label{sec:ldp_overlapping_measurements}

The preparation of Lemma~\ref{thm:ldp_long_range_measurement_depth} applies
the isometry of the blocked tensor of every block separately, which needs the
$q$-site states of distinct blocks to be orthogonal, and cuts the chain into
equal blocks. This subsection removes both restrictions. It applies, on every
block, the isometric factor $V$ of the polar decomposition of the blocked map of
the direct sum $\bigoplus_jA_j$ with weights one and one copy of each block. For
large block lengths this map is injective on the bond pairs of the blocks even
when their states overlap, so $V$ is an isometry there, and the weights of the
canonical form enter only through the amplitudes of the GHZ-type state. The
chain is cut into blocks ``all of the same size, $q_N$, except for the last one,
which may be larger'', as in the proof of Theorem~1 of
\cite{Malz2024Preparation}. The bound on the error is relative to the norm of
the weights, so no factor $\min(1,\sum_j\lvert\beta_j\rvert^2)^{-1/2}$ enters
\cite{gap:mswc24_measurement_preparation_scope}.

Throughout, $A_1,\dots,A_b$ are normal tensors in the gauge
(\ref{eq:ldp_normal_gauge}) with positive definite fixed points $\rho_j$ of trace
one, placed on disjoint bond coordinates by isometries $E_j$, $K_j=E_j\otimes
    E_j$, and $\lambda_2\in\C$ bounds the moduli of the eigenvalues other than one of
every transfer map $\E_{A_j}$ and the moduli of all eigenvalues of the mixed
transfer maps $\E_{jj'}$, $j\neq j'$, with correlation length $\xi$. For a block
length $\ell$, $\mathcal B^{(\ell)}=V^{(\ell)}P^{(\ell)}$ is the polar
decomposition of the $\ell$-site blocked map of $\bigoplus_jA_j$.

\begin{lemma}[Injectivity on the bond pairs of the blocks]
    \label{lem:ldp_overlapping_block_injectivity}
    \lean{MPSTensor.offBlockProj,
        MPSTensor.exists_pos_algebraMap_le_blockSumGramLimit_add,
        MPSTensor.isInjectiveOn_of_eq_zero_of_mulVec_eq_zero,
        MPSTensor.exists_isInjectiveOn_blockTensor_blockSum}
    \leanok
    \uses{lem:ldp_overlapping_positive_part, def:ldpolar_injective_on}
    Let $S$ be the set of bond pairs $(E_ja,E_jc)$ whose two indices lie in the
    same block.
    \begin{enumerate}
        \item If the matrices of a tensor $B$ vanish outside a set $S$ of bond
              pairs and $\mathcal Bx=0$ has no solution $x\neq0$ supported on
              $S$, then $B$ is injective on $S$.
        \item With $|\lambda_2|<1$ there is $L_0$ such that for every $\ell\geq
                  L_0$ the $\ell$-site blocked tensor of $\bigoplus_jA_j$ is
              injective on $S$.
    \end{enumerate}
\end{lemma}
\begin{proof}
    \leanok
    For part~(1), if a matrix unit $\ket\alpha\bra\beta$ with $(\alpha,\beta)\in
        S$ were not in the span of the matrices of $B$, a linear functional $f$
    vanishing on that span but not on $\ket\alpha\bra\beta$ would give the vector
    $x_p=f(\ket{p_1}\bra{p_2})$ for $p\in S$, and $x_p=0$ otherwise, with
    $\mathcal Bx=0$ and $x\neq0$. For part~(2), let
    $Q=\Id-\sum_jK_jK_j^\dagger$ and $G_\infty$ be as in
    Lemma~\ref{lem:ldp_overlapping_positive_part}; that lemma gives $c>0$ with
    $G_\infty+Q\geq c\Id$ and
    $\lVert\mathcal B^\dagger\mathcal B-G_\infty\rVert\leq Ke^{-\ell/(2\xi_t)}$,
    where $\xi_t$ is the correlation length at $t=\max(\lvert\lambda_2\rvert,1/2)$,
    which satisfies $0<t<1$ and bounds the same eigenvalues as $\lvert\lambda_2\rvert$.
    The matrices of the blocked tensor vanish outside $S$, and $Qx=0$ for $x$
    supported on $S$. For $\ell$ with $Ke^{-\ell/(2\xi_t)}<c$ and such $x$ with
    $\mathcal Bx=0$,
    \begin{align*}
        0=x^\dagger(\mathcal B^\dagger\mathcal B+Q)x
        \geq\bigl(c-\lVert\mathcal B^\dagger\mathcal B-G_\infty\rVert\bigr)
        x^\dagger x,
    \end{align*}
    so $x=0$, and part~(1) applies.
\end{proof}

\begin{lemma}[The states of the blocks through the isometries of the direct sum]
    \label{lem:ldp_overlapping_block_states}
    \lean{MPSTensor.blockPosTensor, MPSTensor.blockTensor_blockSum_single,
        MPSTensor.mpvState_eq_blockIsoVector,
        MPSTensor.mem_blockPairs_of_pairProductState_embedPair_ne_zero,
        MPSTensor.blockIsometryState_eq_blockIsoVector,
        MPSTensor.inner_blockIsometryState_embedPair_mpvState,
        MPSTensor.inner_blockIsometryState_embedPair,
        MPSTensor.norm_sum_smul_blockIsometryState}
    \leanok
    \uses{lem:ldp_overlapping_block_injectivity, def:ldp_block_isometry_state,
        def:ldp_chain_pos_state, thm:ldp_chain_block_iso_inner,
        lem:ldp_overlapping_overlaps}
    Cut a ring of $N\geq1$ sites into $M\geq1$ blocks of lengths $\ell_k\geq1$,
    and let $\omega_j$ be the pair of $\rho_j$ placed along $E_j$.
    \begin{enumerate}
        \item The $\ell$-site blocked map of $A_j$, placed along $E_j$, is
              $V^{(\ell)}P^{(\ell)}K_jK_j^\dagger$, so
              $\ket{\phi_N(A_j)}=\bigl(\bigotimes_kV^{(\ell_k)}\bigr)
                  \ket{\phi_M(P^{(\ell_k)}K_jK_j^\dagger)}$, where the second
              factor is the periodic state of the family of tensors
              $P^{(\ell_k)}K_jK_j^\dagger$ on the blocks.
        \item If every blocked map $\mathcal B^{(\ell_k)}$ is injective on the
              bond pairs of the blocks, then the states
              $\ket{\psi_j}=\bigl(\bigotimes_kV^{(\ell_k)}\bigr)
                  \bigotimes_k\ket{\omega_j}$ are orthonormal, so
              $\sum_j\alpha_j\ket{\psi_j}$ is a unit vector for every unit vector
              $\alpha$, and
              $z_{jj'}=\braket{\psi_j}{\phi_N(A_{j'})}$ is the trace of the
              ordered product over the blocks of the mixed transfer matrices of
              $P^{(\ell_k)}K_{j'}K_{j'}^\dagger$ against the fixed-point tensor
              of $E_j\rho_jE_j^\dagger$.
    \end{enumerate}
\end{lemma}
\begin{proof}
    \leanok
    For part~(1), the blocked map of $A_j$ placed along $E_j$ is
    $\mathcal B_jK_j^\dagger$, and that of the direct sum is
    $\sum_{j'}\mathcal B_{j'}K_{j'}^\dagger$; since
    $K_{j'}^\dagger K_j=\delta_{jj'}\Id$, the first is the second followed by
    $K_jK_j^\dagger$. Every word of length $N$ is a product of the blocked
    words. For part~(2), the product of the pairs $\omega_j$ vanishes on every
    configuration of bond pairs leaving the pairs of the blocks, and on these
    $V^{(\ell_k)}$ is an isometry, so the inner products are those of the
    products of pairs, which are $\delta_{jj'}$ because the $\omega_j$ are unit
    vectors on disjoint coordinates. The same isometry property and part~(1)
    give $z_{jj'}$ as the overlap of $\bigotimes_k\ket{\omega_j}$ with the
    periodic state of the family $P^{(\ell_k)}K_{j'}K_{j'}^\dagger$, which is the
    stated trace by Lemma~\ref{lem:ldp_overlapping_overlaps}(1).
\end{proof}

\begin{lemma}[The overlaps of the blocks near the limit]
    \label{lem:ldp_overlapping_block_overlaps_limit}
    \lean{Matrix.exists_norm_trace_prod_ofFn_sub_trace_pow_le,
        MPSTensor.mixedMapLM_blockSumPosLimit_mul_apply,
        MPSTensor.exists_norm_trace_prod_blockPosTensor_sub_le}
    \leanok
    \uses{lem:ldp_overlapping_block_states, lem:ldp_overlapping_positive_part,
        lem:ldp_overlapping_overlaps_weights, lem:ldp_shared_overlap_estimates}
    \begin{enumerate}
        \item For a finite index set there is $K\geq0$ such that for every
              matrix $E$ whose powers have norm at most $c$ and all matrices
              $X_0,\dots,X_{M-1}$ within $\delta$ of $E$,
              $\lvert\Tr(X_0\cdots X_{M-1})-\Tr E^M\rvert\leq Kc^2(M\delta)
                  e^{cM\delta}$.
        \item The mixed transfer map of $P_\infty K_{j'}K_{j'}^\dagger$ against
              the fixed-point tensor of $E_j\rho_jE_j^\dagger$ is
              $X\mapsto\Tr(E_jE_j^\dagger X)\,E_j\rho_jE_j^\dagger$ for $j=j'$
              and zero for $j\neq j'$.
        \item With $|\lambda_2|<1$ and $0<\gamma<1$, for all $j,j'$ there is
              $C>0$ such that for every cut into $M\geq1$ blocks of lengths
              $\ell_k\geq q\geq1$, with $y=Me^{-\gamma q/\xi}$, the trace
              $z_{jj'}$ of Lemma~\ref{lem:ldp_overlapping_block_states}(2)
              satisfies $\lvert z_{jj'}-\delta_{jj'}\rvert\leq Cy\,e^{Cy}$.
    \end{enumerate}
\end{lemma}
\begin{proof}
    \leanok
    Part~(1) is the telescoping estimate of
    Lemma~\ref{lem:ldp_shared_overlap_estimates}(8) for a product of different
    matrices, followed by the bound on the trace; it follows the argument of
    \cite[Supplemental Material]{Piroli2021LOCC}. Part~(2) is
    Lemma~\ref{lem:ldp_overlapping_overlaps_weights}(1) with $t=e_{j'}$, since
    $P_\infty K_{j'}K_{j'}^\dagger$ is the limit for these weights. For
    part~(3), the limit $E$ of the mixed transfer matrices is, by part~(2), an
    idempotent of trace one for $j=j'$ and zero otherwise, so its powers are
    bounded and $\Tr E^M=\delta_{jj'}$. By
    Lemma~\ref{lem:ldp_overlapping_positive_part},
    $\lVert P^{(\ell_k)}-P_\infty\rVert\leq K'e^{-\gamma\ell_k/\xi}\leq
        K'e^{-\gamma q/\xi}$, and part~(1) applies.
\end{proof}

\begin{lemma}[Error for overlapping blocks of unequal lengths and all weights]
    \label{thm:ldp_overlapping_unequal_blocks_error}
    \lean{MPSTensor.exists_one_sub_norm_inner_sum_blockIsometryState_le}
    \leanok
    \uses{lem:ldp_overlapping_block_states, lem:ldp_overlapping_block_overlaps_limit,
        lem:ldp_shared_overlap_estimates, thm:ldp_approximation_error_repeated_overlapping,
        def:correlation_length}
    Let $0<\gamma<1$. There is $C>0$ such that the following holds for all
    weights $\beta\in\C^b\setminus\{0\}$ and every cut of a ring of $N$ sites
    into $M\geq1$ blocks of lengths $\ell_k\geq q$ at which every blocked map
    $\mathcal B^{(\ell_k)}$ is injective on the bond pairs of the blocks. With
    $\alpha_j=\beta_j/(\sum_l\lvert\beta_l\rvert^2)^{1/2}$ and
    $\ket\phi=\sum_j\beta_j\ket{\phi_N(A_j)}$,
    \begin{align*}
        1-\Bigl\lvert\Bigl\langle\sum_j\alpha_j\bigl(\textstyle\bigotimes_kV^{(\ell_k)}\bigr)
        \textstyle\bigotimes_k\omega_j\Bigm|\frac{\phi}{\lVert\phi\rVert}\Bigr\rangle\Bigr\rvert
        \leq CMe^{-\gamma q/\xi}.
    \end{align*}
    Here $-\gamma q/\xi$ is read as $0$ when $\xi=0$, that is when
    $\lambda_2=0$; the bound then gives no decay in $q$. The constant does not
    depend on the weights. For the canonical form
    $A^i=\bigoplus_j\operatorname{diag}(\mu_{j,1},\dots,\mu_{j,m_j})\otimes A_j^i$
    of \cite[eq.~(S2)]{Malz2024Preparation}, $\ket{\phi_N(A)}$ is $\ket\phi$
    with $\beta_j=\sum_k\mu_{j,k}^N$, so this bounds the error for every
    multiplicity and every complex weight. Relative to
    \cite[Lemma~1$'$(ii)]{Malz2024Preparation} for the state of the paragraph
    ``Long-range MPS using measurements'', with blocks of unequal lengths, this
    replaces $\xi_{\mathrm{diag}}$ by
    $\max(\xi_{\mathrm{diag}},\xi_{\mathrm{off\text{-}diag}})$ as recorded in
    \cite{gap:mswc24_block_form_mixed_overlap}, and uses the isometries of the
    direct sum with weights one in place of that of the blocked map of $A$ as
    recorded in \cite{gap:mswc24_measurement_preparation_scope}. Unlike
    Lemma~\ref{thm:ldp_approximation_error_repeated_overlapping}, it has no
    factor $\min(1,\sum_j\lvert\beta_j\rvert^2)^{-1/2}$.
    The source states the bound for $0<\gamma<1/2$; the range $\gamma<1$ is a
    project result.
\end{lemma}
\begin{proof}
    \leanok
    If $\lambda_2=0$, $\lvert\lambda_2\rvert\geq1$ or $q=0$, then
    $Me^{-\gamma q/\xi}\geq1$ and
    the bound holds with $C=1$. Otherwise let $w=\sum_l\lvert\beta_l\rvert^2$
    and $\ket{\psi}=\sum_j\alpha_j\ket{\psi_j}$. By
    Lemma~\ref{lem:ldp_overlapping_block_states}(2),
    $\braket\psi\phi=w^{-1/2}S$ with
    $S=\sum_{j,j'}\overline{\beta_j}\beta_{j'}z_{jj'}$. Since
    $\lvert\beta_j\beta_{j'}\rvert\leq w$,
    Lemma~\ref{lem:ldp_overlapping_block_overlaps_limit}(3) gives
    $\lvert S/w-1\rvert\leq b^2C_0y\,e^{C_0y}$, a bound relative to $w$, with
    $y=Me^{-\gamma q/\xi}$.
    Lemma~\ref{lem:ldp_shared_overlap_estimates}(6) with $R=\sqrt w$ gives
    $\lvert\lVert\phi\rVert^2/w-1\rvert\leq Ke^{-\gamma N/\xi}\leq Ky$. The
    elementary inequalities in the proof of
    Lemma~\ref{thm:ldp_approximation_error_repeated_overlapping}, applied to
    $S/w$ and $\lVert\phi\rVert/\sqrt w$ with $s=1$, give the bound.
\end{proof}

\begin{lemma}[The state of a tensor injective on a set of bond pairs, with
        measurements]
    \label{thm:ldp_injective_on_measurement}
    \lean{MPSPreparation.exists_isPreparedWithMeasurementsAndCircuitInDepth_blockIsometryState_of_isInjectiveOn,
        MPSPreparation.isPreparedWithMeasurementsAndCircuitInDepth_of_isPreparedInDepth}
    \leanok
    \uses{def:qc_measurement_circuit_preparation, thm:ldp_window_ghz_measurement,
        thm:ldp_block_sum_unitary, thm:ldp_sequential_factorization_on,
        def:ldp_block_isometry_state, lem:qc_measurement_unitary_case}
    For $d\geq2$, every bond dimension $D$ and every number $b$ of labels, there
    are $C$ and $L_0$ such that the following holds. Let $A$ be a tensor of bond
    dimension $D$, $S$ a set of bond pairs, $\omega_1,\dots,\omega_b$
    orthonormal pairs whose products vanish on every configuration of bond pairs
    leaving $S$, and $\alpha$ a unit vector. Cut the ring of $N$ sites into
    $M\geq1$ blocks of lengths $L_0\leq\ell_k\leq L$ at which every blocked
    tensor of $A$ is injective on $S$. Then
    $\sum_j\alpha_j\bigl(\bigotimes_kV^{(\ell_k)}\bigr)\bigotimes_k\ket{\omega_j}$
    is prepared with measurements and a circuit in depth at most $CL$. A
    nonzero vector prepared by a local circuit of depth $T$ is prepared with
    measurements and a circuit in depth $T$.
\end{lemma}
\begin{proof}
    \leanok
    As in the proof of Lemma~\ref{thm:ldp_orthogonal_blocks_measurement},
    encoding the labels and the bond indices in $r_1=b+D+2$ sites, preparing the
    GHZ-type state with amplitudes $\alpha$ on the registers by
    Lemma~\ref{thm:ldp_window_ghz_measurement}, and applying on every pair window
    one unitary sending the encoded label $j$ to $\ket{\omega_j}$ sets up the
    inputs. On every block,
    the unitary of Lemma~\ref{thm:ldp_block_sum_unitary} for an ordering
    of the bond pairs in which those of $S$ come first implements, by
    Lemma~\ref{thm:ldp_sequential_factorization_on}, the
    isometric factor $V^{(\ell_k)}$ on the pairs of $S$, which suffices because
    the products of the pairs are supported there. For the second statement, the
    measurement stage measures no site and the circuit after it is empty, by
    Lemma~\ref{lem:qc_measurement_unitary_case}.
\end{proof}

\begin{lemma}[A chain of one block]
    \label{lem:ldp_overlapping_one_block_chain}
    \lean{MPSPreparation.exists_blockIsometryState_eq_smul}
    \leanok
    \uses{lem:ldp_overlapping_block_states, thm:ldpolar_tensor_injective_on,
        thm:ldp_chain_block_iso_inner}
    Let the $N$-site blocked map of $\bigoplus_jA_j$ be injective on the bond
    pairs of the blocks, and let $\ket\phi=\sum_j\beta_j\ket{\phi_N(A_j)}\neq0$.
    Then $\ket\phi/\lVert\phi\rVert=V^{(N)}\ket\omega$ for a unit pair $\omega$
    supported on the bond pairs of the blocks.
\end{lemma}
\begin{proof}
    \leanok
    By Lemma~\ref{lem:ldp_overlapping_block_states}(1) with one block,
    $\ket\phi=V^{(N)}\ket Y$ with
    $Y=\sum_j\beta_j\Tr\bigl(P^{(N)}K_jK_j^\dagger\bigr)$, and the rows of
    $P^{(N)}$ at bond pairs outside the blocks vanish by
    Lemma~\ref{thm:ldpolar_tensor_injective_on}, so $Y$ is supported on the
    bond pairs of the blocks. There $V^{(N)}$ is an isometry, so
    $\lVert\phi\rVert=\lVert Y\rVert$ and $\omega=Y/\lVert Y\rVert$.
\end{proof}

\begin{lemma}[A common rate for the blocks and their mixed transfer maps]
    \label{lem:ldp_common_mixed_gap}
    \lean{MPSPreparation.exists_forall_eigenvalue_norm_le_of_mixed}
    \leanok
    \uses{lem:ldp_common_gap}
    For finitely many normal left-canonical tensors $A_1,\dots,A_b$, $b\geq1$, of
    positive bond dimensions, whose mixed transfer maps $\E_{jj'}$, $j\neq j'$,
    have all eigenvalues of modulus below one, there is $0<t<1$ bounding the
    moduli of the eigenvalues other than one of every $\E_{A_j}$ and of all
    eigenvalues of every $\E_{jj'}$.
\end{lemma}
\begin{proof}
    \leanok
    Lemma~\ref{lem:ldp_common_gap} bounds the transfer maps. Each $\E_{jj'}$ has
    finitely many eigenvalues, all of modulus below one, so their largest
    modulus $t_{jj'}$ is below one; the maximum of these bounds gives $t$.
\end{proof}

\begin{lemma}[Every translation-invariant MPS with measurements, in depth
        $O(\log(N/\varepsilon))$]
    \label{thm:ldp_overlapping_measurement_depth}
    \lean{MPSPreparation.exists_isPreparedWithMeasurementsAndCircuitInDepth_le_log_of_mpvState_ne_zero,
        MPSPreparation.exists_mpvState_repeatedBlockSum_ne_zero,
        MPSPreparation.exists_isPreparedWithMeasurementsAndCircuitInDepth_le_log}
    \leanok
    \uses{thm:ldp_injective_on_measurement, thm:ldp_overlapping_unequal_blocks_error,
        lem:ldp_overlapping_block_injectivity, lem:ldp_overlapping_one_block_chain,
        lem:ldp_common_mixed_gap, def:ldp_repeated_block_sum,
        thm:ldp_chain_depth, thm:ldp_short_chain_preparation,
        lem:ldp_log_threshold_error, lem:ldp_shared_overlap_estimates}
    Let $A$ be the direct sum with multiplicities of
    Definition~\ref{def:ldp_repeated_block_sum}, with nonzero complex weights,
    every block $A_j$ normal in the gauge (\ref{eq:ldp_normal_gauge}) with a
    positive definite fixed point $\rho_j$ of trace one, and let the mixed
    transfer maps $\E_{jj'}$, $j\neq j'$, have all eigenvalues of modulus below
    one. There is $c$, depending only on $A$, such that for every $N\geq2$ and
    $0<\varepsilon\leq1$ at which $\ket{\phi_N(A)}\neq0$, a unit vector
    $\ket\psi$ with $\varepsilon(\psi,\phi_N)\leq\varepsilon$ is prepared with
    measurements and a circuit in depth at most $c\log(N/\varepsilon)$. There is
    $N_0$ such that $\ket{\phi_N(A)}\neq0$ for every $N\geq N_0$ at which the
    weights $\beta_j=\sum_k\mu_{j,k}^N$ are not all zero, so the same holds for
    these $N$. Under the readings recorded in
    \cite{gap:mswc24_measurement_preparation_scope,gap:mswc24_block_form_mixed_overlap,gap:mswc24_depth_upper_bound_nonzero_state},
    this establishes the depth $O(\log(N/\varepsilon))$ of the paragraph
    ``Long-range MPS using measurements'' of \cite{Malz2024Preparation}, with
    measurements used only for the GHZ-type state, for blocks whose $q$-site
    states may overlap, every multiplicity, every nonzero complex weight and
    every chain length.
\end{lemma}
\begin{proof}
    \leanok
    Let $t$ be the rate of Lemma~\ref{lem:ldp_common_mixed_gap},
    $r=-\tfrac14\log t$, $K$ the constant of
    Lemma~\ref{thm:ldp_overlapping_unequal_blocks_error} at $\gamma=1/4$, and
    $L_{\min}$ larger than the lengths of
    Lemma~\ref{lem:ldp_overlapping_block_injectivity}(2) and
    Lemma~\ref{thm:ldp_injective_on_measurement}. Let
    $Q=a\log(N/\varepsilon)+b'$ with $a=1/r$ and
    $b'=\max(\log K,0)/r+L_{\min}$, and $q=\lceil Q\rceil$. When $q\leq N$, cutting the
    chain into $\lfloor N/q\rfloor-1$ blocks of length $q$ and one of length
    $q+(N\bmod q)<2q$ ensures that every blocked map is injective on the bond pairs of the
    blocks, the state of Lemma~\ref{thm:ldp_overlapping_unequal_blocks_error}
    has error at most $KNe^{-rq}\leq\varepsilon$ by
    Lemma~\ref{lem:ldp_log_threshold_error}, and it is prepared in depth
    $O(q)=O(\log(N/\varepsilon))$ by
    Lemma~\ref{thm:ldp_injective_on_measurement}. If $L_{\min}\leq N<q$, the
    normalized target itself is $V^{(N)}\ket\omega$ by
    Lemma~\ref{lem:ldp_overlapping_one_block_chain} and is prepared by a circuit
    of depth $O(N)=O(\log(N/\varepsilon))$ by Theorem~\ref{thm:ldp_chain_depth}.
    If $N<L_{\min}$, it is prepared in bounded depth by
    Lemma~\ref{thm:ldp_short_chain_preparation}. For the second statement,
    $\lvert\lVert\phi_N(A)\rVert^2/w-1\rvert\leq Ke^{-N/(4\xi_t)}$, with
    $\xi_t$ the correlation length at $t$ and
    $w=\sum_j\lvert\beta_j\rvert^2$ by
    Lemma~\ref{lem:ldp_shared_overlap_estimates}(6), which is below one for
    large $N$.
\end{proof}

\begin{lemma}[Translation-invariant MPS from product states with
        measurements, one direction]
    \label{thm:ldp_qccc_classification_forward}
    \lean{MPSPreparation.isAsymptoticallyPreparedWithMeasurementsInDepth_normalizedMPVState}
    \leanok
    \uses{thm:ldp_overlapping_measurement_depth,
        lem:qc_asymptotic_preparation_of_overlap,
        def:qc_asymptotic_measurement_preparation}
    Let $A$ be as in Lemma~\ref{thm:ldp_overlapping_measurement_depth}, and
    let $\ket{\phi_N(A)}\neq0$ for all large $N$. There is $C$ such that the
    sequence of normalized states $\ket{\phi_N}$ is asymptotically prepared
    with measurements in depth $C\log N$: unit vectors $\ket{\psi_N}$,
    prepared with measurements and a circuit in depth at most $C\log N$,
    satisfy
    $\lVert\ket{\psi_N}\bra{\psi_N}-\ket{\phi_N}\bra{\phi_N}\rVert_1\to0$.

    For the equivalence of sequences of states in the paragraph ``Phases of
    matter'' of \cite{Piroli2021LOCC}, this proves the direction from the
    trivial sequence to the sequence of MPS of the theorem stating that in one
    dimension every translation-invariant MPS of fixed bond dimension is in the
    same phase as the trivial state, in polylogarithmic depth $C\log N$.
    As recorded in \cite{gap:psc21_mps_classification_scope}, the converse
    direction and the reduction of an arbitrary translation-invariant tensor to
    the canonical form assumed here are not formalized.
\end{lemma}
\begin{proof}
    \leanok
    Let $c$ be the constant of
    Lemma~\ref{thm:ldp_overlapping_measurement_depth}. For $N\geq2$ with
    $\ket{\phi_N(A)}\neq0$ and $\varepsilon=1/N$, a unit vector of error at
    most $1/N$ is prepared with measurements and a circuit in depth at most
    $c\log(N^2)=2c\log N$. Lemma~\ref{lem:qc_asymptotic_preparation_of_overlap}
    with $C=2c$ and $\varepsilon_N=1/N$ concludes.
\end{proof}


\section{Bounded physical leaves and arbitrary chain lengths}
\label{sec:ldp_unequal_mera}

The isometries in the network need not be equal within a layer.
Allowing unequal positive block lengths removes the divisibility restriction,
while a uniform bound on the physical leaves keeps every local dimension
independent of the chain length and the requested accuracy.

\begin{definition}[Isometry trees with bounded leaves]
    \label{def:ldp_bounded_isometry_tree}
    \lean{MPSPreparation.IsometryTree, MPSPreparation.IsometryTree.matrix,
        MPSPreparation.IsometryTree.pairIndexEquiv,
        MPSPreparation.IsometryTree.joinMatrix, MPSPreparation.IsometryTree.joinMatrix_apply,
        MPSPreparation.IsometryTree.matrix_fork_apply,
        MPSPreparation.IsometryTree.cast, MPSPreparation.IsometryTree.matrix_cast}
    \leanok
    Fix a physical dimension $d$, a virtual dimension $\chi$ and a leaf
    bound $R$. A depth-zero tree on $n$ sites, with $1\le n\le R$, is
    an isometry $V\colon\C^\chi\to(\C^d)^{\otimes n}$.
    A tree of coarse depth $h+1$ consists of two trees $T_1,T_2$ of
    coarse depth $h$, on $n_1,n_2$ sites, and an isometry
    $W\colon\C^\chi\to\C^\chi\otimes\C^\chi$.
    Its contraction is
    \begin{align}
        T & =(T_1\otimes T_2)W
        \colon\C^\chi\longrightarrow(\C^d)^{\otimes(n_1+n_2)}.
        \label{eq:ldp_unequal_tree_contraction}
    \end{align}
    The physical sites keep their order under concatenation.
    A tree of coarse depth $h$ has $2^h$ leaves and $h+1$ isometry
    layers, counting the finest layer.
\end{definition}

\begin{theorem}[Isometric contraction of a tree]
    \label{thm:ldp_bounded_tree_isometry}
    \lean{MPSPreparation.IsometryTree.isIsometry_matrix,
        MPSPreparation.IsometryTree.isIsometry_joinMatrix}
    \leanok
    \uses{def:ldp_bounded_isometry_tree}
    Every tree contraction $T$ satisfies $T^\dagger T=\Id$.
\end{theorem}
\begin{proof}
    \leanok
    By induction on the tree, at a binary vertex one has
    $T^\dagger T=W^\dagger(T_1^\dagger T_1\otimes T_2^\dagger T_2)W
        =W^\dagger W=\Id$.
\end{proof}

\begin{theorem}[Polar trees for unequal block lengths]
    \label{thm:ldp_unequal_polar_tree}
    \lean{MPSPreparation.exists_isometryTree_matrix_eq_cfgPolarIso}
    \leanok
    \uses{def:ldp_bounded_isometry_tree, def:ldpolar_matrix, def:injective}
    Let $s\ge1$, and assume that $A$ blocked over every $m\ge s$ sites
    is injective. If $2^h s\le n\le2^h c$, its polar isometry $V_n$
    is the contraction of a tree of coarse depth $h$ whose physical
    leaves contain at most $c$ sites.
\end{theorem}
\begin{proof}
    \leanok
    \uses{lem:ldp_polar_merge}
    Splitting $n$ into $n_1=\lfloor n/2\rfloor$ and $n_2=n-n_1$ gives
    $2^h s\le n_i\le2^h c$ at a successor depth, with no rounding error
    added to the upper bound because its endpoint is even. The polar-merge
    identity gives $V_{n_1+n_2}=(V_{n_1}\otimes V_{n_2})W_{n_1,n_2}$,
    where injectivity makes $W_{n_1,n_2}$ an isometry. Induction then
    gives the two descendants, and at depth zero the polar isometry is
    itself the bounded leaf.
\end{proof}

\begin{theorem}[Polar trees on a common supported space]
    \label{thm:ldp_supported_polar_tree}
    \lean{MPSPreparation.supportedCfgPolarIso,
        MPSPreparation.isIsometry_supportedCfgPolarIso,
        MPSPreparation.supportedMergeIso,
        MPSPreparation.supportedCfgPolarIso_append,
        MPSPreparation.supportedCfgPolarIso_split,
        MPSPreparation.supportedCfgPolarIso_eq_joinMatrix_mul,
        MPSPreparation.isIsometry_supportedMergeIso,
        MPSPreparation.exists_isometryTree_matrix_eq_supportedCfgPolarIso}
    \leanok
    \uses{def:ldp_bounded_isometry_tree, def:ldpolar_injective_on}
    Let $S$ be a set of bond pairs and let $J\colon\C^\chi\to\C^{D^2}$
    be an injective coordinate inclusion with range $\C^S$.
    Suppose that for every $m\ge s\ge1$, the matrices of $A^{[m]}$
    vanish outside $S$ and span every matrix unit indexed by $S$.
    For $A^{[m]}=V_mP_m$ and $\widehat V_m=V_mJ$, one has
    $\widehat V_m^\dagger\widehat V_m=I_\chi$.
    For $l,r\ge s$, the polar merge $W_{l,r}$ restricts to
    $\widehat W_{l,r}=(J^\dagger\otimes J^\dagger)W_{l,r}J$, and
    \begin{align}
        \widehat V_{l+r}
         & =(\widehat V_l\otimes\widehat V_r)\widehat W_{l,r},
         & \widehat W_{l,r}^\dagger\widehat W_{l,r}            & =I_\chi.
        \label{eq:ldp_supported_polar_merge}
    \end{align}
    Consequently, in the supported-space form of the tree in
    \cite[eq.~(16), Long-range MPS using measurements]{Malz2024Preparation},
    if $2^h s\le n\le2^h c$, then $\widehat V_n$ is a tree of coarse depth $h$
    with at most $c$ sites per physical leaf, and the same maps act on every
    superposition of the pairs in $S$.

    % Formula: Vhat_(l+r) = (Vhat_l tensor Vhat_r) What_(l,r).
    % Source: MSWC24 eq. (16), restricted to the common support S.
    % Ink-to-index map: virtual wires have dimension chi, not D^2.
    % Boundary signature on both sides: one C^chi input and ordered H_l,H_r outputs.
    % Internal contractions: the two C^chi outputs of What and the inputs of Vhat_l,Vhat_r.
    % No projection, normalization scalar, discarded environment or postselection is suppressed.
    \begin{center}
        \tenkzkernel
        \begin{tenkzequation}
            \begin{tenkz}[rows={wire}, cols=1, bonds=none]
                \tn[skin=box, ports={135:physical:$\mathcal H_l$,
                            45:physical:$\mathcal H_r$,270:virtual:$\C^\chi$}]{\widehat V_{l+r}}
            \end{tenkz}
            =
            \begin{tenkz}[rows={wire,wire}, cols=3, bonds=none]
                \tn[at=(1,1), skin=box, name=vl,
                    ports={90:physical:$\mathcal H_l$,270:virtual}]{\widehat V_l}
                \tn[at=(1,3), skin=box, name=vr,
                    ports={90:physical:$\mathcal H_r$,270:virtual}]{\widehat V_r}
                \tn[at=(2,2), skin=box, name=w,
                    ports={135:virtual,45:virtual,270:virtual:$\C^\chi$}]{\widehat W_{l,r}}
                \tnwire{w.135}{vl.270}
                \tnwire{w.45}{vr.270}
            \end{tenkz}
        \end{tenkzequation}
    \end{center}
    Here $\mathcal H_m=(\C^d)^{\otimes m}$.
    The boxes $\widehat V_m$ denote entire subtrees, not free gates on $m$ sites.
\end{theorem}
\begin{proof}
    \leanok
    \uses{thm:ldpolar_tensor_injective_on, lem:ldp_polar_merge}
    The polar support is the coordinate projection $JJ^\dagger$,
    so $V_m$ vanishes outside the selected columns and
    $J^\dagger V_m^\dagger V_mJ=I_\chi$.
    Restricting the polar-merge sum to those columns in each child gives
    (\ref{eq:ldp_supported_polar_merge}). With
    $L=\widehat V_l\otimes\widehat V_r$, this also proves that the
    image of $\widehat V_{l+r}$ is contained in the image of $L$.
    Since $L^\dagger L=I$, the restricted merge satisfies
    $\widehat W^\dagger\widehat W
        =\widehat V_{l+r}^\dagger\widehat V_{l+r}=I_\chi$, and
    repeated floor/ceiling splits give the bounded-leaf tree.
\end{proof}

\begin{definition}[Finite-range networks on unequal blocks]
    \label{def:ldp_unequal_mera}
    \lean{MPSPreparation.UnequalMERA, MPSPreparation.UnequalMERA.topPair,
        MPSPreparation.UnequalMERA.treeMatrix,
        MPSPreparation.UnequalMERA.treeMatrix_apply,
        MPSPreparation.UnequalMERA.state, MPSPreparation.UnequalMERA.state_apply,
        MPSPreparation.UnequalMERA.ofTrees}
    \leanok
    \uses{def:ldp_bounded_isometry_tree}
    Let $D,M\ge1$, let $\ell_0+\cdots+\ell_{M-1}=N$, and let $T_k$
    be trees of common coarse depth $h$, virtual dimension $D^2$ and
    physical leaf bound $R$, on their respective blocks of $\ell_k$
    sites. A unitary $u$ on $\C^D\otimes\C^D$ prepares a pair
    $\ket\omega=u\ket{0,0}$. The network state is
    \begin{align}
        \ket\psi
         & =\left(\bigotimes_{k=0}^{M-1}T_k\right)
        \bigotimes_{k=0}^{M-1}\ket\omega_{R_kL_{k+1}},
        \label{eq:ldp_unequal_mera_state}
    \end{align}
    with cyclic indexing of the pair legs. All lower disentanglers
    are identities. The network has $h+2$ layers, including the finest
    isometries and the top unitary layer. Internal tree legs have
    dimension $D^2$; every physical leaf has dimension at most $d^R$.
\end{definition}

\begin{theorem}[Normalization and arbitrary top pairs]
    \label{thm:ldp_unequal_mera_norm}
    \lean{MPSPreparation.UnequalMERA.isIsometry_treeMatrix,
        MPSPreparation.UnequalMERA.sum_star_topPair_mul,
        MPSPreparation.UnequalMERA.norm_state,
        MPSPreparation.UnequalMERA.exists_disentangler_pair,
        MPSPreparation.UnequalMERA.exists_of_trees}
    \leanok
    \uses{def:ldp_unequal_mera}
    The state in (\ref{eq:ldp_unequal_mera_state}) has norm one.
    Conversely, for any such forest and any unit pair $\ket\omega$,
    a top unitary realizes that pair and hence the forest contraction.
\end{theorem}
\begin{proof}
    \leanok
    \uses{thm:ldp_bounded_tree_isometry, thm:ldp_chain_block_iso_inner,
        thm:ldp_isometric_extension_overlap, thm:ldp_unitary_implementing_isometry}
    Completing the unit column $\ket\omega$ to a unitary taking
    $\ket{0,0}$ to it realizes the top pair, while tensor products of the
    normalized pairs have norm one and the isometric forest preserves
    their norm.
\end{proof}

\begin{theorem}[Polar forests and the block-isometry state]
    \label{thm:ldp_unequal_mera_polar_state}
    \lean{MPSPreparation.UnequalMERA.state_eq_blockIsometryState,
        MPSPreparation.UnequalMERA.exists_of_trees_eq_blockIsometryState}
    \leanok
    \uses{def:ldp_unequal_mera, def:ldp_block_isometry_state}
    If $T_k=V_{\ell_k}$ and the top unitary prepares $\ket\omega$,
    the network state equals
    $(\bigotimes_k V_{\ell_k})\bigotimes_k\ket\omega_{R_kL_{k+1}}$.
    In particular, every normalized top pair and every forest of
    these polar trees give a network representing this block-isometry state.
\end{theorem}
\begin{proof}
    \leanok
    \uses{def:ldp_unequal_mera, def:ldp_block_isometry_state, thm:ldp_unequal_mera_norm}
    Substituting the tree contractions into
    (\ref{eq:ldp_unequal_mera_state}) regroups the physical
    configurations according to the prescribed block partition.
\end{proof}

\begin{theorem}[Bounded single leaves and product states]
    \label{thm:ldp_unequal_mera_caps}
    \lean{MPSPreparation.UnequalMERA.exists_single_leaf,
        MPSPreparation.UnequalMERA.nonempty_bond_one}
    \leanok
    \uses{def:ldp_unequal_mera}
    Every unit state on $1\le N\le R$ sites is represented exactly
    by a single leaf with virtual dimension one and coarse depth zero.
    If $d,R,N\ge1$, a product-state network of coarse depth zero
    also exists, with one physical site at each leaf and virtual
    dimension one.
\end{theorem}
\begin{proof}
    \leanok
    \uses{thm:ldp_unequal_mera_norm}
    The unit column $1\mapsto\ket\psi$ and the scalar top unitary give
    the first assertion, while the column $1\mapsto\ket0$ applied
    separately at every site gives the second.
\end{proof}

\begin{theorem}[Uniform finite-range networks at every positive length]
    \label{thm:ldp_every_length_mera}
    \lean{MPSPreparation.exists_unequalMERA_every_length}
    \leanok
    \uses{def:ldp_unequal_mera, def:normal}
    Let $A$ be a normal tensor of physical dimension $d$ and bond
    dimension $D$. There are constants $R\ge1$ and $C\ge2$, chosen
    before $\varepsilon$ and $N$, such that for every $\varepsilon>0$
    and every positive length $N$ with nonzero periodic state, a
    network with some $1\le D'\le D$, physical leaf bound $R$ and
    coarse depth $h$ satisfies
    \begin{align}
        h+2
         & \le C\left[1+\log\!\left(1+\max\{0,\log(N/\varepsilon)\}\right)\right],
        \label{eq:ldp_every_length_mera_layers}                                    \\
        1-\left|\braket{\psi}{\phi_N(A)}\right|
         & \le\varepsilon.
        \label{eq:ldp_every_length_mera_error}
    \end{align}
    Its internal tree dimensions are at most $D^2$, and its physical
    leaf dimensions are at most $d^R$. Thus the layer count is
    $O(\log\log(N/\varepsilon))$ in the large-size, small-error regime,
    with no divisibility restriction on $N$.
\end{theorem}
\begin{proof}
    \leanok
    \uses{thm:ldp_unequal_polar_tree, thm:ldp_unequal_mera_polar_state,
        thm:ldp_unequal_mera_caps, thm:ldp_block_approximation_error,
        thm:ldp_short_chain_preparation, lem:ldp_normalized_state_rescaling,
        thm:ldp_depth_upper_bound_divisible, thm:prep_pair_norm}
    For $\varepsilon>1$, a product state suffices, so assume
    $0<\varepsilon\le1$. Passing to a normal left-canonical representative $B$ changes
    normalized states only by a phase; let $s\ge1$ be its
    eventual injectivity threshold and set $R=4s$. The unequal-block
    second-order estimate at $\gamma=1/2$ gives constants $K,r>0$
    with error at most $KM e^{-rq}$ for blocks of length at least $q$.
    Let $q=\lceil a\log(N/\varepsilon)+b\rceil$, with $a=1/r$
    and fixed $b$ large enough that $q\ge s$ and $KN e^{-rq}\le\varepsilon$.

    If $q\le N$, partitioning into $M=\lfloor N/q\rfloor$ blocks, the first
    $M-1$ of length $q$ and the last of length $q+N\bmod q$,
    and choosing $h$ with $2^h s\le q<2^h(2s)$ places every block length
    between $2^h s$ and $2^h(4s)$, so
    Theorem~\ref{thm:ldp_unequal_polar_tree} gives its bounded-leaf
    tree, and the fixed-point pair at the top together with the
    unequal-block error estimate yields the stated fidelity bound.

    If $s\le N<q$, one polar tree over $N$, with $h$ chosen
    from $N/s$, together with the normalized trace pair at its root, prepares the
    normalized periodic state exactly; its terminal widths remain at most $4s$
    rather than forming a single leaf of accuracy-dependent width.
    If $N<s$, the bounded single leaf of
    Theorem~\ref{thm:ldp_unequal_mera_caps} gives the state exactly,
    with $D'=1$.

    In the nontrivial branches $2^h\le q\le a\log(N/\varepsilon)+b+1$.
    Taking logarithms and adding the finest and top layers gives
    (\ref{eq:ldp_every_length_mera_layers}). The regularized expression
    also covers $N=1$ and $\varepsilon\ge1$.
\end{proof}

\begin{theorem}[Normal states at every sufficiently large length]
    \label{thm:ldp_eventual_every_length_mera}
    \lean{MPSPreparation.exists_unequalMERA_of_normal}
    \leanok
    \uses{def:ldp_unequal_mera, def:normal}
    For a normal tensor with $D\ge1$, the same fixed dimension and
    layer bounds hold for all $N\ge N_0$, for some fixed $N_0$,
    without a separate nonvanishing hypothesis.
\end{theorem}
\begin{proof}
    \leanok
    \uses{thm:ldp_every_length_mera, lem:ldp_periodic_state_nonzero}
    Normality implies eventual nonvanishing of the periodic state, so
    Theorem~\ref{thm:ldp_every_length_mera} applies beyond that threshold.
\end{proof}


\section{Non-normal canonical states in double-logarithmic depth}
\label{sec:ldp_nonnormal_trees}

The paragraph ``Tree-RG circuit with measurements'' of
\cite{Malz2024Preparation} gives double-logarithmic depth for short-range
correlated states. Its paragraph ``Long-range MPS using measurements''
prepares coherent sector labels before the blocked isometries. Combining
these constructions on their common virtual support extends the
measurement-assisted depth bound to the canonical non-normal class below.
This is an enhancement of those constructions, not a claim that the source
already proves the all-length statement in this form.

\begin{theorem}[Inequivalence gives the mixed eigenvalue gap]
    \label{thm:ldp_inequivalent_mixed_gap}
    \lean{MPSTensor.mixedMap_eigenvalue_norm_lt_one_of_spectralRadius_lt_one,
        MPSTensor.mixedMap_eigenvalue_norm_lt_one_of_normal_inequivalent}
    \leanok
    Let $A_j,A_k$ be normal left-canonical blocks of positive bond dimensions.
    If their dimensions agree, assume they are not gauge-phase equivalent.
    Every eigenvalue of their mixed transfer map has modulus less than one.
\end{theorem}
\begin{proof}
    \leanok
    Normality implies irreducibility. The mixed-transfer spectral-radius
    theorem gives a strict gap, both for inequivalent square blocks and for
    different bond dimensions. An eigenvalue belongs to the spectrum of
    the mixed map and hence its modulus is bounded by that spectral radius.
\end{proof}

\begin{theorem}[Coherent supported block trees from sparse labels]
    \label{thm:ldp_coherent_supported_blocks}
    \lean{QuantumCircuit.IsPreparedWithMeasurementRoundsInDepth.apply_implementation,
        MPSPreparation.exists_pairFamilyUnitary,
        MPSPreparation.pairLayerOp_mulVec_registerCfg_eq,
        MPSPreparation.blockLayerOp_pairLayerOp_mulVec_registerCfg,
        MPSPreparation.exists_isPreparedWithMeasurementRoundsInDepth_sum_blockIsometryState,
        MPSPreparation.exists_supported_block_preparation_constants}
    \leanok
    \uses{thm:ldp_sparse_window_ghz, thm:ldp_supported_tree_rounds,
        lem:ldp_balanced_widths}
    Fix $d\ge2$, a tensor $B$, a set $S$ of virtual pairs, an eventual
    support-injectivity threshold $L$, and a bound on the number of sector
    labels. There are $s\ge\max(2,L)$ and $C$, fixed before the ring and
    amplitudes, with the following property. Suppose
    \begin{align*}
        2^{h+2}s\le\ell_k\le2^{h+1}(8s).
    \end{align*}
    Let the pairs $\omega_j$ be orthonormal and $\sum_j|\alpha_j|^2=1$.
    Assume that, for the number $M$ of blocks, each product of the
    cyclic pairs $\omega_j$ is supported on $S^M$. Then
    \begin{align*}
        \sum_j\alpha_j\Bigl(\bigotimes_kV_{\ell_k}\Bigr)
        \bigotimes_k\ket{\omega_j}_{R_kL_{k+1}}
    \end{align*}
    is prepared in quantum depth at most $C(h+1)$.
\end{theorem}
\begin{proof}
    \leanok
    The sparse GHZ construction has depth independent of the block lengths.
    A single bounded-window unitary maps every encoded label $j$ to
    $\omega_j$, preserving all coherent amplitudes. The bond encodings have
    a final zero site. Thus their full physical configurations agree exactly
    with the tree-root windows even when an odd block leaves one terminal
    site outside the even leaf cutoff. The pair layer leaves all tree
    interior scratch sites zero, so the supported compiler applies.
    Compose the rounds and add their depths. For a fixed outcome
    history, its scalar is common to all inputs in the zero-scratch
    subspace. The endpoint specifies every physical site.
\end{proof}

\begin{theorem}[All-length compilation of exponential block approximations]
    \label{thm:ldp_supported_loglog_compilation}
    \lean{MPSPreparation.exists_log_log_preparation_of_supported_block_approximation,
        MPSPreparation.natCast_mul_succ_le_mul_log,
        MPSPreparation.exists_two_pow_le}
    \leanok
    \uses{thm:ldp_coherent_supported_blocks, lem:ldp_block_length_threshold_error}
    Under the fixed tensor, support and eventual support-injectivity
    assumptions of Theorem~\ref{thm:ldp_coherent_supported_blocks},
    suppose a normalized target family has coherent supported block
    approximations of error at most $KM e^{-rq}$ for block lengths
    $\ell_k\ge q\ge L$, with $K,r>0$ independent of the amplitudes and
    lengths. Suppose also that every admissible target of length $N\ge L$
    has an exact normalized one-block supported pair. Then a fixed $c$
    gives, for every admissible $N\ge2$ and $0<\varepsilon\le1$, a unit
    approximant of error at most $\varepsilon$ prepared in depth
    \begin{align*}
        T\le c\log\bigl(1+\log(N/\varepsilon)\bigr).
    \end{align*}
    The one-block support assumption is required only for $M=1$.
\end{theorem}
\begin{proof}
    \leanok
    With $q=\lceil a\log(N/\varepsilon)+b\rceil$, where $a=1/r$ and
    $b\ge\max(\log K,0)/r+4s+1$, when $q\le N$, partitioning into $M=\lfloor N/q\rfloor$
    blocks, the last of length $q+N\bmod q$ and all others of length $q$, and
    selecting $h$ with $2^{h+2}s\le q<2^{h+1}(4s)$ ensures that all block lengths satisfy
    the previous theorem, and $KM e^{-rq}\le KN e^{-rq}\le\varepsilon$.

    If $4s\le N<q$, a whole-ring supported tree with $h$ chosen from
    $N$ rather than from $q$ has a root pair that gives the target exactly.
    In both branches $2^h\le q\le a\log(N/\varepsilon)+b+1$, so the
    depth $C(h+1)$ has the required bound. For $N<4s$, exact synthesis
    on finitely many possible lengths is absorbed into the constant,
    since $\log(1+\log(N/\varepsilon))\ge\log(1+\log2)>0$.

    % Formula: normalized target = Vhat_N eta_N on the one-block support.
    % Source: supported polar factorization and the exact one-block pair.
    % Boundary signature on each side: no input and one H_N physical output.
    % Internal contraction: a chi-dimensional supported root state into Vhat_N.
    % Vhat_N is the entire bounded-leaf tree, never a free N-site gate.
    \begin{center}
        \tenkzkernel
        \begin{tenkzequation}
            \begin{tenkz}[rows={wire}, cols=1, bonds=none]
                \tn[skin=box, ports={90:physical:$\mathcal H_N$}]{\widehat\phi_N}
            \end{tenkz}
            =
            \begin{tenkz}[rows={wire,wire}, cols=1, bonds=none]
                \tn[at=(1,1), skin=box, name=tree,
                    ports={90:physical:$\mathcal H_N$,270:virtual}]{\widehat V_N}
                \tn[at=(2,1), skin=box, name=root,
                    ports={90:virtual:$\C^\chi$}]{\eta_N}
                \tnwire{root.90}{tree.270}
            \end{tenkz}
        \end{tenkzequation}
    \end{center}
    Here $\widehat\phi_N=\phi_N/\|\phi_N\|$, and $\eta_N$ is the
    normalized supported one-block input, including
    the cyclic ordering of the two virtual legs. The box $\widehat V_N$
    denotes the entire bounded-leaf tree, not one free $N$-site gate.
\end{proof}

\begin{theorem}[Non-normal canonical preparation in double-logarithmic depth]
    \label{thm:ldp_nonnormal_loglog}
    \lean{MPSPreparation.exists_isPreparedWithMeasurementRoundsInDepth_le_log_log_of_mpv_eq_sum,
        MPSPreparation.exists_isPreparedWithMeasurementRoundsInDepth_repeatedBlockSum_le_log_log,
        MPSPreparation.exists_isPreparedWithMeasurementRoundsInDepth_repeatedBlockSum_le_log_log_of_inequivalent}
    \leanok
    \uses{thm:ldp_inequivalent_mixed_gap, thm:ldp_supported_loglog_compilation,
        def:ldp_repeated_block_sum, lem:ldp_common_mixed_gap,
        lem:ldp_overlapping_block_injectivity, thm:ldp_overlapping_unequal_blocks_error,
        lem:ldp_overlapping_one_block_chain}
    Let
    \begin{align*}
        A^i=\bigoplus_j\operatorname{diag}(\mu_{j,1},\ldots,\mu_{j,m_j})
        \otimes A_j^i
    \end{align*}
    be the canonical repeated-block form, with normal inequivalent blocks
    in the left-canonical gauge, and positive definite transfer fixed points
    $\sigma_j$ of trace one. The copy weights are arbitrary nonzero complex
    numbers under the existing convention that zero padding is omitted.
    There is $c$, fixed before $N$ and $\varepsilon$, such that every
    $N\ge2$ with $\phi_N(A)\ne0$ and every $0<\varepsilon\le1$ admit
    a unit vector $\psi$ satisfying
    \begin{align*}
        1-|\langle\psi,\phi_N(A)/\|\phi_N(A)\|\rangle|\le\varepsilon,
        \qquad T\le c\log\bigl(1+\log(N/\varepsilon)\bigr).
    \end{align*}
    The protocol prepares $\psi$ in measurement rounds of quantum depth $T$.
    Finite-block physical ranges may overlap. Individual sums
    $\beta_j(N)=\sum_k\mu_{j,k}^N$ may vanish.
\end{theorem}
\begin{proof}
    \leanok
    For the unweighted direct sum $B=\bigoplus_jA_j$ and its
    diagonal sector-pair support, the mixed gap gives eventual support
    injectivity and the amplitude-uniform overlap estimate. With
    $\alpha_j=\beta_j/(\sum_l|\beta_l|^2)^{1/2}$, nonvanishing of the
    target implies that the coefficient vector is nonzero, without
    requiring each coefficient to be nonzero. The embedded fixed-point
    pairs are orthonormal, their cyclic products have the required
    support, and the exact one-block pair gives the short accuracy branch,
    so Theorem~\ref{thm:ldp_supported_loglog_compilation} applies.

    The intermediate theorem accepts any identity
    $\phi_N(A)=\sum_j\beta_j(N)\phi_N(A_j)$, so the coefficients need not
    be powers for that interface. For the repeated-block form the trace
    identity gives precisely the power sums above.
\end{proof}

Only nearest-neighbour unitary layers contribute to $T$. Measurements,
classical communication and on-site corrections are free, the number of
rounds may grow with the tree height, and no outcome is postselected.
The theorem concerns the given canonical block class. Passing from
an arbitrary original tensor, including its periodicity and blocking data,
is a separate statement.

